% Research contribution by /root/shortening.  This is an original contribution,
% not an extraction or copy of an attached source.
\section{Strict shortening and finite price obstructions}
\label{sec:shortening-prices}

Throughout this section, a spectral point and a price have precisely the
meanings of Definition A.1 and Definition A.4 of the attached, unrefereed
Anthropic manuscript (pp.~19--20).  We use its Lemma A.2 (p.~19) for the
evaluation of products of wires.  All arguments below work over an arbitrary
field unless an additional hypothesis is explicitly stated.

The unconditional low-width conclusion of this section is
Corollary~\ref{cor:strict-shortening-width-at-least-seven}:
every strict shortening, with arbitrary finite term layouts,
has cable width at least seven.  The planar two-summand
case requires width at least eight.  The unrestricted
existence question remains unresolved.

\subsection{Normalization of the enlarged spectrum}

\begin{proposition}[SH-10: power normalization]
\label{prop:spectral-normalization}
Let $I$ be a finite party set with at least two parties, and let
$\varphi$ be a spectral point in the enlarged sense of Definition A.1.
Set $\alpha=g_\varphi(I)$.  Then $0\le\alpha\le1$.
If $\alpha=0$, $\varphi$ equals one on every nonzero tensor and zero on
the zero tensor.  If $\alpha>0$, the function
$\psi(T)=\varphi(T)^{1/\alpha}$ is multiplicative, additive under full
direct sum, monotone under restriction, and satisfies
$\psi(\langle q\rangle_I)=q$.  In particular, $\psi$ is a normalized
Strassen spectral point.  Conversely, every power $\psi^\alpha$ with
$0<\alpha\le1$ of a normalized spectral point has the properties of
Definition A.1.
\end{proposition}
\begin{proof}
The unit tensor $\langle q\rangle_I$ is covered by $q$ copies of the
scalar unit, so Lemma A.2 gives $0\le\alpha\le1$.  Every nonzero tensor
$T$ is a restriction of $\langle R(T)\rangle_I$, where $R(T)$ is its
finite tensor rank, and restricts to the scalar unit.  Thus
$1\le\varphi(T)\le R(T)^\alpha$.  This proves the assertion when
$\alpha=0$.

Suppose $\alpha>0$.  For tensors $S,T$, the full direct sum
$(S\oplus T)^{\otimes N}$ splits by multiplicity into $N+1$ blocks.
The block with $k$ copies of $S$ is isomorphic, by local basis
permutations, to
\[
 \left\langle\binom Nk\right\rangle_I
   \otimes S^{\otimes k}\otimes T^{\otimes(N-k)}.
\]
Every block is a restriction of the whole tensor, and the whole tensor
is covered by the blocks.  Hence, writing
$a=\varphi(S)^{1/\alpha}$ and $b=\varphi(T)^{1/\alpha}$,
\[
 \max_{0\le k\le N}
   \left[\binom Nk a^k b^{N-k}\right]^\alpha
 \le \varphi(S\oplus T)^N
 \le \sum_{k=0}^N
   \left[\binom Nk a^k b^{N-k}\right]^\alpha.
\]
By the binomial theorem the unpowered bracket terms have sum
$(a+b)^N$.  Their maximum lies between $(a+b)^N/(N+1)$ and
$(a+b)^N$.  The displayed lower and upper bounds consequently differ
from $(a+b)^{\alpha N}$ only by factors between
$(N+1)^{-\alpha}$ and $N+1$.  Taking $N$th roots and letting
$N\to\infty$ gives
\[
 \varphi(S\oplus T)=(a+b)^\alpha.
\]
This proves additivity of $\psi$; its other properties follow directly
from those of $\varphi$.  A tensor covered by $T_1,\ldots,T_r$ is a
restriction of their full direct sum, by using local fold maps, so
ordinary additivity and monotonicity give cover subadditivity of
$\psi$.  Finally
$(\sum_j a_j)^\alpha\le\sum_j a_j^\alpha$ for nonnegative numbers
and $0<\alpha\le1$, proving the converse.
\end{proof}

% Further statements in this section are appended below; stable SH labels
% are used rather than relying on their eventual theorem numbers.



\begin{corollary}[SH-11: finite prices need no limiting closure]
\label{cor:price-normal-form}
Let $|I|\ge2$, and let $\Delta_I$ denote the normalized Strassen
spectrum on all tensors on $I$.  The finite price cone is exactly
\[
 \{s\,d_\psi:s\ge0,\ \psi\in\Delta_I\}.
\]
Thus every finite price has an exact normalized spectral realization,
although finding it may be nonconstructive.  For every spanning tree
$\mathcal T$ on $I$, the scale in such a representation satisfies
$s\le\sum_{e\in\mathcal T}p(e)$.
\end{corollary}
\begin{proof}
The normalized spectrum is a compact set in the topology of pointwise
convergence.  Indeed every coordinate $\psi(T)$ lies between zero and
the finite tensor rank $R(T)$, and normalization, multiplicativity,
additivity, and monotonicity are closed requirements in the resulting
compact product.  A bipartition flattening rank supplies a point, so
the set is nonempty.  Its wire-metric image is therefore compact.

By Proposition~\ref{prop:spectral-normalization}, any nonzero scaled
enlarged spectral metric has the form $s d_\psi$ with
$\psi\in\Delta_I$.  Lemma A.2 gives
$1=g_\psi(I)\le\sum_{e\in\mathcal T}d_\psi(e)$, whence
$s\le\sum_{e\in\mathcal T}s d_\psi(e)$.  In any convergent sequence
of such metrics, these scales are bounded.  The compactness just
proved allows a convergent subnet of the scales and normalized wire
metrics.  Its limit has the form $s d_\psi$.  Zero prices are already
obtained by $s=0$.  This proves the claimed closedness and the scale
bound.  The reverse inclusion is immediate from the definition of
prices.
\end{proof}

\paragraph{An alternative derivation of the selector inequality.}
For a spectral point on $V$, restrict it to tensors supported on the
terminal set $P$ and let $\alpha=g_\varphi(P)$.  Power normalization
there gives
\[
 \varphi\!\left(\sum_{a=1}^A T_a\right)
 \le A^\alpha\max_a\varphi(T_a).
\]
When $\alpha=0$ this follows from triviality; otherwise it follows by
the normalized cover inequality.  A unit tensor on $P$ is a restriction
of the star with center $o$ and one equally sized wire to each terminal:
the center imposes equality of all indices and is then contracted.
Consequently
$\alpha\le\sum_{v\in P}d_\varphi(o,v)$.
Together with the cable bounds, this proves
\eqref{eq:selector-price} directly for $d_\varphi$, and homogeneity and
limits give it for all prices.  This also identifies the missing
normalization in the original scale argument of Fact~5.

\subsection{An individual network cannot shorten diameters}

\begin{proposition}[SH-01: a single network lower bound]
\label{prop:single-network}
Suppose a finite bit tensor network, with vertices in the Euclidean unit ball
of $\mathbb R^D$, has output $C_{D,m}$.  Its sum of cable widths times
Euclidean lengths is at least $2Dm$.  The assertion holds over every field.
\end{proposition}
\begin{proof}
For a unit vector $u$ and a real number $t$, split all vertices according to
the sign of $u\cdot x-t$.  Exclude hyperplanes through a vertex, a null set
for integration in $t$.  Let $b(u,t)$ be the total width of cables crossing
this split.  Contracting everything on each side separately expresses the
output matrix across this split as a sum of at most $2^{b(u,t)}$ rank-one
matrices.  Its rank is therefore at most $2^{b(u,t)}$ over any field.

The same flattening of $C_{D,m}$ is, up to tensor factors of rank one, a
tensor product of identity matrices of order $2^m$, one for each coordinate
diameter whose endpoints the hyperplane separates.  If $k(u,t)$ is their
number, its rank is exactly $2^{m k(u,t)}$.  Consequently
$b(u,t)\ge m k(u,t)$.

For fixed $u$, the integral in $t$ of the crossing indicator of a segment
from $x$ to $y$ is $|u\cdot(x-y)|$.  For uniform probability measure
$\sigma$ on the sphere, rotational invariance gives
\[
 \int |u\cdot(x-y)|\,d\sigma(u)=\kappa_D\|x-y\|_2,
 \qquad \kappa_D:=\int|u_1|\,d\sigma(u)>0.
\]
Integrating $b(u,t)\ge m k(u,t)$ first in $t$ and then in $u$, and using
finiteness to interchange sums and integrals, gives
$\kappa_D\operatorname{cost}(N)\ge\kappa_D(2Dm)$.  Division by
$\kappa_D$ proves the claim.
\end{proof}

This proof checks the substantive content of Fact~4 in the prompt.
Hyperplanes that fail to cross any diameter contribute zero; the wording
that every hyperplane through the ball crosses a diameter must not be read
literally for arbitrary offsets.

\subsection{Removing the apparent normalization gap}

\begin{lemma}[SH-02: compression by multiplicity profiles]
\label{lem:profiles}
Let a tensor $T$ on a fixed terminal set $P$ be a sum of the outputs of
$A\ge1$ fixed finite bit networks $N_1,\ldots,N_A$.  Put their vertices in
one common finite set $V$, and let $o\in V$ be a further selector vertex.
For each positive integer $j$, the tensor $T^{\otimes j}$ is a sum of at
most $\binom{j+A-1}{A-1}$ networks on this same vertex set.  Each summand
uses, for some $n_1+\cdots+n_A=j$, $n_a$ copies of $N_a$ and one selector
cable of width $b_j=\lceil j\log_2 A\rceil$ from $o$ to each terminal.
Thus, for any nonnegative edge weights $p$ on $V$,
\[
 \operatorname{cost}_p(N_{\boldsymbol n})
 \le j\max_a\operatorname{cost}_p(N_a)
       +b_j\sum_{v\in P}p(o,v).
\]
\end{lemma}
\begin{proof}
Expand $T^{\otimes j}$ as a sum indexed by words
$s=(s_1,\ldots,s_j)\in\{1,\ldots,A\}^j$.  Partition the words by their
multiplicity vector $\boldsymbol n$.  For a fixed vector, put $n_a$ copies
of network $N_a$ at their original vertex positions.  A binary word of
length $b_j$ can encode all $A^j$ possible words $s$.  The local table at
$o$ sends the same encoded word along all terminal selector cables and
is zero on codes that are invalid or have multiplicity vector different
from $\boldsymbol n$.

At each terminal, the local table decodes $s$ and permutes the $j$ output
blocks: the $r$th occurrence of symbol $a$ is attached to the $r$th copy of
$N_a$.  This is an arbitrary finite local table, as permitted by the
network definition.  No internal vertex of any $N_a$ needs to receive the
selector.  After contraction, the resulting tensor is precisely the sum
of the word terms with profile $\boldsymbol n$.  Each word contributes
once.  In particular, this reasoning is valid in positive characteristic;
it never divides by a multinomial coefficient.  Summing over profiles
recovers $T^{\otimes j}$.  Copies with the same spatial position may be
fused into one local table.  The claimed cost bound follows by addition
of cable widths.
\end{proof}

\begin{theorem}[SH-03: the selector inequality for prices]
\label{thm:selector-price}
Let $C$ be a product of binary wires, with $c_e$ bits on each pair $e$ of
terminals $P$.  Suppose $C=\sum_{a=1}^A T_a$, where $T_a$ is the output of
a finite network $N_a$ on a common set of positions $V$, with
$P\cup\{o\}\subseteq V$.  Every price $p$ on $V$ satisfies
\begin{equation}
\label{eq:selector-price}
 \sum_e c_e p(e)
 \le \max_a\operatorname{cost}_p(N_a)
       +(\log_2 A)\sum_{v\in P}p(o,v).
\end{equation}
\end{theorem}
\begin{proof}
First set $p=d_\varphi$ for a spectral point on $V$.  An output of a
network is a restriction of the product of wires on its cables, with
trivial indices at its internal parties.  Lemma A.2 and monotonicity
therefore give $\varphi(T_a)\le2^{\operatorname{cost}_{d_\varphi}(N_a)}$.
Apply Lemma~\ref{lem:profiles} and cover subadditivity to obtain
\[
 2^{j\sum_e c_e d_\varphi(e)}
 \le \binom{j+A-1}{A-1}
 2^{j\max_a\operatorname{cost}_{d_\varphi}(N_a)
       +b_j\sum_{v\in P}d_\varphi(o,v)}.
\]
Take base-two logarithms, divide by $j$, and let $j\to\infty$.
The logarithm of the profile count is $O_A(\log(j+1))$, and
$b_j/j\to\log_2 A$.  This proves the inequality for $d_\varphi$.
It is homogeneous in $p$, so it also holds for every nonnegative multiple
of $d_\varphi$.  Its two sides are continuous functions of the finitely
many values of $p$.  Definition A.4 therefore proves it for every price.
\end{proof}

\begin{corollary}[SH-04: shortening excludes flatness]
\label{cor:shortening-not-flat}
A strict shortening in the sense of OpenAI's Definition~2.1 (attached
source, Section~2, lines~203--211) yields a finite Euclidean configuration
whose Euclidean metric is not a price.  Consequently the flat hypothesis
of Anthropic's Definition~2.1 cannot hold over that field.  If the claimed
alternative theorem of that manuscript is correct, its contracting
hypothesis then holds.
\end{corollary}
\begin{proof}
Take $V$ to contain all positions in all $A$ shortening layouts, the
terminals $P=\{\pm e_i:1\le i\le D\}$, and $o=0$.  All of these form a
single finite set independent of $j$ in Lemma~\ref{lem:profiles}.
For Euclidean distance, the left side of
\eqref{eq:selector-price} equals $2Dm$, while the right side is at most
$w+2D\log_2 A$.  The strict shortening inequality contradicts
\eqref{eq:selector-price}.  Flatness would assert that this Euclidean
metric is a price, including the required limits and scales.  The final
conditional implication is the contrapositive of the alternative theorem.
\end{proof}

Fact~5 in the prompt is thus insufficient as an obstruction analysis.
Its inequality for an individual, unnormalized spectral point is valid,
but its asserted limitation to scales at most $2D$ disappears after
multiplicity-profile compression.  The contradiction above is homogeneous
and therefore survives arbitrary scaling and the limiting definition of
a price.

\subsection{A converse reduction to a metric rigidity statement}

\begin{proposition}[SH-05: no shortening forces saturated diameters]
\label{prop:saturated-diameters}
Assume the duality statement of Theorem A.6 in the attached Anthropic
manuscript (pp.~20--21).  Fix $D\ge2$.  If no strict shortening exists in
this dimension over the field, then for every finite set $V$ in the Euclidean unit ball
containing $0$ and $\pm e_1,\ldots,\pm e_D$ there is a price $p$ on $V$
such that
\[
 p(x,y)\le\|x-y\|_2\quad(x,y\in V),\qquad
 p(e_i,-e_i)=2\quad(1\le i\le D).
\]
\end{proposition}
\begin{proof}
Write $\ell(x,y)=\|x-y\|_2$ and let $K$ be the set of prices $p$ on $V$
with $0\le p\le\ell$.  This is a nonempty compact set: it contains the
zero price and is closed by Definition A.4, inside a finite product of
compact intervals.  The continuous function
$F(p)=\sum_i p(e_i,-e_i)$ has maximum at most $2D$ on $K$.  Suppose this
maximum is $2D-\varepsilon$, for some $\varepsilon>0$.

For any spectral point $\varphi$ on $V$, put
\[
 M=\max_{x\ne y}\frac{d_\varphi(x,y)}{\ell(x,y)}.
\]
If $M>0$, the scaled metric $d_\varphi/M$ belongs to $K$, so
$F(d_\varphi)\le(2D-\varepsilon)M$.  If $M=0$, all wire values are one
and the comparison below holds without division.

Choose an integer $m\ge4/\varepsilon$.  For each unordered pair
$e=\{x,y\}$ of distinct points, set
\[
 q_e=\left\lceil\frac{m(2D-\varepsilon)}{\ell(e)}\right\rceil,
 \qquad A_e=\langle2^{q_e}\rangle_e.
\]
The target $C=C_{D,m}$ is viewed as a tensor on $V$ by adjoining trivial
parties.  Since $\ell(e)\le2$,
\[
 q_e\ell(e)\le m(2D-\varepsilon)+2
       \le m(2D-\varepsilon/2).
\]
For every spectral point, choose an edge attaining $M$; then
\[
 \log_2\varphi(C)=mF(d_\varphi)
 \le m(2D-\varepsilon)M
 \le\max_e q_e d_\varphi(e)=\log_2\max_e\varphi(A_e).
\]
By Theorem A.6, for any $\eta>0$ and all large $N$, $C^N$ is covered by
at most $2^{\eta N}$ tensors, each a product of $N$ of the $A_e$.
Each cover term is a restriction of such a product, so it has a bit
network realization on $V$ with cost at most
$Nm(2D-\varepsilon/2)$; straight segments remain in the ball.  Take
$\eta=m\varepsilon/(8D)$.  For the resulting identity with $A$ terms,
\[
 w+2D\log_2 A
 \le Nm(2D-\varepsilon/2)+2D\eta N
 =2DmN-\frac14m\varepsilon N<2DmN.
\]
This is a strict shortening for $C_{D,mN}$, a contradiction.
Thus $\max_KF=2D$.  Every antipodal summand is at most two, so a maximizer
has all the asserted equalities.
\end{proof}

\begin{corollary}[SH-06: a saturated price on the whole ball]
\label{cor:ball-price}
Under the hypotheses of Proposition~\ref{prop:saturated-diameters}, there
is a price $p$ on the entire unit ball with $p\le\ell$ and
$p(e_i,-e_i)=2$ for every coordinate diameter.
\end{corollary}
\begin{proof}
In the compact product of intervals
$\prod_{\{x,y\}\subset B^D}[0,\ell(x,y)]$, impose the closed requirements
that each finite restriction be a price and that the listed antipodal
equalities hold.  Every finite collection of requirements concerns a
finite union of points and is satisfied by
Proposition~\ref{prop:saturated-diameters}.  Compactness gives a common
point.  By the definition of an infinite price, it is the asserted price.
\end{proof}

\subsection{Rigidity and the exact relationship with flatness}

\begin{lemma}[SH-07: rigidity of saturated coordinate diameters]
\label{lem:rigid-ball}
Let $f:B^D\longrightarrow X$ be a $1$-Lipschitz map from the Euclidean
unit ball to a CAT(0) space.  Suppose
$d_X(f(e_i),f(-e_i))=2$ for every $i$.  Then $f$ restricts to an isometry
on the ball of radius $2^{-(D-1)/2}$ centered at zero.
\end{lemma}
\begin{proof}
We use the roundness-two inequality and the flat-rectangle property for
CAT(0) spaces, in the precise forms of the attached manuscript's
Lemma C.4(b),(f), p.~29.  The latter is a form of the Flat Quadrilateral
Theorem of Bridson and Haefliger (1999, II.2.11).

First observe that a nonexpansive map of a Euclidean line segment which
preserves its endpoint distance preserves all distances on that segment.
Indeed the triangle inequality, combined with the upper bounds on each
successive subsegment, must be an equality.  In a CAT(0) space its image
is consequently the uniquely parametrized geodesic between the endpoint
images.  It follows that $f$ is isometric on every coordinate-axis
segment $[-e_i,e_i]$.

We will also use the following consequence of the flat-rectangle
property.  If a nonexpansive map of a Euclidean square sends its four
corners to points with the same six mutual distances, then it is
isometric on the square.  For completeness, the reduction to the stated
property is as follows.  The preceding segment argument fixes the
images of each boundary edge.  Lemma C.4(f) applied to opposite boundary
edges gives the Euclidean distances between all their points.  In
particular, each segment joining corresponding points on these edges
has its full endpoint distance preserved, and so is mapped to the
appropriate geodesic.  For two such segments their four endpoint images
have all the distances of a Euclidean rectangle.  Another application
of C.4(f) gives the Euclidean distance between every point on one and
every point on the other.  These segments fill the square.

Write $E_k=\operatorname{span}\{e_1,\ldots,e_k\}$ and
$r_k=2^{-(k-1)/2}$.  We prove inductively that $f$ is isometric on
$E_k\cap B(0,r_k)$.  The case $k=1$ was proved above.  Suppose it holds
for $k-1$ and set $r=r_{k-1}$.  For any unit $u\in E_{k-1}$, consider
the four points $ru,re_k,-ru,-re_k$.  The two diagonal lengths in their
images are $2r$ by the inductive hypothesis and the coordinate-axis
observation.  Each side length is at most $\sqrt2r$ by nonexpansiveness.
The roundness-two inequality reads
\[
 8r^2\le\sum_{\text{four sides}}d_X(f(v),f(w))^2\le8r^2.
\]
Every side therefore has length $\sqrt2r$.  These are all six distances
of a Euclidean square, so the preceding consequence of C.4(f) shows
that $f$ is isometric on
\[
 S_u=\{a u+b e_k:|a|+|b|\le r\}.
\]

Every $z\in E_k\cap B(0,r/\sqrt2)$ has the form $a u+b e_k$ with
$|a|+|b|\le\sqrt2\|z\|_2\le r$, unless its projection on $E_{k-1}$
vanishes, in which case it lies on the already treated coordinate axis.
Thus the diametral distance from $z$ to $-z$ is preserved throughout
this smaller ball.  Given any $x,y$ in it, apply the roundness-two
inequality to the cyclic quadrilateral $x,y,-x,-y$.  Its image has the
sum of squared diagonal lengths
$4\|x\|_2^2+4\|y\|_2^2$.  The sum of its squared side lengths is at
most the Euclidean value, which by the parallelogram identity is the
same number.  Hence each of the four nonexpansive side bounds is an
equality, including $d_X(f(x),f(y))=\|x-y\|_2$.  This proves the
inductive step and the lemma.
\end{proof}

\begin{theorem}[SH-08: shortening is exactly failure of flatness]
\label{thm:shortening-flat-equivalence}
Assume the duality and CAT(0)-embedding statements of the attached
Anthropic manuscript, Theorem A.6 and Proposition C.3, respectively.
Over any fixed field, a strict shortening exists if and only if the flat
hypothesis fails.
\end{theorem}
\begin{proof}
Strict shortening implies failure of flatness by
Corollary~\ref{cor:shortening-not-flat}.  For the converse, suppose no
strict shortening exists.  Fix $D\ge2$.  By
Corollary~\ref{cor:ball-price}, the entire unit ball has a price $p$ with
$p\le\ell$ and all coordinate diameters saturated.  Proposition C.3
represents this price as distances under a map $f$ to a CAT(0) space.
The inequality $p\le\ell$ says that $f$ is $1$-Lipschitz, so
Lemma~\ref{lem:rigid-ball} makes it an isometry on a positive-radius
$D$-dimensional Euclidean ball.  The restriction of $p$ to this ball
is therefore Euclidean distance.  Every finite subset of $\mathbb R^D$
fits into this ball after translation and positive rescaling.  Pullback
and scaling preserve prices, so every such finite subset is a price.
Taking $D$ arbitrarily large gives the flat hypothesis.
\end{proof}

\begin{corollary}[SH-09: the remaining equivalence gate]
\label{cor:exclusivity-gate}
In addition to the hypotheses of Theorem~\ref{thm:shortening-flat-equivalence},
assume the alternative theorem of the attached manuscript
(Theorem D.2, p.~30).  For a fixed field the assertion
\[
 \text{a strict shortening exists}\quad\Longleftrightarrow\quad
 \text{the contracting hypothesis holds}
\]
is equivalent to mutual exclusion of the flat and contracting
hypotheses.  In particular, strict shortening always implies contracting
under these structural claims.
\end{corollary}
\begin{proof}
Theorem~\ref{thm:shortening-flat-equivalence} identifies strict shortening
with the negation of flatness.  Theorem D.2 says that this negation implies
contracting.  The reverse implication is exactly the assertion that
contracting excludes flatness.  This is also precisely mutual exclusion,
since both hypotheses are propositions about the same fixed field.
\end{proof}


\subsection{A failed route: a four-point obstruction for the rank metric}
\label{sec:rank-roundness-search}
The proof of the alternative theorem produces the putative price
$p(A,B)=\sqrt{\operatorname{rank}(A-B)}$ if contraction fails.  A possible
route to contraction would be to show that this metric violates a
necessary CAT(0) inequality.  The integer computation
\texttt{programs/sh\_rank\_roundness.py} exhaustively tested the
roundness-two inequality for this metric on $\operatorname{Mat}_n(\mathbb
F_2)$ for $n=2,3$.  Translation invariance sets the first matrix to zero,
leaving $2^{3n^2}$ ordered triples.  The maximum of
\[
 \operatorname{rank}(C)+\operatorname{rank}(B-D)
 -\operatorname{rank}(B)-\operatorname{rank}(B-C)
 -\operatorname{rank}(C-D)-\operatorname{rank}(D)
\]
was zero in both cases.  Thus this necessary inequality supplies no
obstruction in these two sizes.  This computation neither proves
CAT(0) embeddability nor certifies that the metric is a price.

\subsection{Exact experiments on the two-crossing-wire coefficients}
\label{sec:crossing-computations}

\paragraph{SH-COMP-02: calculation and exact checks.}
We implemented the determinant construction of the attached Anthropic
Proposition B.2 (pp.~22--24) for the whole Frobenius-symbol classes
\[
 \Sigma_{n,E}=\{(I,J):|I|=|J|=n,
       \sum_{i\in I}(2i+1)+\sum_{j\in J}(2j+1)=E\}.
\]
Here the two crossing cables use the same class, of size $q$, and
$\Theta=4E$.  The executable source is
\texttt{programs/sh\_crossing\_ranks.cpp}; the case list is
\texttt{inputs/sh\_crossing\_ranks.tsv}.  The driver and completion
summary are generated by the two accompanying Python programs.

The implementation evaluates the $4n$ by $4n$ matrix in the cited
construction and recovers all determinant coefficients by exact
Vandermonde interpolation.  For audit purposes, it interpolates through
degree $4E+4n$, which is a direct upper bound: each matrix entry has
degree at most $2i+2j+3$, so a determinant term has degree at most
$2\sum i+2\sum j+12n=4E+4n$.  It then verifies that every coefficient
above the claimed degree $\Theta=4E$ is zero.  It also checks the value
at $X=1$ against the specified pair of identity cables, the coefficient
reflection $N\leftrightarrow\Theta-N$ with the corresponding party
swap, and evaluation at one point not used for interpolation.

For characteristic two, interpolation takes place in the explicitly
implemented field $\mathbb F_{256}=\mathbb F_2[z]/(z^8+z^4+z^3+z+1)$.
Every recovered coefficient is checked to lie in $\{0,1\}$ before
the three flattening ranks are calculated.  The other runs use the
prime $1{,}000{,}003$.  A full matrix rank modulo this prime certifies
full matrix rank over $\mathbb Q$, hence over $\mathbb C$, for the
same integral coefficient tensor.  A smaller modular rank alone is
not claimed to determine the characteristic-zero rank.

The saved, completed cases include size-one classes with
$E=4,8,12,16$, size-two classes through $E=20$, and the additional
cases indicated in the completed-output table.  Every completed run
has zero failures in each of the exact checks just described.  In the
size-one cases the balanced coefficients have full adjacent-flattening
rank $q^2$ modulo the odd prime.  Characteristic two does reduce some
ranks: for $n=1,E=8,q=4$ the central coefficient has rank six in both
adjacent cuts, whereas for $n=1,E=16,q=8$ it has rank twenty-eight.
These isolated reductions do not produce a shortening.

\paragraph{SH-12: exact geometry of the tested two-hub layouts.}
For clarity, the geometric test grants ideal cable width $m=\log_2 q$
and bridge width $h=\log_2 r$.  These real widths are an optimistic
asymptotic relaxation of binary widths.  Put the four terminals at
$\pm e_1,\pm e_2$.  For the adjacent split, take symmetric hubs at
$s(e_1+e_2)/\sqrt2$ and its negative.  Their cable cost is
\[
 F(s)=4m\sqrt{1-\sqrt2s+s^2}+2hs.
\]
Minimizing for $s\ge0$ gives
\[
 \min F=
 \begin{cases}
 \sqrt2\bigl(h+\sqrt{4m^2-h^2}\bigr),&0\le h<\sqrt2m,\\
 4m,&h\ge\sqrt2m.
 \end{cases}
\]
Indeed, after writing $a=1/\sqrt2$, the vanishing derivative equation
is $(a-s)/\sqrt{(a-s)^2+a^2}=h/(2m)$.  It has a positive-displacement
solution exactly when $h<\sqrt2m$, at
$s=a(1-h/\sqrt{4m^2-h^2})$.  Substitution gives the first expression;
otherwise convexity gives the minimum at zero.  Thus the strict
first-order rank threshold is $r<q^{\sqrt2}$.

For each nonzero coefficient the experiment uses the cheaper of its
two adjacent ranks.  The completion summary records the worst such
rank, the resulting optimistic hub cost, and the full selection
penalty $4\log_2 A$, where $A$ is the number of nonzero coefficient
tensors.  In every completed case in the listed finite range,
$A\ge q$.  Consequently the selector penalty alone is at least the
direct cost $4\log_2 q$, before any positive network cost is charged.
These finite coefficient decompositions cannot satisfy the strict
shortening inequality.  This is a refutation of those particular
certificates, not a bound on arbitrary larger symbol classes,
regroupings of coefficients, or alternative layouts.

\paragraph{SH-COMP-03: two-layout regrouping and coefficient-span search.}
The preceding selector obstruction was explicitly removed in a second
experiment.  For every cutoff $K$, we formed
$S_K=\sum_{N\le K}\Psi_N$ and $U_K=\sum_{N>K}\Psi_N$.  Their sum is
the same target, but there are only two layouts.  The program computes
both adjacent flattening ranks of each of $S_K,U_K$, so each term may
choose its cheaper adjacent cut independently.  All fourteen cases in
\texttt{inputs/sh\_crossing\_grouped.tsv} completed with the same exact
identity certificates passing.  The cost summary is
\texttt{outputs/sh\_crossing\_grouped\_summary.tsv}.

For the size-one classes over $\mathbb F_2$, the minimum, over all
cutoffs, of the larger of the two terms' cheaper ranks is respectively
$1,6,28,120$ at $q=2,4,8,16$.  It is attained by taking the coefficient
$\Psi_0$ alone versus its complement.  At $q=8,16$ these ranks already
exceed $q^{\sqrt2}$, so even the geometric part of this two-hub routing
cannot save length.  Across all fourteen cases, every nontrivial
two-group candidate has a negative selector-paid margin.  The largest
margin in this finite search is $-2\sqrt2$, attained in the $q=2$
characteristic-two case.  Zero summands, which leave a single network,
are excluded from that maximum and are covered by
Proposition~\ref{prop:single-network}.

We then allowed arbitrary $\mathbb F_2$ linear combinations of the
coefficient tensors, rather than only threshold sums.  The exact binary
span has dimension four for $q=2$ and twelve for $q=4$.  Exhaustive
enumeration of all sixteen and all $4096$ tensors in these respective
spans proves that the best possible larger cheaper-adjacent rank of
$S,C-S$ is one and six, respectively.  For $q=8$ the span has dimension
twenty-seven.  A fixed-seed search and coordinate descent evaluated
$8391$ distinct tensors and found best larger rank twenty-eight; this
last statement is solely a finite search outcome, not an optimum proof.
The program, fixed input, packed coefficient tensors, and complete
best-witness records are included as
\texttt{programs/sh\_crossing\_subsets.py},
\texttt{inputs/sh\_crossing\_subsets.json}, and the corresponding files
in \texttt{outputs}.  Each recorded witness was reconstructed from its
original coefficient degrees and checked against its two flattenings.

\paragraph{A further attempted modification and its precise failure.}
The square-root symbol-count growth in Proposition B.1 explains the
critical $\sqrt2$ bridge exponent.  Replacing the linear mode weights
$2r+1$ by superlinear weights would change the counting exponent, but
it does not preserve the crossing identity.  Even at the constant
polynomial $B(r,0)=1$, replacing $X^{2r+1}$ by $X^{(2r+1)^2}$ while
leaving the Meixner kernel unchanged changes its orthogonality sum to
\[
 (1-X^2)\sum_{r\ge0}X^{(2r+1)^2}
   =X-X^3+X^9-X^{11}+\cdots,
\]
which differs from the required $X$ already in degree three, in every
characteristic.  A new superlinear-weight construction would therefore
need a new orthogonality kernel, an exact reflection identity, and a
new degree accounting argument.  None of those requirements follows
from the existing determinant construction.  This route remains open;
the displayed calculation refutes only the naive weight substitution.

\begin{proposition}[SH-13: width one has no strict shortening]
\label{prop:width-one}
Over any field and in every dimension $D\ge2$, $C_{D,1}$ has no strict
shortening.  More generally, any strict shortening must have
$2\le A<2^m$ and hence $m\ge2$.
\end{proposition}
\begin{proof}
Network cost is nonnegative.  If $A\ge2^m$, the selector term alone
is $2D\log_2 A\ge2Dm$, violating strictness.  If $A=1$,
Proposition~\ref{prop:single-network} rules out the shortening.
For $m=1$ these alternatives exhaust all positive integers $A$.
\end{proof}

\begin{proposition}[SH-14: a finite-width obstruction for the full-width two-hub model]
\label{prop:full-width-obstruction}
Let $q=2^m$, and write the two-crossing-cable tensor $C_{2,m}$ as a sum
of $A$ nonzero tensors.  Suppose each term is realized by the standard
two-hub construction for one of the two adjacent-pair cuts, with
\emph{all four terminal cables of width $m$} and one bridge of some
width $h_j$.  The hub positions may be arbitrary in the unit ball.
A strict shortening by these layouts requires
\[
 \frac{m}{\log_2 A}
 >\frac{1}{1-2/\sqrt{10-4\sqrt2}}>24.
\]
In particular, no such certificate has $m\le24$.  This assertion does
not cover layouts that compress the local terminal supports and use
smaller terminal-cable widths, nor networks of other forms.
\end{proposition}
\begin{proof}
For the adjacent cut $ab\mid cd$, the cost as a function of hub
positions $x,y$ is convex.  It is invariant under
$(x,y)\mapsto(-y,-x)$.  Averaging these two layouts makes the hubs
opposite without increasing cost.  Reflection interchanging $e_1,e_2$
and projection to their coordinate plane then put the hubs on their
bisector, again by convexity and symmetry.  Reversing a negative
bisector displacement decreases the lead cost.  Hence the global
minimum is exactly the function $F(m,h)$ computed in SH-12, not merely
an upper bound supplied by that symmetric family.  The same argument
applies to the other adjacent cut.  This minimum is nondecreasing in
$h$.

Fix the external indices at the two opposite parties $a,c$ to the
same basis vector $e_i$.  The target becomes the $q$ by $q$ identity
matrix on $b,d$.  A summand routed across either adjacent cut with a
bridge of width $h_j$ becomes a matrix of rank at most $2^{h_j}$:
fixing $a,c$ in its cut-rank decomposition leaves a sum of at most
that many rank-one matrices.  Rank subadditivity gives
\[
 q\le\sum_{j=1}^A2^{h_j},\qquad
 \max_j h_j\ge m-\log_2 A.
\]
By SH-13, a strict shortening would have $1<A<q$.  Put
$a=\log_2 A$ and $h=m-a>0$.  Its maximum layout cost is at least
$F(m,h)$, so the strict shortening inequality requires
\[
 \sqrt2\bigl(h+\sqrt{4m^2-h^2}\bigr)+4a<4m,
 \quad\text{or equivalently}\quad
 \sqrt2\bigl(h+\sqrt{4m^2-h^2}\bigr)<4h.
\]
Rearranging the positive quantities and squaring gives
$4m^2<(10-4\sqrt2)h^2$.  Substituting $h=m-a$ yields the stated
necessary condition.  Its right-hand constant exceeds twenty-four,
and $a\ge1$, so $m\le24$ is impossible.
\end{proof}

\begin{corollary}[SH-15: the fixed-dimension equivalence]
\label{cor:fixed-dimension}
Under the structural hypotheses of Theorem~\ref{thm:shortening-flat-equivalence},
for every fixed $D\ge2$ the following are equivalent:
\begin{enumerate}
\item no strict shortening exists in dimension $D$ at any width;
\item every finite Euclidean configuration of dimension at most $D$
is a price.
\end{enumerate}
\end{corollary}
\begin{proof}
The selector inequality excludes a shortening whenever all its
positions carry their Euclidean price.  Conversely, the proof of
Proposition~\ref{prop:saturated-diameters} would produce a shortening
in the same fixed dimension $D$.  Thus absence of such a shortening
gives the saturated ball price of Corollary~\ref{cor:ball-price} in
that dimension.  Lemma~\ref{lem:rigid-ball} and the final scaling
argument in Theorem~\ref{thm:shortening-flat-equivalence} give all
finite Euclidean configurations of dimension at most $D$.
\end{proof}

\subsection{A limitation of cut-mixture spectral constructions}

\begin{proposition}[SH-16: the rank square-root metric is outside $L_1$ in size three]
\label{prop:rank-not-lone}
The metric
$d(A,B)=\sqrt{\operatorname{rank}(A-B)}$ on
$G=\operatorname{Mat}_3(\mathbb F_2)$ is not conditionally negative
definite as a distance kernel.  In particular, it has no isometric
embedding into $L_1$.
\end{proposition}
\begin{proof}
Let $H=\operatorname{diag}(1,1,0)$ and
$\chi(A)=(-1)^{\sum_{ij}H_{ij}A_{ij}}$.  This nontrivial character
satisfies $\sum_A\chi(A)=0$.  Exact enumeration of the $512$ binary
matrices gives the signed counts
\[
 \left(\sum_{\operatorname{rank}A=r}\chi(A)\right)_{r=0}^3
 =(1,1,-10,8).
\]
The executable certificate
\texttt{programs/sh\_rank\_metric\_fourier.py} uses binary Gaussian
elimination and integer addition only; its input and output are
included.  Translation invariance and the character law yield
\[
 \sum_{A,B\in G}\chi(A)\chi(B)d(A,B)
 =512\bigl(1-10\sqrt2+8\sqrt3\bigr)>0.
\]
The positivity is exact: $(1+8\sqrt3)^2=193+16\sqrt3>200$.
This violates the defining condition for a conditionally negative
definite kernel.

To see the stated $L_1$ consequence in the finite form needed here,
every cut metric $\delta_f(A,B)=|f(A)-f(B)|$, $f:G\to\{0,1\}$,
satisfies
\[
 \sum_{A,B}c_Ac_B\delta_f(A,B)
 =-2\left(\sum_Ac_Af(A)\right)^2\le0
 \quad\text{when }\sum_Ac_A=0.
\]
Absolute distance on the real line is an integral of threshold cut
metrics, and an $L_1$ distance is an integral of such coordinate
distances.  Thus every finite $L_1$ metric obeys the same inequality,
whereas the character coefficients above violate it.
\end{proof}

This obstruction does not establish that the rank square-root metric
is not a price: no inclusion of all prices into conditionally negative
definite metrics has been proved here.  Its role is to identify the
limitation of $L_1$ or cut-mixture constructions for the equality metric
arising in the alternative theorem.  A proposed proof of such an
inclusion would be a distinct substantive step, not a consequence of
CAT(0) embeddability alone.

\subsection{A concrete graph-tensor test for conditional negative type}

Let $G_r$ be the graph tensor with one binary wire on every edge of
$K_{r,r}$, with labelled parts $L=\{L_1,\ldots,L_r\}$ and
$R=\{R_1,\ldots,R_r\}$.  Let $H_r$ have one binary wire on every
edge of the two disjoint cliques on $L$ and on $R$.

\begin{proposition}[SH-17: the conditional-negative-type gate]
\label{prop:cnd-graph-gate}
Assume the all-tensor duality in Anthropic Theorem~A.6, pp.~20--21,
and the pullback property for prices in Definition~A.4 and its
immediately following paragraph, p.~20.
The following assertions are equivalent over a fixed field:
\begin{enumerate}
\item every finite price $p$ is conditionally negative definite as a
distance kernel, meaning
$\sum_{i,j}c_ic_jp(i,j)\le0$ whenever $\sum_i c_i=0$;
\item for every $r\ge1$, $H_r^{\otimes N}$ can be covered by
$2^{o(N)}$ copies of $G_r^{\otimes N}$ along an unbounded sequence
of integers $N$.
\end{enumerate}
The second assertion permits arbitrary local restrictions of each
covering copy.  It is not an assertion of a single exact restriction.
\end{proposition}
\begin{proof}
For any spectral point $\varphi$ with wire metric $d=d_\varphi$,
Lemma~A.2 gives
\[
 \log_2\varphi(H_r)=
 \sum_{i<j}d(L_i,L_j)+\sum_{i<j}d(R_i,R_j),\qquad
 \log_2\varphi(G_r)=\sum_{i,j}d(L_i,R_j).
\]
The conditional-negative-type inequality for the coefficients
$+1$ on $L$ and $-1$ on $R$ says exactly that the first expression
does not exceed the second.  Applying Theorem~A.6 with its source
family consisting of the single tensor $G_r$ supplies the covers in
(2): for every positive $\eta$, eventually the cover count is at most
$2^{\eta N}$, and a slow diagonal choice makes $\eta$ tend to zero.

Conversely, the cover inequality and multiplicativity imply
$\varphi(H_r)^N\le A_N\varphi(G_r)^N$.  Taking $N$th roots,
then scaling and passing to limits, proves the displayed wire
inequality for every price on the $2r$ labelled parties.
Let $p$ be any finite price and let integer coefficients $c_i$ sum
to zero.  Make $c_i$ positively labelled copies of point $i$ when
$c_i>0$, and $-c_i$ negatively labelled copies when $c_i<0$.
There are equally many copies of each sign.  Pull back $p$ to
these labels; distinct copies of the same point have distance zero.
The wire inequality becomes
$\frac12\sum_{i,j}c_ic_jp(i,j)\le0$.
Clearing denominators proves this for rational coefficients, and
continuity on the finite-dimensional hyperplane $\sum_i c_i=0$
proves it for real coefficients.
\end{proof}

Proposition~\ref{prop:rank-not-lone} shows why this gate is relevant.
The full-matrix character witness there has $256$ points of each
sign.  Proposition~\ref{prop:eight-point-cnd} below strengthens it
to four points of each sign.  Thus the single assertion in (2)
for $r=4$ would already exclude the rank square-root equality
metric on $\operatorname{Mat}_3(\mathbb F_2)$ from the price cone.
No such cover is constructed here.

\begin{proposition}[SH-18: exact routing threshold]
\label{prop:cnd-routing}
Among restrictions constructed by routing independent wires along
paths, the optimal rate from $G_r$ to $H_r$ is
\[
 \frac{r}{2(r-1)}\qquad(r\ge2).
\]
In particular, $H_2^{\otimes2}\le G_2^{\otimes2}$ over every field.
For $r\ge3$, routing cannot establish the rate-one gate in
Proposition~\ref{prop:cnd-graph-gate}.
\end{proposition}
\begin{proof}
Every demanded wire has endpoints in the same part of the bipartite
source graph, so every routing path uses at least two source edges.
There are $r(r-1)$ target edges and $r^2$ source edges.  At target
rate $\rho$, summing all source edge loads gives
$2r(r-1)\rho\le r^2$, proving the upper bound.

For the matching construction, take source width $2(r-1)$ and
target width $r$.  For each target edge in $L$, route one of its
$r$ bits through each of the $r$ possible intermediate vertices
in $R$.  Do the same with the roles exchanged for target edges
in $R$.  Each source edge carries $r-1$ bits for the first group
of demands and $r-1$ for the second.  The local maps forward
endpoint indices and impose equality on the two halves of every
path at its intermediate vertex.  They have entries zero or one,
and contraction gives each target assignment exactly once.
This proves the exact restriction
$H_r^{\otimes r}\le G_r^{\otimes 2(r-1)}$ over every field.
\end{proof}

\begin{proposition}[SH-19: the first nontrivial direct restriction fails]
\label{prop:k33-no-direct}
Let $q>1$ be an integer and work over any field.  Write $G_3(q)$
and $H_3(q)$ for the preceding graph tensors with wire dimension
$q$ in place of two.  Then
\[
 H_3(q)\not\le G_3(q).
\]
Consequently there is no exact rate-one restriction
$H_3^{\otimes N}\le G_3^{\otimes N}$ for any $N\ge1$.
This statement does not rule out covers with subexponentially
many terms or restrictions with asymptotically negligible extra
source width.
\end{proposition}
\begin{proof}
Suppose local maps $A_i$ at $L_i$ and $B_i$ at $R_i$ give the
claimed restriction.  Group $L_i$ and $R_i$ as one of three
superparties.  Contract their internal source edge, of dimension
$q$, into their two local maps.  The resulting map $F_i$ has
domain and codomain of dimension $q^4$.  On the remaining edges,
the grouped source and target are both triangle graph tensors
with edge dimension $q^2$, hence both are matrix-multiplication
tensors $\langle q^2,q^2,q^2\rangle$, up to local basis
permutations.  Their flattening rank at each superparty is $q^4$.
Therefore every $F_i$ is invertible.

Use the order $(L\text{-half},R\text{-half})$ for the two
$q$-dimensional factors in each grouped edge.  Its source pairing
is the swap matrix $P$ on $\mathbb F^q\otimes\mathbb F^q$,
whereas its target pairing is the identity matrix.  Turning each
source pairing into the identity by swapping its factors at one
endpoint gives a usual matrix-multiplication tensor.  The
fixed-party isotropy theorem says that each resulting local
automorphism factors as a tensor product over its two incident
edges.  The exact result used is Burichenko (2016),
\emph{The isotropy group of the matrix multiplication tensor},
Proposition~5, p.~5 of the author reprint
\texttt{arXiv:2210.16565v1}, with proof in \S8, pp.~12--13;
its ground field is arbitrary.  Undoing the endpoint swaps
preserves this edgewise factorization.  Thus, after harmless
nonzero scalar rescaling of factors, there are invertible
$q^2\times q^2$ matrices $V_{ij}$ such that
\[
 F_i=\bigotimes_{j\ne i}V_{ij},\qquad
 V_{ij}P V_{ji}^{\mathsf T}=I\quad(i<j).
 \tag{\(*\)}
\]

For an operator on $U\otimes W$, define its operator Schmidt rank
$\operatorname{osr}$ to be its tensor rank in
$\operatorname{End}(U)\otimes\operatorname{End}(W)$, equivalently
the ordinary rank of that two-factor flattening.  Contracting
$A_i$ and $B_i$ over their shared $q$-state index expresses
$F_i$ as a sum of $q$ product operators across its physical
$L_i:R_i$ split.  Hence
\[
 \operatorname{osr}(F_i)\le q.
\]
On the other hand, the first equality in $(*)$ and multiplicativity
of matrix rank under tensor product give
\[
 \operatorname{osr}(F_i)=\prod_{j\ne i}
 \operatorname{osr}(V_{ij}).
\]
The swap matrix has operator Schmidt rank $q^2$, since
$P=\sum_{a,b}E_{ab}\otimes E_{ba}$ has an invertible
$q^2\times q^2$ operator flattening.  The second equality in
$(*)$ implies $P=V_{ji}^{\mathsf T}V_{ij}$.  Operator Schmidt
rank is invariant under full transpose and submultiplicative
under composition, as is seen by multiplying product-operator
expansions.  It follows that
\[
 q^2\le \operatorname{osr}(V_{ji})\operatorname{osr}(V_{ij})
 \quad(i<j).
\]
Multiplication over the three edges contradicts the local bounds:
\[
 q^6\le\prod_{i<j}
 \operatorname{osr}(V_{ij})\operatorname{osr}(V_{ji})
 =\prod_{i=1}^3\operatorname{osr}(F_i)\le q^3.
\]
All ranks and factorizations in this argument are over the
specified field, so the obstruction applies in characteristic two
as well as in characteristic zero.
\end{proof}

The gap between Proposition~\ref{prop:k33-no-direct} and
Proposition~\ref{prop:cnd-graph-gate} is essential: an asymptotic
cover is a sum of restricted source tensors.  Neither an individual
term's full local rank nor the invertibility of its grouped maps
follows from the full local rank of that sum.

\begin{proposition}[SH-20: an eight-point obstruction, minimal within its support]
\label{prop:eight-point-cnd}
In $\operatorname{Mat}_3(\mathbb F_2)$ put
\[
 A=E_{11},\qquad B=E_{22},\qquad
 C=\begin{pmatrix}0&0&1\\0&0&1\\1&1&0\end{pmatrix},
 \qquad V=\operatorname{span}_{\mathbb F_2}\{A,B,C\}.
\]
The eight-point metric $d(X,Y)=\sqrt{\operatorname{rank}(X-Y)}$
on $V$ is not conditionally negative definite.  Every proper subset
of this same eight-point metric is strictly conditionally negative
definite: its quadratic form is negative for every nonzero real
coefficient vector with sum zero.  The latter assertion does not
claim that no smaller witness exists elsewhere in the full matrix space.
\end{proposition}
\begin{proof}
The matrices are linearly independent.  Direct binary row reduction,
in the order $0,A,B,A+B,C,A+C,B+C,A+B+C$, gives ranks
\[
 (0,1,1,2,2,3,3,2).
\]
For the character $\chi(xA+yB+zC)=(-1)^{x+y+z}$, the signed
rank-stratum counts are therefore $(1,-2,-1,2)$.  Consequently
\[
 \sum_{X,Y\in V}\chi(X)\chi(Y)d(X,Y)
 =8(2\sqrt3-\sqrt2-2)>0.
\]
The last inequality follows from
$12>(2+\sqrt2)^2=6+4\sqrt2$, since $6>4\sqrt2$.

For the proper-subset assertion, identify $V$ with $\mathbb F_2^3$
in the displayed basis.  Translation acts by isometries, so it is
enough to consider coefficients $c_x$ with $c_0=0$ and
$\sum_xc_x=0$.  The eigenvalues of the translation-invariant
distance matrix in the seven nonconstant character directions are
\[
 -\sqrt2\quad\text{(four times)},\qquad
 -2\sqrt3-2+3\sqrt2,\qquad
 2-2\sqrt3-\sqrt2,\qquad
 \lambda_+=2\sqrt3-\sqrt2-2.
\]
Each of the six negative eigenvalues is less than $-1$.
For the only less immediate one this is
$2\sqrt3+1>3\sqrt2$, obtained by squaring and using
$\sqrt3>5/4$.  Moreover
$0<\lambda_+<1/6$, since $\sqrt3<7/4$ and $\sqrt2>4/3$.

Let $\widehat c(s)=\sum_xc_x(-1)^{s\cdot x}$.
The coefficient assumptions give $\widehat c(0)=0$ and
$\sum_{s\ne0}\widehat c(s)=0$.  If $t$ is the coefficient in
the $\lambda_+$ direction, the other six coefficients sum to
$-t$, so their squared sum is at least $t^2/6$.  Fourier
diagonalization now gives
\[
 \sum_{x,y}c_xc_yd(x,y)
 =\frac18\sum_{s\ne0}\lambda_s\widehat c(s)^2<0
\]
for every nonzero $c$: when $t\ne0$ the upper bound
$(\lambda_+-1/6)t^2/8$ is strictly negative, and when $t=0$
all nonzero Fourier coordinates have negative eigenvalues.
Every proper subset omits at least one point, which can be
translated to zero, proving the claim.

The complete exact arithmetic certificate is
\path{programs/sh_cnd_eight_verify.py}, with its matching input
and output files.  The preceding exploratory linear-subspace
search is separately preserved in
\path{programs/sh_cnd_small_witness.py}; absence of a hit in any
random-search sample is not used in the proof.
\end{proof}

\begin{proposition}[SH-21: the direct obstruction for every larger bipartition]
\label{prop:krr-no-direct}
For every $r\ge3$, every integer $q>1$, and every field,
$H_r(q)$ is not a restriction of $G_r(q)$.
\end{proposition}
\begin{proof}
We first record the required extension of the fixed-party isotropy
statement.  A local automorphism of a complete-graph tensor on
$r\ge3$ parties factors at every party over its incident edges.
To see this, fix an edge $ij$ and group all other parties into one.
The tensor becomes a rectangular matrix-multiplication tensor,
apart from a fixed nonzero tensor on the edges internal to the
third group.  The product of the local maps in that group
preserves its tensor support, because the maps at $i,j$ are
invertible and the global tensor is fixed.  Its restriction to
that support is therefore invertible.  Burichenko's
Proposition~5, with the rectangular dimensions, forces the map
at $i$ to factor between edge $ij$ and the bundle of its other
edges.  Doing this for every $j$ forces full edgewise
factorization: a nonzero tensor having rank one across each
single-factor-versus-rest flattening is decomposable, by
successive extraction of those factors.  This argument uses
neither division by the wire dimension nor a characteristic-zero
assumption.

Now group $L_i,R_i$ in the proposed restriction exactly as in
Proposition~\ref{prop:k33-no-direct}.  Each contracted map $F_i$
is square of dimension $q^{2(r-1)}$, and full flattening rank
makes it invertible.  The grouped source and target are
complete-graph tensors with edge dimension $q^2$; source edge
pairings are swaps and target pairings are identities.  Undoing
one swap per edge and applying the preceding isotropy statement
gives $F_i=\bigotimes_{j\ne i}V_{ij}$ and
$P=V_{ji}^{\mathsf T}V_{ij}$, up to irrelevant nonzero scalars.
The same operator-Schmidt-rank bounds as before now imply
\[
 q^{r(r-1)}
 \le\prod_{i<j}\operatorname{osr}(V_{ij})
                     \operatorname{osr}(V_{ji})
 =\prod_i\operatorname{osr}(F_i)\le q^r,
\]
which is impossible when $r\ge3$ and $q>1$.
\end{proof}

\subsection{A quantitative consequence of the single eight-party gate}

\begin{lemma}[SH-22a: the finite linear equality system]
\label{lem:linear-rank-rigidity}
Let $u(a,b)=u(b,a)$ be a real array on
$M=\operatorname{Mat}_3(\mathbb F_2)$, with $u(a,a)=0$.
For every transversal pair $W,W'$ of ranks $(1,1)$ or $(1,2)$,
including the reversed ordering, suppose that in the quadrilateral
$a,a+W,a+W+W',a+W'$ opposite $u$-values are equal and the
sum of the four side values equals the sum of the two diagonal
values.  Then $u(a,b)=c\operatorname{rank}(a-b)$ for one
real constant $c$.
\end{lemma}
\begin{proof}
This is the purely linear form of Anthropic's Theorem~D.2,
Steps~4--7, pp.~30--31.  The following verification makes clear
that no positivity or metric assumption is needed in this form.
Put $n_a(W)=u(a,a+W)$ and
$B_a(W,W')=(n_a(W+W')-n_a(W)-n_a(W'))/2$.
The hypotheses give, for transversal pairs,
\[
 n_{a+W'}(W)=n_a(W),\qquad
 u(a+W,a+W')=2n_a(W)+2n_a(W')-n_a(W+W').
 \tag{L1}
\]
Apply the second identity at $a+W_2$ to
$(W_1,W_2+W_3)$, where $W_1,W_2,W_3$ are transversal
rank-one matrices, and substitute both identities at $a$.
It gives
\[
 n_a(W_1+W_2+W_3)-\sum_{i<j}n_a(W_i+W_j)
                  +\sum_i n_a(W_i)=0,
\]
or
$B_a(W_1,W_2+W_3)=B_a(W_1,W_2)+B_a(W_1,W_3)$.

Fix rank-one $W$ and write
$\xi(v,x)=B_a(W,vx^{\mathsf T})$ for matrices transversal
to $W$.  Comparing the two transversal decompositions
\[
 v_1x_1^{\mathsf T}+v_2x_2^{\mathsf T}
 =v_1(x_1+x_2)^{\mathsf T}+(v_1+v_2)x_2^{\mathsf T}
\]
gives
$\xi(v_1,x_1)-\xi(v_1,x_1+x_2)
 =\xi(v_1+v_2,x_2)-\xi(v_2,x_2)$.
Replacing $x_1$ by $x_1+x_2$ changes the left side's sign
and leaves the right side unchanged, so both vanish.  The graph
of row vectors outside the one-dimensional row space of $W$,
joining vectors independent modulo that space, is the connected
complete tripartite graph with two vertices in each class.
The same holds for column vectors.  Thus $B_a(W,V)$ is
constant on rank-one $V$ transversal to $W$.  Any two rank-one
matrices have a third transversal to both; symmetry of $B_a$
therefore gives one value $\lambda(a)$ on all transversal
rank-one pairs.

Symmetry of $u$ and (L1) give
$B_{a+V}(V,W')=-B_a(V,W')$ for every transversal pair.
For ranks $(1,1)$ this says
$\lambda(a+V)=-\lambda(a)$.
For ranks $(2,1)$ the additivity just proved gives the same
identity.  A rank-two matrix is also a sum of two transversal
rank-one matrices, whose two sign changes instead give
$\lambda(a+V)=\lambda(a)$.  Hence $\lambda=0$.
It follows that $n_a$ is additive on all transversal sums of
total rank at most three.

Apply the same two-decomposition exchange, now to $n_a$ itself.
It proves that $n_a(vx^{\mathsf T})$ is unchanged when one
nonzero row or column vector is replaced by another.  Thus it
has one value $c(a)$ on all rank-one matrices.  The first
identity in (L1), using a rank-one matrix transversal to any
given rank-one translation, makes $c(a)$ invariant under all
rank-one translations.  Those translations span $M$, so $c(a)$
is a constant $c$.  Transversal rank-one decomposition of an
arbitrary matrix gives $n_a(W)=c\operatorname{rank}W$.
\end{proof}

\begin{proposition}[SH-22b: an explicit conditional contraction factor]
\label{prop:cnd-explicit-contraction}
Assume prices embed in CAT(0) spaces as in Anthropic
Proposition~C.3, p.~28.  Suppose that every price on eight
labelled parties satisfies the conditional-negative-type
inequality with four coefficients $+1$ and four coefficients
$-1$.  Then every price on $\operatorname{Mat}_3(\mathbb F_2)$
satisfies
\[
 D(3)^2\le \left(3-2^{-523336}\right)D(1)^2.
\]
In particular, under Theorem~A.6 the single subexponential-cover
assertion in Proposition~\ref{prop:cnd-graph-gate} for $r=4$
implies the contracting hypothesis with the explicit parameters
$\mu=3$ and $s=\sqrt{3-2^{-523336}}$.
\end{proposition}
\begin{proof}
If $D(1)=0$, every rank-one step has distance zero and their
connectivity makes the entire pseudometric zero.  Otherwise
scale to $D(1)^2=1$ and put $\delta=3-D(3)^2\ge0$.
The last inequality is the averaged CAT(0) quadrilateral
inequality, as in Anthropic Theorem~D.2, Step~1, p.~30.
It also gives
\[
 0\le2-D(2)^2\le\delta,\qquad
 0\le1+D(2)^2-D(3)^2\le\delta.
 \tag{L2}
\]
If $\delta\ge1$ the desired estimate already holds, so assume
$0\le\delta\le1$.

For a transversal quadrilateral let $\Delta$ be its squared
side sum minus its squared diagonal sum.  The mean of $\Delta$
in each of the rank types $(1,1)$ and $(1,2)$ is at most
$2\delta$ by (L2).  There are fewer than
$512^3=2^{27}$ ordered choices of base point and two matrices
in either family.  Since $\Delta\ge0$ individually,
\[
 \Delta\le2^{28}\delta
 \quad\text{for every such quadrilateral}.                 \tag{L3}
\]
There are $49$ rank-one matrices and hence $512\cdot49/2=12544$
undirected rank-one edges.  Their squared lengths sum to $12544$,
so each has length at most $112$.  A difference of rank at most
three is a sum of at most three rank-one differences.  The
diameter of the price is consequently at most $336$.

The strong quadrilateral inequality in Anthropic
Lemma~C.4(b), p.~29, is the CAT(0) Reshetnyak inequality.
It implies that the squared difference between the lengths of
either pair of opposite sides is at most $\Delta$.
Thus the difference between their squared lengths has absolute
value at most
$672\sqrt\Delta\le672\,2^{14}\sqrt\delta
 <2^{28}\sqrt\delta$.
The squared-side-sum equality itself has residual
$\Delta\le2^{28}\sqrt\delta$ by (L3).

There are
$N=\binom{512}{2}=130816$ unordered-pair variables $u(a,b)$.
Take all the homogeneous integer linear equations of
Lemma~\ref{lem:linear-rank-rigidity}, and add the equation
that the sum on the $12544$ rank-one edges is $12544$.
The lemma shows that this affine system has the unique solution
$u(a,b)=\operatorname{rank}(a-b)$.  Select $N$ linearly
independent rows, forming a square integer matrix $K$.
Every homogeneous row has Euclidean norm at most $\sqrt6$;
the single normalization row has norm $112$.
Since $\det K$ is a nonzero integer, Cramer's rule and
Hadamard's determinant bound give
\[
 \|K^{-1}\|_{\infty\to\infty}
 \le112N6^{N/2}<2^{2N}.
\]
Apply this to the actual squared distances, whose normalization
residual is zero and whose other residuals have just been
bounded.  We obtain the explicit stability estimate
\[
 \max_{a\ne b}
 \bigl|d(a,b)^2-\operatorname{rank}(a-b)\bigr|
 \le2^{2N+28}\sqrt\delta.                               \tag{L4}
\]

Suppose $\delta<2^{-(4N+72)}$.  Then (L4) is less than
$1/256$.  Since every nonzero rank is at least one,
$|\sqrt x-\sqrt r|\le|x-r|$ for $x\ge0,r\ge1$,
so the distance errors themselves are less than $1/256$.
On the eight matrices of Proposition~\ref{prop:eight-point-cnd},
the parity-character quadratic form for the rank square-root
metric equals $8\lambda_+$, where
\[
 \lambda_+=2\sqrt3-\sqrt2-2>1/25.
\]
For this last bound use $\sqrt3>173/100$ and
$\sqrt2<142/100$, verified by squaring.  Perturbing each of
the $56$ off-diagonal distances by less than $1/256$ changes
the quadratic form by less than $56/256$.  It remains positive,
because $8/25>56/256$.  This contradicts the assumed eight-party
inequality.  Therefore
$\delta\ge2^{-(4N+72)}=2^{-523336}$, as asserted.
\end{proof}

The numerical factor is deliberately coarse.  Its purpose is to
make the consequence of the finite graph-cover gate effective;
it is not a proof that the gate holds.  The proof uses only the
specified eight-party inequality, not a general assertion that
all price metrics have conditional negative type.

\subsection{Unconditional exclusions at small cable widths}

\begin{lemma}[SH-23a: geometric bounds for sign cuts]
\label{lem:sign-cut-cost}
Let $N$ be any finite tensor network in the Euclidean unit ball,
with terminals $\{\pm e_i:1\le i\le D\}$ and integer cable
widths $b_e\ge0$.  Write $c(N)=\sum_e b_e\ell_e$ for its cost.
For each sign vector $s\in\{\pm1\}^D$, modulo simultaneous
negation, let $k_s$ be the minimum cable capacity of a cut
separating $\{s_i e_i\}_i$ from $\{-s_i e_i\}_i$.
There are $M=2^{D-1}$ such cuts.  Then
\[
 c(N)\ge\sqrt2\,k_s\quad\text{for every }s,
 \qquad
 c(N)\ge\frac2M\sum_s k_s.                         \tag{S1}
\]
If $S$ is the tensor computed by $N$, its flattening rank on
the terminal partition specified by $s$ is at most $2^{k_s}$.
\end{lemma}
\begin{proof}
Minimum cut capacity is an integer because every $b_e$ is an
integer.  The maximum-flow/minimum-cut theorem supplies
nonnegative weights on paths between the two terminal sets,
with total weight $k_s$ and total load on edge $e$ at most $b_e$.
One may attach a source and a sink by sufficiently large
capacity edges to formulate the standard two-terminal theorem.
After deleting these additional edges, all the paths connect
opposite terminal sets.  Every such pair of terminals has
Euclidean distance at least $\sqrt2$.  Thus
\[
 \sqrt2\,k_s
 \le\sum_P f_P\operatorname{length}(P)
 \le\sum_e b_e\ell_e=c(N).
\]
This uses precisely the weighted-path formulation of
Ford and Fulkerson (1956), Theorem~1, printed p.~400;
integrality of the flow itself is unnecessary.

For the second inequality put $u_s=s/\sqrt D$.  Projection
onto $u_s$ sends the two terminal sets to $\pm1/\sqrt D$.
Every hyperplane $\langle u_s,x\rangle=t$ with
$-1/\sqrt D<t<1/\sqrt D$, except the finitely many thresholds
through network vertices, cuts cable capacity at least $k_s$.
Integrating the crossing count in $t$ yields
\[
 \sum_e b_e|\langle u_s,\Delta_e\rangle|
 \ge\frac2{\sqrt D}k_s.
\]
Here $\Delta_e$ is the displacement between the endpoints of
edge $e$; straight edges suffice, and a curved edge can only
increase its length.  The sign vectors form a tight frame:
$\sum_s u_su_s^{\mathsf T}=(M/D)I$.
Cauchy--Schwarz consequently gives
\[
 \sum_s|\langle u_s,v\rangle|
 \le\sqrt{M\sum_s\langle u_s,v\rangle^2}
 =\frac M{\sqrt D}\|v\|.
\]
Summing the projection inequalities proves the second
inequality in (S1).  Finally, cutting a set of cables of total
width $k_s$ expresses the terminal flattening as a sum of at
most $2^{k_s}$ rank-one matrices.  This proves its rank bound
over every field.
\end{proof}

\begin{proposition}[SH-23b: no strict shortening at width at most three]
\label{prop:no-small-width-shortening}
Over every field, a strict shortening of $C_{D,m}$ has $m\ge4$.
More generally, if a strict shortening uses $A\ge2$ terms of
cost at most $w$, then
\[
 w\ge\sqrt2\,\big\lceil Dm-\log_2 A\big\rceil,
 \qquad
 \frac{m}{\log_2 A}>
 \frac{2D-\sqrt2}{D(2-\sqrt2)}
 \ge\frac{3+\sqrt2}{2}.                              \tag{S2}
\]
\end{proposition}
\begin{proof}
Across every sign cut the target flattening is an identity
matrix of rank $2^{Dm}$.  Rank subadditivity and
Lemma~\ref{lem:sign-cut-cost} therefore give
\[
 2^{Dm}\le\sum_{a=1}^A\operatorname{rank}_s(T_a)
 \le A2^{\lfloor w/\sqrt2\rfloor}.
\]
This implies the first inequality in (S2).  Combining its
weaker version $w\ge\sqrt2(Dm-\log_2 A)$ with
$w<2D(m-\log_2 A)$ gives the strict ratio bound.
Its right side is increasing in $D\ge2$, with value
$(3+\sqrt2)/2$ at $D=2$.

A single term is excluded by the Crofton bound
of Proposition~\ref{prop:single-network}.
For $A\ge2$, (S2) excludes $m\le2$.  If $m=3$,
it also excludes $A\ge3$, since
$3/\log_2 3<2<(3+\sqrt2)/2$.
It remains to exclude $m=3$, $A=2$.  Strict shortening
would require $w<4D$.  For every term and sign cut,
\[
 k_s\le\lfloor w/\sqrt2\rfloor\le3D-1,
\]
because $2\sqrt2D<3D$.  The two term ranks on each such
cut are consequently at most $2^{3D-1}$.  Their sum must
reach the target rank $2^{3D}$, so both ranks must equal
$2^{3D-1}$.  Every term thus has $k_s\ge3D-1$ for every
sign cut.  The averaged bound in (S1) gives
\[
 c(N)\ge2(3D-1)=6D-2>4D\qquad(D\ge2),
\]
a contradiction.
\end{proof}

\begin{lemma}[SH-24a: a planar Crofton rank bound]
\label{lem:planar-crofton-profile}
Let $S$ be a nonzero tensor on four terminals
$a=e_1,b=e_2,c=-e_1,d=-e_2$, each of dimension $q$.
Write $r_i$ for its four singleton flattening ranks and
$r_1,r_2$ for the ranks of the cuts $ab\mid cd$ and
$ad\mid bc$.  Every network computing $S$ has cost at least
\[
 (\sqrt2-1)\sum_{i\in\{a,b,c,d\}}\log_2 r_i
 +(2-\sqrt2)(\log_2 r_1+\log_2 r_2).                 \tag{S3}
\]
Its integer minimum cut capacities on the two adjacent cuts
also satisfy $k_1+k_2\le c(N)$.
\end{lemma}
\begin{proof}
Apply the same Crofton integration as in
Proposition~\ref{prop:single-network}, now retaining
the actual flattening rank on each terminal cut.  The only
nontrivial terminal partitions induced by a line are the
four singleton cuts and the two adjacent balanced cuts.
By symmetry, their respective Crofton weights are constants
$a_0$ and $b_0$.  An antipodal segment is separated by two
cuts of each type, whereas an adjacent segment is separated
by two singleton cuts and one balanced cut.  The Crofton
normalization therefore gives
\[
 2a_0+2b_0=2,\qquad 2a_0+b_0=\sqrt2.
\]
Thus $a_0=\sqrt2-1$ and $b_0=2-\sqrt2$, proving (S3).
Internal vertices do not change these terminal cut classes
or their measures.  The last assertion is (S1) at $D=2$,
where $M=2$.
\end{proof}

\begin{proposition}[SH-24b: the remaining width-four parameters]
\label{prop:width-four-shortening}
Over every field there is no planar strict shortening with
$m=4$.  Any strict shortening with $m=4$ must have exactly
$A=2$ terms and dimension $D\ge3$.
\end{proposition}
\begin{proof}
First take $D=2,A=2$, so the putative cost bound is $w<12$.
Put $q=16$ and write the target as $T=S+U$.
The integer capacities of either term on the two adjacent
cuts sum to at most $11$.  On each such cut the target has
rank $q^2=256$, so at least one term has capacity at least
$7$.  The same term cannot have capacity at least $7$ on
both cuts.  Thus, after exchanging $S,U$ and the cut names,
$S$ has its high capacity on the first cut and $U$ on the
second; their other capacities are at most $4$.

On either cut, the high-rank term consequently has rank
at least $256-16=240>128$, and hence capacity at least $8$.
The other capacity of that same term is therefore at most
$3$.  It follows that the two low flattening ranks are at
most $8$, and the two high ranks are at least $248$.
Fixing the two opposite indices $a=c=e_j$ in
$T_{abcd}=\delta_{ac}\delta_{bd}$ leaves the $16\times16$
identity matrix on $b,d$.  The ranks of the corresponding
slices of $S$ and $U$ are bounded by their respective low
adjacent flattening ranks.  Rank subadditivity forces the
sum of those two low ranks to be at least $16$.
Both low ranks must therefore equal $8$.

Each high flattening rank is at least $248>16\cdot15$.
A flattening rank is bounded by the product of the two
singleton ranks on either side.  Since every singleton
rank is at most $16$, all four singleton ranks of both
$S$ and $U$ must equal $16$.  Formula (S3) now gives,
even using only the weaker high-rank lower bound $128$,
\[
 c(N)\ge16(\sqrt2-1)+10(2-\sqrt2)
       =4+6\sqrt2>12,
\]
a contradiction.

It remains to dispose of other term counts when $m=4$.
The single-term case was already excluded, and (S2)
excludes $A\ge4$.  Suppose $A=3$ and put $a=\log_2 3$.
The necessary continuous inequality is
\[
 4D(2-\sqrt2)>(2D-\sqrt2)a.                         \tag{S4}
\]
Since $3^{12}>2^{19}$, we have $a>19/12$.
For $D\ge3$, the left side minus the right side in (S4)
is strictly less than
\[
 D(29/6-4\sqrt2)+19\sqrt2/12
 \le(174-125\sqrt2)/12<0.
\]
Here both inequalities follow directly by squaring the
positive rational comparisons.  Thus $D\ge3$ is impossible.
For $D=2$, strict shortening requires
$w<16-4a<29/3<7\sqrt2$.  Lemma~\ref{lem:sign-cut-cost}
bounds every term rank on each sign cut by $2^6=64$.
Three terms cannot sum to a rank-$256$ target.  This
excludes $A=3$ as well and proves the proposition.
\end{proof}

\paragraph{A numerical search for a border restriction.}
The exact obstruction of Proposition~\ref{prop:k33-no-direct}
does not exclude a border restriction.  As an exploratory
test over $\mathbb C$, we optimized the normalized overlap
of the grouped $K_{3,3}$ output with the two-triangle target
at wire dimension $q=2$.  Each grouped local map was
parametrized as a sum of two operator product terms, so its
physical operator Schmidt rank was at most two.  The
program \path{programs/sh_border_k33_search.py}, with input
\path{inputs/sh_border_k33_search.json}, used eight complex
random restarts and at most $2500$ L-BFGS iterations each.
The analytic real gradient was checked on twelve coordinates
against finite differences; the maximum error was less than
$4\cdot10^{-12}$.  The best attained normalized squared
overlap was approximately $0.25$.  Complete iteration
summaries and the best parameter vector are retained in
\path{outputs/sh_border_k33_search_summary.json} and
\path{outputs/sh_border_k33_search_best.npz}.

This is a local numerical search, not a certificate of a
global optimum.  Some restarts reached the iteration limit,
and the homogeneous parametrization permits large parameter
norms that can weaken a gradient stopping criterion.  The
search produced neither an exact identity nor a border curve,
and supplies no new obstruction to subexponential covers.

\subsection{A finite certificate for every dimension at width four}

\begin{lemma}[SH-25a: rank profiles on the folded sign cube]
\label{lem:folded-rank-profiles}
Suppose $C_{D,m}=S+U$, with both tensors nonzero and networks
for each term of cost less than $2D(m-1)$.  For a sign cut $s$
put
\[
 r_V(s)=\operatorname{rank}_s(V),\qquad
 \kappa_V(s)=\lceil\log_2r_V(s)\rceil
 \quad(V=S,U).
\]
Identify opposite sign vectors, and let
$h(s,t)=\min\{\operatorname{Ham}(s,t),
D-\operatorname{Ham}(s,t)\}$ be their folded Hamming distance.
Writing $M=2^{D-1}$, these integer profiles satisfy
\begin{align}
 0\le\kappa_V(s)&\le Dm,&
 \sum_s\kappa_V(s)&\le D(m-1)M-1,                 \tag{R1}\\
 |\kappa_V(s)-\kappa_V(t)|&\le2m h(s,t),          \tag{R2}\\
 r_S(s)+r_U(t)&\ge2^{m(D-h(s,t))}.               \tag{R3}
\end{align}
For $m\ge3$, after exchanging the terms and translating
the folded cube, one may also assume
\[
 \kappa_S(0)\le(m-2)D-1.                         \tag{R4}
\]
\end{lemma}
\begin{proof}
Every sign flattening has row and column dimension $2^{Dm}$.
Its log-rank ceiling is at most its integer network minimum
cut capacity.  Formula (R1) therefore follows from (S1) and
the strict cost bound.

Moving a tensor factor of dimension $2^m$ across a flattening
can increase matrix rank by a factor of at most $2^m$:
expand each rank-one matrix term in a basis of the moved
factor.  Changing the sign on one axis moves two such factors,
so the rank ratio in either direction is at most $2^{2m}$.
This proves (R2), including after taking ceilings.

Choose orientations of $s,t$ that disagree on exactly $h(s,t)$
axes.  On each disagreeing diameter fix both endpoint indices
to the same basis value.  The remaining target is an identity
flattening on $D-h(s,t)$ diameters.  The two restricted term
ranks are bounded respectively by $r_S(s)$ and $r_U(t)$,
and their terminal bipartitions now agree.  Rank subadditivity
proves (R3).

Finally (R1) gives a cut at which
$\kappa_S+\kappa_U\le2D(m-1)-1$.  On this same cut the
target rank is $2^{Dm}$.  If both ceilings are at most $Dm-1$,
both must equal $Dm-1$ to allow their two rank bounds to sum
to $2^{Dm}$.  But their sum $2Dm-2$ exceeds
$2D(m-1)-1$ for $D\ge2$.  Thus one ceiling is $Dm$,
and the other is at most $(m-2)D-1$.  Exchange the terms
and apply a sign translation to put the latter cut at zero.
\end{proof}

\begin{proposition}[SH-25b: width four in dimensions three through five]
\label{prop:width-four-three-five}
There is no strict shortening with $m=4$ and $3\le D\le5$,
over any field.
\end{proposition}
\begin{proof}
Proposition~\ref{prop:width-four-shortening} leaves only
$A=2$.  By (R4), choose a cut $s_0$ for which
$r:=r_S(s_0)\le2^{2D-1}$.  For a cut at distance $h$ from
$s_0$, (R2) before taking logarithms and (R3) imply
\[
 r_U(t)\ge
 \max\{2^{4D}-2^{8h}r,\ 2^{4(D-h)}-r\}.          \tag{R5}
\]
The first quantity uses the ordinary target-rank inequality
on cut $t$ and the upper bound $r_S(t)\le2^{8h}r$.

For $D=3$, the distance layers have sizes $1,3$ and
$r\le32$.  Formula (R5) gives lower ceilings $12,8$,
respectively.  Thus
$\sum_s\kappa_U(s)\ge12+3\cdot8=36$, contradicting
(R1), whose upper bound is $35$.

For $D=4$, the layer sizes are $1,4,3$ and $r\le128$.
At distance one the first bound in (R5) is $2^{15}$;
at distance two the second is $128=2^7$.
The three lower ceilings are consequently $16,15,7$.
Their weighted sum is $97$, exceeding the upper bound
$3DM-1=95$.

For $D=5$, the layer sizes are $1,5,10$ and $r\le512$.
At distance one the first bound is
$2^{20}-2^{17}>2^{19}$; at distance two the second is
$2^{12}-512=3584>2^{11}$.
The lower ceilings are $20,20,12$, whose weighted sum is
$240$, exceeding the upper bound $239$.
\end{proof}

\begin{proposition}[SH-26a: a dimension bound for every fixed width and term count]
\label{prop:shortening-dimension-bound}
If a strict shortening has $A\ge2$ terms, put $a=\log_2 A$.
Then necessarily
\[
 a^2(D-1)^2<\frac{(A-1)m^2D}{2}.                 \tag{R6}
\]
In particular, a two-term width-four shortening must have
$D\le9$.
\end{proposition}
\begin{proof}
For each nonzero term use the log-rank ceiling $\kappa_j$
on the full sign cube, and for a zero term set
$\kappa_j\equiv0$.  A strict shortening necessarily has
$a<m$, so zero terms never meet the rank threshold used below.
The profiles are invariant under
simultaneous sign reversal and are $2m$-Lipschitz along
cube edges, by the proof of (R2).

Here is the needed variance estimate in full.  Expand a
real function $f$ on the uniform sign cube in its orthonormal
Walsh basis, $f=\sum_J\widehat f(J)\chi_J$.
Complement invariance makes every odd-degree coefficient
zero.  If $s^{(i)}$ flips coordinate $i$, Parseval gives
\[
 \sum_{i=1}^D\mathbb E(f(s)-f(s^{(i)}))^2
 =4\sum_J|J|\widehat f(J)^2
 \ge8\operatorname{Var}f.
\]
Applying the edge Lipschitz bound yields
\[
 \operatorname{Var}\kappa_j\le\frac{m^2D}{2}=:V.
 \tag{R7}
\]
The frame inequality (S1) gives
$\mathbb E\kappa_j\le c(N_j)/2<D(m-a)$.
At every cut, rank subadditivity for the sum of $A$ terms
forces at least one rank to be at least $2^{Dm}/A$.
Thus the events $\{\kappa_j\ge Dm-a\}$ cover the cube.

For completeness, the one-sided variance bound used here
follows by applying Markov's inequality to $(X+b)^2$ when
$\mathbb EX=0$: on $X\ge t$, it is at least $(t+b)^2$,
and minimizing $(\operatorname{Var}X+b^2)/(t+b)^2$ at
$b=\operatorname{Var}X/t$ gives
$\mathbb P(X\ge t)\le\operatorname{Var}X/
(\operatorname{Var}X+t^2)$.
The gap between $Dm-a$ and the mean of $\kappa_j$ is
strictly larger than $a(D-1)$.  Therefore each covering
event has probability strictly less than
\[
 \frac{V}{V+a^2(D-1)^2}.
\]
The union bound gives $1<AV/(V+a^2(D-1)^2)$, which
is (R6).  For $m=4,A=2$, it requires
$(D-1)^2<8D$, which fails for every $D\ge10$.
\end{proof}

\begin{proposition}[SH-26b: exact finite certificates in dimensions six through nine]
\label{prop:width-four-finite-certificates}
For each $D\in\{6,7,8,9\}$, no pair of integer profiles
satisfies (R1)--(R4) with $m=4$.  Consequently there is
no strict shortening with $m=4$ in these dimensions,
over any field.
\end{proposition}
\begin{proof}[Proof by rational certificates]
We specify the finite system and the exact certificates;
no floating-point infeasibility assertion is used.
Let $d=\lfloor D/2\rfloor$, $n=4D$, and let
$w_h$ be the number of folded sign vectors at distance $h$
from zero:
\[
 w_h=\binom Dh\ (2h<D),\qquad
 w_{D/2}=\tfrac12\binom D{D/2}\ (D\text{ even}).
\]
For $j\in\{S,U\}$, $0\le h\le d$, and $1\le t\le n$,
let $x_{j,h,t}$ be the average of the indicator
$[\kappa_j(s)\ge t]$ over that distance layer.
There are $2(d+1)n$ real variables.  They necessarily obey
\begin{align}
 0\le x_{j,h,t}&\le1,&
 x_{j,h,t+1}&\le x_{j,h,t},                         \tag{R8}\\
 \sum_{h=0}^d w_h\sum_{t=1}^n x_{j,h,t}
 &\le3D2^{D-1}-1,&x_{S,0,2D}&=0.                  \tag{R9}
\end{align}

Consider every triple $(h,k,\ell)$ realized as the distances
from zero of two folded sign vectors, and their mutual
distance.  Put $L=4(D-\ell)$.
By (R3), $2^{\kappa_S(s)}+2^{\kappa_U(t)}\ge2^L$.
For integer exponents this requires either one exponent
at least $L$, or both exponents at least $L-1$.
It follows that
\begin{align}
 x_{S,h,L-1}+x_{U,k,L}&\ge1,\\
 x_{S,h,L}+x_{U,k,L-1}&\ge1.                       \tag{R10}
\end{align}
To justify averaging even when a triple has several pair
orbits, average over any one orbit under coordinate
permutations fixing zero.  Both marginal distributions
are uniform on their respective layers.
For each realized triple with $\ell=1$, (R2) similarly
gives, for $t>8$,
\[
 x_{j,h,t}\le x_{j,k,t-8}.                         \tag{R11}
\]

Write every inequality (R8)--(R11) and the low-cut equation
as a row $A_i x\le b_i$, using
$x_{S,0,2D}\le0$ together with nonnegativity for the equation.
The exact row order and all coefficients are generated by
\path{programs/sh_radial_rank_certificate.py}.
For each of the four dimensions, the corresponding file
\path{outputs/sh_radial_rank_D6_m4.json},
\path{outputs/sh_radial_rank_D7_m4.json},
\path{outputs/sh_radial_rank_D8_m4.json}, or
\path{outputs/sh_radial_rank_D9_m4.json}
lists nonnegative rational coefficients $\lambda_i$ such that
\[
 \sum_i\lambda_i A_i=0,\qquad
 \sum_i\lambda_i b_i=-1.                           \tag{R12}
\]
Their sizes and common denominator bounds are as follows.
\[
\begin{array}{c|r|r|r|r}
 D&\text{variables}&\text{rows}&
 \text{nonzero coefficients}&\text{common denominator}\\\hline
 6&192&809&193&4410\\
 7&224&1007&225&11926\\
 8&320&1421&321&26846\\
 9&360&1687&361&127447
\end{array}
\]
The verifier reconstructs every row, checks every coefficient
is nonnegative, adds coefficients with exact rational
arithmetic, checks every variable coefficient is exactly
zero, and checks the resulting right side is exactly $-1$.
Thus (R12) proves the contradiction $0\le-1$.
The solver was used only to propose coefficients; the saved
certificates and exact verifier are sufficient to reproduce
this proof without optimization.
An independently written second verifier,
\path{programs/rt_verify_rank_profile_certificates.py},
reconstructs rows directly from their mathematical labels,
checks layer-pair existence by an intersection-size formula,
and verifies (R12) without importing the generator or using
an optimization solver.  Its four successful checks are
recorded in \path{outputs/rt_rank_profile_verification.json}.
\end{proof}

\begin{theorem}[SH-26c: all widths at most four are excluded]
\label{thm:no-shortening-width-four}
Over every field, every strict shortening of the coordinate
diameter tensor has cable width $m\ge5$.
\end{theorem}
\begin{proof}
Proposition~\ref{prop:no-small-width-shortening} handles
$m\le3$.  At $m=4$,
Proposition~\ref{prop:width-four-shortening} leaves only
two terms and $D\ge3$; Proposition~\ref{prop:width-four-three-five}
excludes $3\le D\le5$;
Proposition~\ref{prop:width-four-finite-certificates}
excludes $6\le D\le9$; and
Proposition~\ref{prop:shortening-dimension-bound}
excludes $D\ge10$.
\end{proof}

This is a finite-width exclusion, not an answer to the
unbounded-width existence question.  It gives no general
lower bound on the cost of arbitrary sums once $m\ge5$.

\begin{proposition}[SH-27: remaining parameters at width five]
\label{prop:width-five-parameters}
Over every field, a strict shortening at width $m=5$ must
have exactly two terms and satisfy $4\le D\le14$.
\end{proposition}
\begin{proof}
Write $a=\log_2 A$.  A single term is excluded by SH-01.
For $A\ge5$, (S2) fails: $\log_2 5>16/7$ follows from
$5^7>2^{16}$, while $35/16<(3+\sqrt2)/2$.
Thus only $A=2,3,4$ need consideration.

\emph{Four terms.}
When $A=4$, the necessary continuous bound (S2) requires
$D(6-5\sqrt2)+2\sqrt2>0$.  This already fails at
$D=3$, since $18-13\sqrt2<0$, and decreases thereafter.
At $D=2$ the cost bound is $w<12$.  Every term has
capacity at most $8$ on each adjacent cut by (S1).
The target has rank $1024$, so the sum of four term
ranks forces all four ranks on both cuts to equal $256$.
Each term then has cost at least $8+8=16$, a contradiction.

\emph{Three terms outside dimensions three through five.}
For $A=3$, use $a>19/12$.  If $D\ge6$, then
\[
 \frac w{\sqrt2}<\frac{41\sqrt2}{12}D
 <\frac{29}{6}D\le5D-1,
\]
where $41\sqrt2<58$ follows by squaring.
Every term rank on a sign cut is therefore at most
$2^{5D-2}$; three terms cannot reach $2^{5D}$.
For $D=2$, $w<41/3<14$.  The two adjacent integer
capacities of a term have sum at most $13$ and are
individually at most $9$.  Consequently the sum of its
two adjacent flattening ranks is at most
$2^9+2^4=528$.  Three terms give at most $1584$,
whereas the two target ranks sum to $2048$.

\emph{Three terms in dimensions three through five.}
Use the log-rank ceiling profiles $\kappa_j$ as above.
Set $K=5D-1$ and $M=2^{D-1}$.
The individual flow bound gives $\kappa_j(s)\le K$.
The frame bound and $a>19/12$ give the following
integer budgets $B$ for each term:
\[
\begin{array}{c|r|r|r|r}
 D&M&K&B&3B-2KM\\\hline
 3&4&14&40&8\\
 4&8&19&109&23\\
 5&16&24&273&51
\end{array}
\]
On every cut, at least two terms have ceiling $K$, or
exactly one has ceiling $K$ and both others have ceiling
$K-1$.  Indeed the three rank bounds must sum to
$2^{K+1}$, and all are at most $2^K$.
The former case costs at least $2K$ in summed ceilings;
the latter costs $3K-2$, an excess of $K-2$.
The table shows that the latter case can occur at most
$0,1,2$ times, respectively.  Hence the total numbers
of occurrences of the ceiling $K$ are at least
$8,15,30$.

For $D=3$, a term can have ceiling $14$ at at most two
cuts because its budget is $40$.  Thus the three terms
have at most six high occurrences, a contradiction.
For $D=4$, the folded cube graph is $K_{4,4}$.
Five high vertices meet both bipartition classes, so
every other vertex neighbors a high vertex.  The
$10$-Lipschitz bound then gives ceiling at least $9$
at each of the other three vertices.  The budget would
be at least $5\cdot19+3\cdot9=122>109$.
Each term therefore has at most four high occurrences,
giving at most twelve in total.

For $D=5$, the folded cube is a triangle-free
$5$-regular graph on $16$ vertices.  Among any ten
chosen high vertices, at least five of the remaining
six vertices neighbor a chosen high vertex.  A vertex
with no such neighbor must be adjacent to all the other
five remaining vertices.  Two such vertices, together
with any third remaining vertex, would form a triangle.
Thus ten high vertices force budget at least
$10\cdot24+5\cdot14=310>273$ by the same Lipschitz
bound.  Each term has at most nine high occurrences,
giving at most twenty-seven in total.  This finishes
the three-term exclusions.

\emph{Two terms in the plane.}
Here $q=32$ and strict shortening would require $w<16$.
The adjacent capacities of either term sum to at most
$15$.  Rank $q^2=1024$ on both target cuts forces the
high cuts to belong to different terms.  Initially
the high capacities are at least $9$ and the opposite
low capacities at most $6$.  The high ranks are then
at least $1024-64=960>512$, forcing high capacities
at least $10$ and low capacities at most $5$.
Thus the low ranks are at most $32$ and the high
ranks at least $992$.  Every singleton rank is at
least $31$, since $30\cdot32<992$.
The opposite-terminal slice used in SH-24 shows that
the two low ranks sum to at least $32$, so one is
at least $16$.  For that term, (S3) gives
\[
 c(N)\ge4(\sqrt2-1)\log_2 31
 +(2-\sqrt2)(4+\log_2 992)
 >\frac{42+27\sqrt2}{5}>16.
\]
For the first strict comparison use
$\log_2 31>24/5$, since $31^5>2^{24}$,
and $992=32\cdot31$; the second follows from
$\sqrt2>38/27$.  This is a contradiction.

\emph{Two terms in dimension three.}
Now $n=15$, $M=4$, and (R1) gives budget $47$ for
each term.  The low-cut bound gives
$\kappa_S(0)\le8$.  Mixed slicing gives
$\kappa_U\ge10$ on the other three cuts, and the
ordinary target rank gives $\kappa_U(0)=15$.
The budget $47$ consequently makes each of the
other three ceilings at most $12$.  On those cuts,
rank subadditivity therefore forces $\kappa_S=15$.
The budget for $S$ implies $\kappa_S(0)\le2$.
But every other folded-three-cube vertex is an axis
flip from zero, and the $10$-Lipschitz bound requires
$\kappa_S(0)\ge5$.  This is the last contradiction.

Finally, for $m=5,A=2$, (R6) requires
$(D-1)^2<25D/2$, which fails for all $D\ge15$.
The only remaining parameters are as stated.
\end{proof}

\begin{proposition}[SH-28a: width five in dimension four]
\label{prop:width-five-dimension-four}
There is no strict shortening with $m=5,D=4$ over any field.
\end{proposition}
\begin{proof}
By SH-27 only two terms remain.  Their sign-cut ceiling
profiles on the folded cube $K_{4,4}$ each have total at
most $127$, are $10$-Lipschitz, and lie between $0$ and $20$.
Choose the absolute minimum $x$ of both profiles.  By (R4),
$x\le11$; relabel and translate so $\kappa_S(0)=x$.
Call the parity class containing zero even.  Mixed slicing
gives $\kappa_U(0)=20$ and $\kappa_U\ge15$ at all four
odd vertices.

If $x\le8$, Lipschitz continuity gives
$\kappa_S\le18$ on the odd class, so $\kappa_U=20$
there.  On the other three even vertices, mixed slicing
at distance two gives $\kappa_U\ge10$.
Thus the sum of the $U$ profile is at least
$20+4\cdot20+3\cdot10=130>127$.
If $x=9$, the odd $U$ values are at least $19$ and
the other even $U$ values at least $9$.
The latter three values sum to at most $31$, so each
is at most $13$.  Hence $S$ is $20$ at these three
vertices.  Mixed slicing from any of them forces
$S\ge15$ throughout the odd class.  Its profile sum
is at least $9+3\cdot20+4\cdot15=129>127$.

It remains to consider $x=10$ or $11$.
Each term has value $20$ at at most four vertices,
by its budget and the global lower bound $x$.
At every vertex one of the two values is $20$, or
both are $19$.  Relative to the baseline $20+x$
per vertex, the latter exception costs $18-x$
extra units; a vertex with both values $20$ costs
$20-x$ extra units.  The total available excess is
\[
 R=254-8(20+x)=
 \begin{cases}14,&x=10,\\6,&x=11.\end{cases}
\]
For $x=11$ there are no exceptions.  For $x=10$
there is at most one exception, and an exception
cannot coexist with a double-$20$ vertex because
their excess $8+10$ exceeds $R$.

With no exceptions, the sets where $S$ and $U$
equal $20$ partition the eight vertices into two
sets of size four.  The edges joining these sets
have a matching of size four.  Indeed, if the
first set has $a$ even and $4-a$ odd vertices,
the cross-edge graph is the disjoint union of
$K_{a,a}$ and $K_{4-a,4-a}$.
With one exception, the two sets have sizes four
and three and the same argument gives a matching
of size three after deleting the exceptional vertex.

On every matched edge the two opposite, potentially
low, profile values obey a mixed-slice rank bound
with target $2^{15}$.  At least one of these two
ceilings is at least $14$.  Since the matching has
disjoint endpoints, it consumes at least $14-x$
extra units per edge.  With no exception this gives
$4(14-x)>R$ for both $x=10,11$.  With one exception,
$x=10$ and the required excess is at least
$8+3(14-10)=20>14$.  Every case is impossible.
\end{proof}

\begin{proposition}[SH-28b: complete minimum branches in dimensions five, seven, and nine]
\label{prop:width-five-odd-certificates}
There is no strict shortening with $m=5$ and
$D\in\{5,7,9\}$ over any field.
\end{proposition}
\begin{proof}[Proof by exact rational certificates]
Choose the absolute minimum $t$ of both term profiles,
relabel its term $S$, and translate its cut to zero.
By (R4), $t$ is one of the integers $0,\ldots,3D-1$.
Use the radial threshold system (R8)--(R11), replacing
$n=4D$ by $n=5D$, the budget by $4D2^{D-1}-1$,
the mixed-cut exponent by $L=5(D-\ell)$, the low-cut
condition by $x_{S,0,3D}=0$, and the Lipschitz shift
$8$ by $10$.
In the branch with minimum $t$, impose
\[
 x_{S,0,t+1}=0,\qquad x_{S,h,t}=1\quad(t>0)
 \quad\text{for every layer }h.
 \tag{R13}
\]
Some branches also use $x_{U,h,t}=1$, which follows
from choosing the absolute minimum of both profiles.
All these are linear inequalities using the existing
box bounds; for $t=0$ the lower conditions are vacuous.

The file
\path{outputs/sh_width5_single_minimum_coverage.json}
lists one exact certificate for every branch: fifteen
for $D=5$, twenty-one for $D=7$, and twenty-seven for
$D=9$.  Each certificate is a nonnegative rational
combination of the stated inequalities that cancels
every variable and has right side exactly $-1$.
The independent verifier
\path{programs/rt_verify_rank_profile_certificates.py}
reconstructs these additional minimum rows from their
labels and validates their branch metadata, as well
as the original mixed-slice and distance constraints.
Its checks are recorded in
\path{outputs/rt_rank_profile_width5_verification.json}.

For a few certificates, the numerical proposal did not
admit direct rational reconstruction at the chosen
denominator bound.  The generator then rounded every
nonnegative proposed coefficient to a common rational
grid, computed the exact residual coefficient of each
variable, and cancelled that residual using the
available rows $-x\le0$ or $x\le1$ with nonnegative
weights.  It checked that the corrected right side
was strictly negative, and rescaled it to $-1$.
This is an exact use of the box inequalities; no
approximate cancellation is accepted by either verifier.
Thus every integer minimum branch is excluded.
\end{proof}

\begin{proposition}[SH-28c: all remaining width-five minimum branches]
\label{prop:width-five-two-minima}
There is no strict shortening with $m=5$ and
$D\in\{6,8,10,11,12,13,14\}$ over any field.
\end{proposition}
\begin{proof}[Proof by exact rational certificates]
Continue the absolute-minimum branching of SH-28b.
Every branch is excluded by a radial certificate except
for the first-minimum values in the following table.
These values are handled by additionally fixing the
minimum of the other term and the distance between
the two marked minimum cuts.
\[
\begin{array}{c|l|r|r}
 D&\text{first minima }t&\text{admissible cases}&
 \text{cases excluded by the rank bound}\\\hline
 6&15&13&23\\
 8&20&21&39\\
 10&25&30&60\\
 11&25,30&54&144\\
 12&30,31&83&162\\
 13&30,35&76&197\\
 14&35,36,37&162&318
\end{array}
\]
For each listed $t$, let $y$ be the minimum of the
$U$ profile and let $h$ be the folded Hamming distance
between its chosen minimum and the chosen minimum
of $S$.  Necessarily
\[
 t\le y\le4D-1,\qquad0\le h\le\lfloor D/2\rfloor.
 \tag{R14}
\]
The upper bound on $y$ follows from the strict mean
bound $\mathbb E\kappa_U<4D$.
Equation (R3) immediately excludes every case with
$2^t+2^y<2^{5(D-h)}$; these are exactly the cases
in the last column of the table.

For each remaining case, translate the first minimum
to zero and permute coordinates so the second minimum
is the cut with its first $h$ signs reversed.
Average the threshold variables over the group
$\mathfrak S_h\times\mathfrak S_{D-h}$, which preserves
both marked cuts.  Its folded-cube orbits are indexed
by pairs
\[
 (a,b),\quad0\le a\le h,\quad0\le b\le D-h,
 \qquad(a,b)\sim(h-a,D-h-b).
\]
Take the lexicographically smaller pair as representative.
Its orbit has size
$\binom ha\binom{D-h}b$, except that this number is
halved when the pair equals its complement.
The weights sum to $M=2^{D-1}$.

For two representative pairs $(a,b)$ and $(c,d)$,
the allowed mutual distances are
\[
 \min\{a+c-2i+b+d-2j,
 D-a-c+2i-b-d+2j\},                               \tag{R15}
\]
where
\[
 \max(0,a+c-h)\le i\le\min(a,c),\quad
 \max(0,b+d-(D-h))\le j\le\min(b,d).
\]
Thus every needed orbit size and distance relation has
an explicit finite combinatorial description.

Use one variable $x_{V,o,z}$ for the threshold average
at level $z$ on orbit $o$.  Impose monotonicity and the
box constraints, the weighted budget $4DM-1$ for each
term, and the same mixed-cut and Lipschitz inequalities
as in SH-28b for every allowed orbit-pair distance.
At the two marked orbits impose upper thresholds
$\kappa_S\le t$ and $\kappa_U\le y$, respectively;
at every orbit impose lower thresholds
$\kappa_S\ge t$ and $\kappa_U\ge y$.
These constraints follow from the definitions of the
two minima.  Averaging over each orbit of pairs is
valid because its marginals are uniform on the
corresponding vertex orbits.

The program
\path{programs/sh_two_minimum_certificate.py}
generates this finite system.  All $439$ admissible
cases in the table have saved nonnegative rational
coefficients whose sum of constraint rows is exactly
$0\le-1$.  The batch manifests
\path{outputs/sh_two_min_D6_m5_manifest.json} through
\path{outputs/sh_two_min_D14_m5_manifest.json}, for
the dimensions appearing in the table, list every
admissible case and every exact rank-bound exclusion.

An independent verifier,
\path{programs/rt_verify_two_minimum_certificates.py},
does not import the generator or use its intersection
formulas.  It enumerates the full folded cube to
recover orbit sizes, marked indices, and allowed
mutual distances, reconstructs each inequality from
its mathematical label, and checks every coefficient
and cancellation with exact rational arithmetic.
It also verifies complete coverage of (R14) and the
integer inequalities used to prune cases.

The authoritative combined input is
\path{inputs/rt_width5_full_coverage.json}.
The independent final checker
\path{programs/rt_verify_width5_coverage.py}
rechecks all $273$ selected radial certificates and
all $439$ two-minimum certificates, verifies their
SHA-256 identities, and confirms that all $285$
possible first-minimum branches for $5\le D\le14$
are covered.  Its successful output is
\path{outputs/rt_width5_full_coverage.json}.
This establishes the stated exclusions using exact
finite certificates, independently of the numerical
optimization that proposed them.
\end{proof}

\begin{theorem}[SH-28d: every strict shortening has width at least six]
\label{thm:no-shortening-width-five}
Over every field, there is no strict shortening of
$C_{D,m}$ for any $D\ge2$ and any $m\le5$.
Equivalently, every strict shortening in the exact
definition of the question must have $m\ge6$.
\end{theorem}
\begin{proof}
SH-26c excludes $m\le4$.
At $m=5$, SH-27 leaves only two terms and
$4\le D\le14$.  SH-28a excludes $D=4$,
SH-28b excludes $D=5,7,9$, and SH-28c excludes
all other dimensions in this interval.
\end{proof}

The theorem is uniform in the field and allows arbitrary
finite network layouts and internal local tensors.
It is a lower bound on the first possible cable width;
it does not settle the existence question at unbounded width.

\subsection{A seven-party sufficient gate for contraction}

\begin{proposition}[SH-29: a seven-point hypermetric obstruction]
\label{prop:seven-point-hypermetric}
Delete any one negative-parity point from the eight matrices
in Proposition~\ref{prop:eight-point-cnd}.  On the remaining
seven points take the four coefficients $+1$ and three
coefficients $-1$ inherited from parity.  For the metric
$d(a,b)=\sqrt{\operatorname{rank}(a-b)}$, their quadratic
form is
\[
 \sum_{i,j=1}^7c_i c_j d(i,j)
 =6(2\sqrt3-\sqrt2-2)>\frac6{25}>0,
 \qquad\sum_{i=1}^7c_i=1.                         \tag{H1}
\]
Suppose that, over the chosen tensor field, every seven-party
price satisfies the corresponding hypermetric inequality
$\sum c_i c_j d(i,j)\le0$ for four positive and three
negative coefficients.  Under the CAT(0) price embedding,
every price on $\operatorname{Mat}_3(\mathbb F_2)$ then obeys
\[
 D(3)^2\le(3-2^{-523336})D(1)^2.
 \tag{H2}
\]
In particular, the following single asymptotic graph-cover
assertion is sufficient for the contracting hypothesis:
\[
 \bigl(K_4(2)\otimes K_3(2)\bigr)^{\otimes N}
 \text{ is covered by }2^{o(N)}\text{ restrictions of }
 K_{4,3}(2)^{\otimes N}.                          \tag{H3}
\]
The two cliques in (H3) occupy the two respective parts of
the complete bipartite graph.  No assertion that (H3) holds
is made here.
\end{proposition}
\begin{proof}
Let $K$ be the eight-by-eight distance matrix from SH-20,
and let $c$ be its parity character.  That proposition gives
$Kc=\lambda_+c$, $c^{\mathsf T}Kc=8\lambda_+$, and
$\lambda_+=2\sqrt3-\sqrt2-2>1/25$.
Deleting a coordinate $v$ with $c_v=-1$ changes the
quadratic form by $-2c_v(Kc)_v$, because $K_{vv}=0$.
The remaining value is therefore $6\lambda_+$ and the
coefficient sum becomes one.  This proves (H1).
It does not contradict the strict conditional negative
type of proper subsets in SH-20: those inequalities
require coefficient sum zero.

Repeat the stability proof of SH-22b up to its estimate
(L4).  If $\delta=3-D(3)^2<2^{-523336}$ after normalizing
$D(1)^2=1$, every relevant distance differs from the
rank square-root metric by less than $1/256$.
There are only $42$ ordered off-diagonal entries on the
seven points.  Their hypermetric quadratic form therefore
remains positive, since $6/25>42/256$.
This contradicts the assumed seven-party inequality and
proves (H2), with exactly the same contraction factor.

Finally (H3) and spectral subadditivity imply
$\varphi(K_4(2)\otimes K_3(2))\le\varphi(K_{4,3}(2))$
after taking $N$th roots.  Logarithms and the wire
evaluation formula give the required hypermetric
inequality for every spectral wire metric, hence for
every price by scaling and limits.  The argument works
over either field in the question, or any other field.
\end{proof}

\paragraph{The exact and border obstructions do not settle this gate.}
For every $r\ge3$ and $q>1$, there is no exact restriction
\[
 K_r(q)\otimes K_{r-1}(q)\le K_{r,r-1}(q)
 \tag{H4}
\]
over any field.  Indeed each left-party source and target
space has dimension $q^{r-1}$.  Full target singleton
ranks force all left local maps to be invertible.
Each right local map has full row rank $q^{r-2}$.
Across the entire left part versus the entire right
part, the source is a full identity pairing.
The restricted tensor therefore has matrix rank
$q^{(r-1)(r-2)}>1$, whereas the target in (H4) has rank one.
This is a direct contradiction and applies to every equal
tensor power by replacing $q$ by $q^N$.

Over $\mathbb C$, Corollary~\ref{pg:cap-hyper}
strengthens this to a border-restriction obstruction at
every equal power, by a multiplicative capacity argument.
Both statements concern one restriction with unchanged
source exponent.  Assertion (H3) permits a subexponential
sum of restrictions, and neither obstruction proves its
failure or success.  Thus the smaller seven-party gate
improves the finite sufficient test for contraction while
leaving the original alternative undecided.

\subsection{The current rank relaxation permits width six}

\begin{proposition}[SH-30: an exact feasible abstract rank profile]
\label{prop:width-six-rank-relaxation}
The sign-cut rank inequalities, mixed-cut slicing bounds,
rank-ceiling Lipschitz bounds, and planar Crofton bounds
used above do not exclude $m=6,D=2,A=2$ by themselves.
They have the following exact feasible numerical profile:
\[
\begin{array}{c|c|c|c|c}
 &\text{four singleton ranks}&r_{ab|cd}&r_{ad|bc}&r_{ac|bd}\\\hline
 S&(64,64,64,64)&32&4064&128\\
 U&(64,64,64,64)&4064&32&128
\end{array}
\]
This is a statement about necessary inequalities only.
It asserts neither a pair of tensors with these ranks,
nor a tensor identity, nor a geometric network realization.
\end{proposition}
\begin{proof}
The target has rank $4096$ on both adjacent cuts, and
$32+4064=4096$.  The mixed low-cut slice has target
rank $64$, and $32+32=64$.
The two sign-cut ceiling profiles are $(5,12)$ and
$(12,5)$.  They satisfy the integer frame budgets
$5+12\le19$ and the rank-change bound $|12-5|\le12$.
All ranks respect their terminal dimension bounds.
The remaining balanced rank $128$ also satisfies
every pairwise product inequality between the three
balanced flattening ranks, for either term.

Because cable capacities are integers, the proof of
(S3) even permits replacing every log-rank by its
ceiling before integrating.  For the displayed profile
this stronger Crofton lower bound is
\[
 24(\sqrt2-1)+17(2-\sqrt2)=10+7\sqrt2.
\]
It is strictly less than the rational cost allowance
$w=199/10$, because $99^2>2\cdot70^2$.
The individual geometric cut bounds also fit below
$w$, while
\[
 w+4\log_2 A=\frac{239}{10}<24=2Dm.
\]
Thus all the listed necessary bounds are simultaneously
compatible with the numerical allowance.  Every check
is recorded with exact integer or symbolic arithmetic in
\path{programs/sh_width6_rank_relaxation.py}, using
\path{inputs/sh_width6_rank_relaxation.json}; the result
is \path{outputs/sh_width6_rank_relaxation.json}.
\end{proof}

Further exclusions at this width need information about
tensor or network realizability beyond these rank profiles,
or a stronger geometric inequality.  The feasibility of the
relaxation supplies no positive shortening construction.

\paragraph{The continuous-capacity geometric bound is already sharp.}
There is also an exact real-capacity graph attaining this
Crofton lower bound for the $U$ cut profile in the table.
Put capacity $7/2$ on the terminal
segments $ad$ and $bc$, capacity $5/2$ on the diameters
$ac$ and $bd$, and no other edges.  Its four singleton
cuts have capacity $6$, its adjacent cuts $ad\mid bc$
and $ab\mid cd$ have capacities $5$ and $12$, and its
remaining balanced cut has capacity $7$.  Its cost is
\[
 2\cdot\frac72\sqrt2+2\cdot\frac52\cdot2
 =7\sqrt2+10.
\]
These values are checked exactly in the same saved program.
The half-integer capacities are forbidden by the bit-cable
definition.  Thus this graph is a geometric relaxation,
not an admissible shortening network.  Nevertheless it
shows that a homogeneous geometric dual valid for all
real capacities cannot improve the bound for this cut
profile.  Any remaining obstruction must use integer
capacities, tensor realizability, or further information
absent from this relaxation.

\subsection{An exact binary pentagon cover and its compression limit}

\begin{proposition}[SH-31: an exact two-term cover on five parties]
\label{prop:binary-pentagon-cover}
Write $G_q=K_{3,2}(q)$, with left parties $L_0,L_1,L_2$
and right parties $R_0,R_1$, and write
$H_q=K_3(q)\otimes K_2(q)$ for the two cliques on these
parts.  Over every field, $H_2$ is a sum of two exact
restrictions of $G_2$.  Two is the minimum possible number
of exact restrictions.  More generally $H_q$ is covered
by $1+\binom q2$ restrictions of $G_q$; for $q=2^N$
there is a cover with $q$ terms.
\end{proposition}
\begin{proof}
For $0\le a<b<q$, let $J^{ab}$ have entries $+1$ at
$(a,b)$ and $-1$ at $(b,a)$, with all other entries zero.
Direct expansion gives the integer coefficient identity
\[
 \delta_{xy}\delta_{uv}
 =\delta_{xv}\delta_{yu}
   +\sum_{a<b}J^{ab}_{xu}J^{ab}_{yv}.              \tag{P1}
\]
Indeed the sum on the right equals
$\delta_{xy}\delta_{uv}-\delta_{xv}\delta_{yu}$:
when $x=u$ or $y=v$ its skew-matrix contribution vanishes,
and otherwise the unique possible unordered pair gives
the indicated sign.  This proves (P1) over the integers,
including all coincident-index cases, and hence over
every field.

Here are all the local maps implementing the identity.
Let the source edge variables at $L_i$ be $(a_i,b_i)$,
where $a_i$ is incident to $R_0$ and $b_i$ to $R_1$.
Use target left coordinates
\[
 L_0=(X,Y_0),\qquad L_1=(X',Z),\qquad
 L_2=(Y_2,Z').
\]
For the first term use the identity at each left party
and the two right maps
\[
 M_0[u,(a_0,a_1,a_2)]
       =\delta_{a_0a_1}\delta_{ua_2},\qquad
 M_1[v,(b_0,b_1,b_2)]
       =\delta_{b_1b_2}\delta_{vb_0}.              \tag{P2}
\]
Its output coefficient is
$\delta_{XX'}\delta_{ZZ'}\delta_{Y_0v}\delta_{Y_2u}$.
For each $J=J^{ab}$ use the swap
$(a_i,b_i)\mapsto(b_i,a_i)$ at every left party and
\[
 M_0[u,(a_0,a_1,a_2)]
       =\delta_{a_1a_2}J_{a_0u},\qquad
 M_1[v,(b_0,b_1,b_2)]
       =\delta_{b_0b_1}J_{b_2v}.                \tag{P3}
\]
This coefficient is
$\delta_{XX'}\delta_{ZZ'}J_{Y_0u}J_{Y_2v}$.
For fixed left output indices the source edge assignment
is unique, so there is no additional contraction
multiplicity.  Summing (P2)--(P3) and applying (P1) gives
\[
 \delta_{XX'}\delta_{ZZ'}\delta_{Y_0Y_2}\delta_{uv},
\]
which is exactly $H_q$.

At $q=2$ there is only one skew matrix, giving two
terms.  The exact one-term obstruction (H4), with
$r=3$, proves minimality.  Tensoring the binary identity
$N$ times and grouping equal parties gives $2^N=q$
restrictions of $G_2^{\otimes N}=G_q$ whose sum is
$H_2^{\otimes N}=H_q$.
\end{proof}

\paragraph{Coefficient certificates and costs.}
The full maps (P2)--(P3), source tensor, each term and
their sum are preserved in
\path{outputs/sh_pentagon_exact_cover_q2_tables.npz}
and \path{outputs/sh_pentagon_exact_cover_q3_tables.npz}.
The program \path{programs/sh_pentagon_exact_cover.py},
with \path{inputs/sh_pentagon_exact_cover.json},
checks every output coefficient: $256$ entries at $q=2$
and $6561$ at $q=3$, with maximum error zero.
It also checks all $q^4$ entries of (P1) for
$q=2,3,4,5,8$.  The recorded result is
\path{outputs/sh_pentagon_exact_cover.json}.
All local coefficients are in $\{-1,0,1\}$.

For $q=2^N$, every source layout has six width-$N$
edges and exact cost
\[
 N\sum_{i=0}^2\sum_{j=0}^1|p_{L_i}-p_{R_j}|.
\]
The target has cost
$N(\sum_{i<j}|p_{L_i}-p_{L_j}|+|p_{R_0}-p_{R_1}|)$.
The cover just proved has $A=2^N$, so $\log_2A=N$;
its overhead is not $o(N)$.  A selector distributed
from a centre $o$ by five width-$N$ cables costs
$N\sum_v|p_v-o|$.  In normalized spectral notation
the binary cover yields only
$\log_2\psi(H_2)-\log_2\psi(G_2)\le1$.
No subexponential cover, price inequality without this
overhead, or strict shortening follows from this
tensor-power expansion alone.

\begin{proposition}[SH-32: the first mixed profile cannot be folded
into one unchanged source]
\label{prop:pentagon-profile-obstruction}
Let $H_2=S+U$ be the two terms of SH-31 in their stated
coordinates, and group equal parties in
\[
 W=S\otimes U+U\otimes S.
\]
Over every field, $W$ is not an exact restriction of
$G_4=G_2^{\otimes2}$.
This concerns the unchanged source graph and dimensions;
it does not concern profiles augmented by selector wires.
\end{proposition}
\begin{proof}
Across the three left parties versus the two right
parties, $S$ is an injective matrix with four columns.
In right coordinates $(u,v)$, write
$v_0=(1,0,0,1)^{\mathsf T}$ and
\[
 R=v_0v_0^{\mathsf T}-I_4.
\]
The coefficient formulas (P2)--(P3) give $U=SR$.
Consequently the global left--right matrix rank of $W$
is the rank of $R\otimes I+I\otimes R$.
In characteristic other than two, $R$ has eigenvalue
$+1$ with multiplicity one and $-1$ with multiplicity
three.  The tensor sum therefore has rank $1+9=10$.
In characteristic two, $R=I+N$ with
$N=v_0v_0^{\mathsf T}$, $\operatorname{rank}N=1$
and $N^2=0$.  Write its four-dimensional space as one
two-dimensional nilpotent Jordan block and two trivial
blocks.  On the corresponding tensor products,
$N\otimes I+I\otimes N$ has ranks $2$, $2$, $2$ and
$0$, respectively, and hence total rank six.

All five singleton ranks of $W$ are
$(16,16,16,4,4)$ over every field.  For completeness,
the two shared left path wires factor from a four-party
core with dimensions $(4,4,4,4)$, whose four singleton
flattenings have signed-permutation minors with
determinants $(1,-1,1,-1)$.  In the order
$(Y_0,Y_2,u,v)$, with binary words in lexicographic
order, their zero-based column sets are
\[
 (1,2,5,18),\quad(4,5,8,24),\quad
 (1,2,6,17),\quad(1,2,5,18).
\]
Thus the first and third left singleton ranks are
$4\cdot4=16$, the middle one is $4^2=16$, and both
right ranks are four.  The saved exact certificate
includes the four minors themselves and verifies the
entire flattenings by rational and binary elimination.

Suppose $W$ were a restriction of $G_4$.  The source
left spaces also have dimension $16$, so their local
maps would be invertible.  The two right local maps
would each have full row rank four.  The source
flattening across all left versus all right parties is
a full identity pairing.  After these maps it must have
rank $4\cdot4=16$, contradicting the rank ten or six
computed above.
\end{proof}

The program \path{programs/sh_pentagon_middle_profile.py}
uses \path{inputs/sh_pentagon_middle_profile.json}
and produces \path{outputs/sh_pentagon_middle_profile.json}
and the complete $65536$-entry tensor in the matching
\path{.npz} file.  It checks $448$ nonzero integer
coefficients, all singleton ranks, the two global ranks,
and the four unit minors.  The failure here is precisely
the stronger claim that counting multiplicity profiles
alone turns each profile into one restriction of the
original graph.  SH-02 adds selector resources to the
profile network and makes no such claim.  Neither this
rank obstruction nor the one-term border obstruction
excludes subexponential sums with more general local maps
and extra resources.

\begin{proposition}[SH-33: positive-overhead and growing-party
contraction gates]
\label{prop:positive-overhead-gates}
The subexponential hypothesis in (H3) can be weakened as follows.
Either of the next two assertions, over the chosen tensor field,
implies the same explicit contraction bound (H2):
\begin{enumerate}
\item For $q=2^N$ along an unbounded sequence, $H_q=K_4(q)
\otimes K_3(q)$ is a sum of at most
$q^{1/128+o(1)}$ restrictions of $K_{4,3}(q)$.
\item For some fixed $r=4t$, $K_r(q)\otimes K_r(q)$ is
a sum of at most $q^{\beta_r+o(1)}$ restrictions of
$K_{r,r}(q)$, along an unbounded sequence $q=2^N$, where
$\beta_r\le r^2/2048$.
\end{enumerate}
In particular, a family with $\beta_r=o(r^2)$ as
$r\to\infty$ through multiples of four is sufficient.
There is also a wholly finite version: it suffices to
find one identity at $q=2^N$ whose term count satisfies
$\log_2 A/N\le1/128$ in the first case, or
$\log_2 A/N\le r^2/2048$ in the second case.
No such sufficiently inexpensive identity is asserted here.
\end{proposition}
\begin{proof}
Normalize a nonzero price on
$\operatorname{Mat}_3(\mathbb F_2)$ by $D(1)^2=1$.
Assume toward contradiction that
$3-D(3)^2<2^{-523336}$.  SH-22b then puts every distance
on the eight-point configuration of SH-20 within
$\varepsilon=1/256$ of $\sqrt{\operatorname{rank}}$.
The elementary rational comparisons
$\sqrt3>433/250$ and $\sqrt2<283/200$ give
\[
 \lambda_+=2\sqrt3-\sqrt2-2>49/1000.              \tag{P4}
\]

For any finite nonzero price $p=s d_\psi$ in the
normalized form SH-11 and any spanning tree $\mathcal T$,
\[
 s\le\sum_{e\in\mathcal T}p(e).                 \tag{P5}
\]
Indeed, the product of its binary tree wires restricts
to the unit tensor of size two on all parties: copy
the common wire index at every vertex.  Evaluation by
the normalized spectrum gives
$1\le\sum_{e\in\mathcal T}d_\psi(e)$, proving (P5).

On the seven-point configuration of SH-29 there is a
spanning tree with five rank-one edges and one rank-two
edge.  The rank-one edges form two four-cycles before
deletion; after removing one vertex, use a spanning
path of the surviving three vertices, a spanning tree
of the untouched four-cycle, and one rank-two edge
between them.  Thus
\[
 s\le5+\sqrt2+6/256<13/2.                       \tag{P6}
\]
The unordered difference between within-part and
cross-part costs satisfies
\[
 \Delta_7=p(K_4)+p(K_3)-p(K_{4,3})
 >3(49/1000)-21/256=2079/32000.                 \tag{P7}
\]
Here each of the $21$ unordered nonzero coefficients
incurs error at most $1/256$.
If the first cover assertion holds, evaluation by
$\psi$, division by $N$, and passage to its unbounded
sequence give
$\Delta_7\le s/128<13/256$.
This contradicts (P7), since
$2079/32000-13/256=227/16000>0$.

For the second assertion use all eight points.
Their spanning tree has six rank-one edges and one
rank-two edge, so
\[
 s\le6+\sqrt2+7/256<15/2.
\]
Replicate every positive-parity point $t$ times and
every negative-parity point $t$ times.  This is a valid
pullback price on $8t$ parties.  Explicitly, merging
each fibre into its original party pulls a normalized
spectral point back to a normalized spectral point:
the full unit tensor remains locally isomorphic to
the same unit tensor after merging.  Wires within a
fibre become one-party nonzero vectors and have
spectral value one.  This verifies both the zero
distances between copies and preservation of the
normalization; it is the pullback operation of
Definition A.4 and the following paragraph in the
unrefereed Anthropic manuscript, p.~20.

The two parts now have $r=4t$ parties each, and the
unordered graph-cost difference is $t^2\Delta_8$, where
\[
 \Delta_8>4(49/1000)-28/256=693/8000.
\]
Its scale $s$ is unchanged; equivalently one may
extend the old spanning tree by zero-distance edges
joining the duplicate labels.  The proposed cover
would imply
\[
 t^2\Delta_8\le s\beta_r
 \le\frac{15}{2}\frac{t^2}{128}
 =\frac{15}{256}t^2,
\]
contrary to $693/8000-15/256=897/32000>0$.
Thus either cover assertion rules out the assumed
near-equality and proves (H2).

For a single finite identity, spectral evaluation
already gives the same inequalities with
$\beta=\log_2A/N$, so no limiting assertion is
needed.  Alternatively, its tensor powers give the
required unbounded sequence.  Finally
$\beta_r=o(r^2)$ eventually satisfies the fixed
threshold $r^2/2048$, proving the family statement.
\end{proof}

The second gate permits a growing number of parties
while keeping the binary edge dimension fixed: covers
of $K_r(2)\otimes K_r(2)$ by restrictions of
$K_{r,r}(2)$ with $\log_2 A_r=o(r^2)$ would suffice.
For example, an exact construction with
$A_r=2^{O(r\log r)}$ would settle this contraction
gate.  This is a weaker sufficient hypothesis than
requiring zero overhead at a fixed number of parties.
It supplies a concrete amplification target for
determinantal or matching identities, rather than
an inference that the two-term five-party identity
already has the required amplification.

The rational margins, root comparisons and spanning-tree
upper bounds in SH-33 are independently reproducible with
\path{programs/sh_positive_overhead_constants.py},
\path{inputs/sh_positive_overhead_constants.json}, and
\path{outputs/sh_positive_overhead_constants.json}.

\subsection{Exploratory cover searches: no further identity certified}

The finite identity SH-31 was first suggested by a complex
local optimization at $q=2$, using
\path{programs/sh_pentagon_two_cover_search.py} and
\path{inputs/sh_pentagon_two_cover_search_q2.json}.
Its best normalized squared overlap was
$0.9999999999999939$.  The retained numerical maps motivated
the exact reconstruction, but the proof and coefficient
certificate of SH-31 do not depend on this overlap.

Two subsequent searches used the six-edge factored contraction
in \path{programs/sh_pentagon_cover_search_factor.py}.
At $q=3$, three two-term starts, each with $1000$ iterations,
reached best squared overlap $0.4182498769$; four three-term
starts, each with $3000$ iterations, reached
$0.6274025508$.  The inputs and outputs have basenames
\path{sh_pentagon_cover_q3_A2} and
\path{sh_pentagon_cover_q3_A3} in their respective directories.
Analytic gradients were checked against centred finite
differences, with maximum errors below $5\cdot10^{-12}$.

For the actual seven-party graph gate, the factored evaluator
\path{programs/sh_hypermetric_cover_search.py} searched
two-term covers of $K_4(2)\otimes K_3(2)$ by restrictions
of $K_{4,3}(2)$.  Two complex starts with $750$ iterations
each reached best squared overlap $0.1171874099$.
The input \path{inputs/sh_hypermetric_cover_r4_q2_A2.json},
complete iteration summaries and best local tables under
\path{outputs/sh_hypermetric_cover_r4_q2_A2} are preserved.
Its gradient check had maximum error below
$3.1\cdot10^{-12}$.

All three later searches terminated at their iteration limits.
They provide neither exact nor border identities, and local
optimization gives no global impossibility certificate.
Their purpose was to test concrete small identities before
attempting symbolic reconstruction or amplification; only
the binary five-party identity passed an exact reconstruction.

\subsection{The integral geometric gap at the first remaining width}

The real-capacity witness of SH-30 has cost $10+7\sqrt2<20$.
The following integer network realizes the same four singleton
and two adjacent cut values, but at a larger cost.  It is a
geometric construction only; no tensor identity is claimed.

\begin{proposition}[SH-34a: an exact integer network]
\label{prop:integer-critical-network}
Put $a=(1,0)$, $b=(0,1)$, $c=(-1,0)$ and $d=(0,-1)$, and let
\[
 h=\frac{1-1/\sqrt3}{2},\qquad x=(h,-h),\qquad y=(-h,h).
\]
Take edges $ad,bc$ of capacity $3$, edges $ac,bd$ of capacity
$2$, and the five edges $ax,dx,by,cy,xy$ of capacity $1$.
The four terminal singleton mincuts are $6$, and
\[
 \lambda(ab)=12,\qquad \lambda(ad)=5,\qquad \lambda(ac)=8.
\]
Its Euclidean cost, minimized over the two internal positions
for this fixed graph, is
\[
 8+7\sqrt2+\sqrt6
   =10+7\sqrt2+(\sqrt6-2)>20.
                                                        \tag{SH-34.1}
\]
Both internal vertices have odd weighted degree $3$.
\end{proposition}

\begin{proof}
The four direct edges contribute singleton cuts $5$, and
balanced cuts $10,4,6$ in the displayed order.  The unit
binary tree with cherries $ad$ and $bc$ contributes singleton
cuts $1$, and balanced cuts $2,1,2$.  These mincuts add because
the two subgraphs meet only at their terminals: once the
terminal side of a cut is fixed, each internal assignment
can be minimized independently.  This proves the cut values.

Each of the four arms of the tree has length $\sqrt{2/3}$,
and its middle edge has length $\sqrt2-\sqrt6/3$.  Thus the
tree costs $\sqrt2+\sqrt6$, while the direct edges cost
$8+6\sqrt2$.  At each of $x,y$ the three unit vectors along
the incident edges sum to zero.  The total tree length is a
convex function of $(x,y)$, so this stationary point is a
global minimum for the fixed graph.  Finally
$\sqrt2>7/5$ and $\sqrt6>12/5$ already give a value greater
than $20$ in (SH-34.1).  All cut assignments, the exact cost,
and the two zero gradients are independently reproducible
with \path{programs/sh_integer_critical_network.py},
\path{inputs/sh_integer_critical_network.json}, and
\path{outputs/sh_integer_critical_network.json}.
\end{proof}

\begin{lemma}[SH-34b: an integral flow and the odd low cut]
\label{lem:critical-integral-flow}
Suppose an integer-capacity graph on these terminals has
terminal degrees $6$, singleton mincuts $6$, and
$\lambda(ab)=12$, $\lambda(ad)=5$.  There is an integral
$12$-unit flow from $\{a,b\}$ to $\{c,d\}$, saturating
six units at each terminal.  Its terminal-to-terminal path
counts have the form
\[
 n_{ac}=n_{bd}=z,\qquad n_{ad}=n_{bc}=6-z,
 \qquad z\in\{0,1,2\}.
                                                        \tag{SH-34.2}
\]
Every minimum $ad$ cut has at least one unused unit of
capacity in this flow.  If $z=2$, each of the four diagonal
paths crosses such a cut exactly once, and none of the
other eight paths crosses it.
\end{lemma}

\begin{proof}
Attach a source to $a,b$ and a sink to $c,d$, using capacity
$6$ for each new edge.  For an original terminal subset $S$
on the source side, the augmented cut has capacity at least
\[
 \lambda(S)+6\bigl(2-|S\cap\{a,b\}|
                         +|S\cap\{c,d\}|\bigr).
\]
The only nontrivial conditions for this to be at least $12$
are $\lambda(ab)\ge12$ and the four singleton conditions
$\lambda(a),\lambda(b),\lambda(c),\lambda(d)\ge6$.
In particular, no unstated condition on the other balanced
cuts is needed.  The integral maxflow--mincut theorem
(Ford and Fulkerson, 1956, Theorem~1, p.~400; see SH-23)
therefore supplies the flow.  Remove cycles and decompose
it into twelve integral paths.  Terminal degree $6$ and
saturation prohibit any terminal from being internal to
a path.  The four terminal balance equations give (SH-34.2)
with an initially nonnegative integer $z$.

For any minimum $ad$ cut, the total divergence inside is
$6-6=0$.  The number of traversals in the two directions is
therefore equal.  Its total used capacity is even and at
most $5$, hence at most $4$.  The $2z$ diagonal paths must
each cross the cut at least once, proving $z\le2$.  If
$z=2$, these required traversals already use all four
available traversals, proving the final assertion.
\end{proof}

The terminal-degree hypothesis causes no loss in an attempted
width-six counterexample with these cut lower bounds: attach
a zero-length capacity-$6$ edge at each terminal and declare
its new endpoint to be the terminal.  The new singleton
mincuts remain $6$, the $ab$ mincut remains $12$, and the
$ad$ mincut remains at least $5$.  Indeed, changing the
terminal side of a cut costs $6$ per changed label.  For
$ab$, a single changed label leaves a singleton or its
complement, whose old cut costs at least $6$; two changes
already cost $12$.  For a singleton or $ad$, one change
already costs at least its required lower bound.  If the
old $ad$ mincut was exactly $5$, keeping all terminal labels
on their old sides also gives the matching upper bound.

An attractive next step would be to extract three $ad$ paths,
three $bc$ paths, two $ac$ paths and two $bd$ paths, leaving
a connected four-terminal residual network.  Its ten paths
and a four-terminal Steiner tree would yield the lower cost
in (SH-34.1).  The following proposition shows that this
packing statement is false in the original graph.  Consequently
the integral-flow parity observation does not by itself prove
the desired Euclidean lower bound.

\begin{proposition}[SH-35: the residual-packing route fails]
\label{prop:residual-packing-counterexample}
There is an integer-capacity graph with terminal degrees $6$,
singleton mincuts $6$, and balanced mincuts
\[
 \lambda(ab)=12,\qquad \lambda(ad)=5,\qquad \lambda(ac)=12,
\]
which does not contain a capacity-disjoint packing of the ten
paths just described and a connected subnetwork spanning all
four terminals.
\end{proposition}

\begin{proof}
Construct three subgraphs meeting only at $a,b,c,d$.
The first has edges $ad,bc$ of capacity $3$.
The second is the unit binary tree with cherries $ad,bc$.
For the third, take a private unit-capacity $K_4$ on
$a',b',c',d'$ and attach each terminal $v$ to $v'$ by an
edge of capacity $2$.  The three subgraphs contribute the
following terminal cut values:
\[
\begin{array}{c|cccc}
 &\text{each singleton}&ab\mid cd&ad\mid bc&ac\mid bd\\\hline
 \text{two capacity-three edges}&3&6&0&6\\
 \text{unit binary tree}&1&2&1&2\\
 \text{private }K_4\text{ and spokes}&2&4&4&4
\end{array}
\]
For the last row, an internal subset $T$ contributes
$|T|(4-|T|)+2|S\mathbin{\triangle}T|$ to the cut specified
by terminal subset $S$; minimization gives $2$ for a
singleton and $4$ for a balanced pair.  Additivity over
the three independent interiors proves the displayed
mincuts.  Every terminal has degree $3+1+2=6$.

Suppose the proposed packing existed.  At each terminal,
five units are required as endpoints of the ten paths, and
the connected residual requires at least one more incident
unit.  Degree $6$ is therefore saturated.  No path can pass
through a terminal, and each terminal must be a leaf of a
spanning tree extracted from the residual.  Each path, and
this entire residual tree, must consequently lie in a single
one of the three subgraphs.  The two isolated padding edges
cannot contain the residual tree.

If the residual tree is the unit binary tree, it uses that
whole tree.  All six padding paths are forced, leaving two
$ac$ paths and two $bd$ paths to be routed by the private
$K_4$.  Every private vertex then has endpoint load $2$
but internal degree $3$.  A path transit would consume two
additional incidences, which is impossible.  Thus all four
paths would have to use direct opposite edges; each such
edge has capacity only $1$, a contradiction.

If the residual tree is in the private-$K_4$ subgraph, it
uses one unit of every spoke.  Terminal-degree saturation
forces the remaining spoke unit and the unit binary-tree
arm at each terminal to be used by paths.  The padding
again supplies all three $ad$ and all three $bc$ paths.
The binary tree must therefore route one $ac$ and one $bd$
path, both of which require its single middle edge.  This
is again impossible.
\end{proof}

This counterexample has a Euclidean realization of cost
$8+7\sqrt2+\sqrt6$: place its binary-tree vertices as in
SH-34a and all four private-$K_4$ vertices at the origin.
If distinct positions are required, their arbitrarily small
perturbations approach the same cost.  In particular it is
not a network of cost below $20$ and does not refute the
width-six obstruction sought here.

There is also a useful cost-preserving repair of precisely
this counterexample.  For arbitrary private positions
$p_1,\ldots,p_4$, averaging over $j$ gives
\[
 \frac14\sum_{j=1}^4 2\sum_{i=1}^4\|p_i-p_j\|
       =\sum_{i<j}\|p_i-p_j\|.
\]
Choose $j$ no larger than this average.  Replacing the entire
private-$K_4$ subgraph by a capacity-$2$ star centred at $p_j$
does not increase cost, because the triangle inequality gives
\[
 2\sum_i\|v_i-p_j\|
 \le 2\sum_i\|v_i-p_i\|+\sum_{i<k}\|p_i-p_k\|.
\]
This star has the same terminal cut function, namely singleton
values $2$ and balanced values $4$, and it can route two
copies of each opposite terminal pair.  Thus a theorem
allowing cost-reducing replacements remains a possible route;
SH-35 refutes the stronger requirement of finding the packing
inside the given graph.

The counterexample, all seven nontrivial terminal cuts and
its exact cost are recorded in
\path{programs/sh_residual_packing_counterexample.py},
\path{inputs/sh_residual_packing_counterexample.json}, and
\path{outputs/sh_residual_packing_counterexample.json}.
The saved mixed-integer infeasibility result is only an
independent computational cross-check; the preceding
combinatorial proof does not depend on a numerical solver.

Before this counterexample was found, an exploratory search
alternated integer capacity optimization at fixed positions
with convex position optimization for $2,4,6$ internal
vertices and three starts per size.  Its best cost was
$8+7\sqrt2+\sqrt6$; it gives no global lower bound.  The
complete program, input and traces are preserved under
\path{sh_integer_network_search}.  A separate abstract
packing search on graphs with at most five internal vertices
found feasible packings in every generated feasible critical
instance; its data under \path{sh_residual_packing_critical}
also give no universal theorem.  SH-35 has six internal
vertices and supplies an exact obstruction beyond those
tests.

\subsection{Planarity and a weaker residual-packing gate}

\begin{lemma}[SH-36a: parity and planarization]
\label{lem:critical-planarization}
In the setting of SH-34b, $\lambda(ac)\ge8$.
More generally, terminal degrees $m$, singleton mincuts $m$
and $\lambda(ab)=2m$ imply
\[
 \left\lfloor\frac{\lambda(ad)}2\right\rfloor+
 \left\lfloor\frac{\lambda(ac)}2\right\rfloor\ge m.
                                                        \tag{SH-36.1}
\]
For geometric lower bounds, a putative network of cost below
$20$ with the width-six critical cut requirements may be
assumed planar in the terminal diamond, with $a,b,c,d$ on
the outer face in cyclic order.  After attaching the
capacity-$6$ terminal edges described above, its cut values
then satisfy
\[
 \lambda(a)=\lambda(b)=\lambda(c)=\lambda(d)=6,
 \quad \lambda(ab)=12,\quad \lambda(ad)=5,
 \quad \lambda(ac)=12.                                \tag{SH-36.2}
\]
\end{lemma}

\begin{proof}
Use the saturated integral $ab$-to-$cd$ flow.  Its diagonal
path counts are $z,z$ and its other counts are $m-z,m-z$.
Each $ad$ cut has zero net divergence, so its used capacity
is even and at least $2z$.  Each $ac$ cut similarly has
even used capacity at least $2(m-z)$.  This proves
(SH-36.1); with $m=6$ and $\lambda(ad)=5$, it gives
$\lambda(ac)\ge8$.

Project all vertices to the closed terminal diamond using
Euclidean nearest-point projection, which is nonexpansive.
Straighten each edge.  At each segment crossing, split the
edges and identify the intersection vertices.  This costs
nothing.  It can only increase terminal mincuts: a cut in
the new graph assigns one side to every original vertex,
and each original edge whose endpoints lie on different
sides contributes at least its original capacity along
its subdivided path.  Coincident edges can be combined by
adding their capacities.  These operations produce a
planar graph with all four terminals on the outer face.
Attach degree-$6$ terminal leaves on that face at their
terminal positions, using zero-length edges if necessary.
As proved following SH-34b, the singleton and $ab$ cut
values are now exactly $6$ and $12$.

If the $ad$ mincut were at least $6$, (RT-W6.1) would give
cost at least
$12\sqrt2+(2-\sqrt2)6=12+6\sqrt2>20$.
Thus it is exactly $5$ in the putative sub-$20$ case.

Finally let $S$ separate $a,c$ from $b,d$ in this planar
graph.  If the induced graph on $S$ connects $a$ to $c$,
the induced graph on its complement cannot connect $b$
to $d$: the two disjoint connecting paths would join
alternating boundary points, contradicting planar separation.
Consequently at least one of these two induced graphs has
its two terminals in different connected components.
Each of those components contains exactly one terminal,
and its boundary capacity is at least $6$.  Their edge
boundaries are disjoint subsets of $\delta S$, so
$c(\delta S)\ge12$.  The union of the incident edges at
$a,c$ gives the reverse bound, proving the last equality
in (SH-36.2).
\end{proof}

These stronger planar cut values do not remove the integral
issue.  The real-capacity network with $ad=bc=7/2$ and a
central four-arm star of capacity $5/2$ has exactly
(SH-36.2), but still costs $10+7\sqrt2<20$.
It is excluded only by its half-integer capacities.

\begin{lemma}[SH-36b: a sufficient fractional packing]
\label{lem:fractional-residual-gate}
Let $G$ be any Euclidean network on the four terminal points.
Suppose it supports the pair demands
\[
 d_{ad}=d_{bc}=3,\qquad d_{ac}=d_{bd}=2
                                                        \tag{SH-36.3}
\]
by fractional path flows, and, disjoint in capacity from
those flows, a connected unit-capacity subnetwork spanning
all four terminals.  Then its cost is at least
$8+7\sqrt2+\sqrt6>20$.

For an integer graph, existence of such a packing with a
chosen unit tree $T$ is equivalent to
\[
 |\delta T(U)|\le c(\delta G(U))
        -d\bigl(\delta(U\cap\{a,b,c,d\})\bigr)
 \quad\text{for every }U\subseteq V(G),                 \tag{SH-36.4}
\]
provided $T$ uses only positive-capacity edges of $G$.
The existence of a suitable $T$ for all the planar critical
graphs in SH-36a is not established here.
\end{lemma}

\begin{proof}
The fractional pair flows alone cost at least
$6\sqrt2+8$, by the distance between their endpoints.
Any connected residual contains a terminal-spanning tree.
The minimum Euclidean length of such a tree is
$\sqrt2+\sqrt6$.  For completeness, discard nonterminal
leaves and suppress degree-two vertices, straightening
paths.  Treat a terminal that is internal as a leaf joined
by a zero-length edge.  Splitting higher-degree vertices
at zero cost leaves either one degree-four star or two
degree-three vertices.  A star has length at least $4$
by the two antipodal triangle inequalities.  Of the three
binary pairings, the opposite pairing also costs at least
$4$ by those inequalities.  The two adjacent pairings
have the common minimum $\sqrt2+\sqrt6$, proved by the
convex stationary calculation in SH-34a and its rotation.
This proves the length bound and hence the claimed cost.

After reserving one capacity unit on each edge of $T$,
(SH-36.4) is exactly the cut condition for the six-pair
four-terminal multiflow, with the two unlisted demands
zero.  For four terminals this cut condition is sufficient
for a fractional multiflow.  One complete primary proof
is Schrijver (1991), Corollary~2, printed p.~37, with the
$K_4$ case of condition (1) on p.~32; this is also the
four-terminus theorem of Seymour (1980).  Necessity follows
by summing the required path flows across each cut.
\end{proof}

The terminal demand cut values in (SH-36.4) are $5$ for
each singleton, $10$ for $ab$, $4$ for $ad$ and $6$ for
$ac$.  Thus the critical minimum-cut budgets for a tree
are respectively $1,2,1,6$ in the planar case.  These
budgets must hold for \emph{every full graph cut}, not
only for a selected minimizer of each terminal partition.

SH-35 does not refute this weaker gate.  Reserve its binary
tree.  The direct padding edges supply the six side paths,
and the private $K_4$ fractionally routes the four diagonal
units: use each opposite edge once and split the remaining
unit of each opposite demand equally over its two
length-two alternatives.  Every private cycle edge then
carries one unit.  The exact rational certificate is
\path{outputs/sh_fractional_residual_counterexample_repair.json},
generated by \path{programs/sh_fractional_residual_packing.py}
from the SH-35 input.  A subsequent exploratory search
generated $22$ feasible critical graphs with six or seven
internal vertices, all admitting an exactly checked
fractional-flow/integral-residual packing.  Its input and
outputs use the basename \path{sh_fractional_residual_critical}.
These finite tests do not prove a universal packing theorem.

\subsection{An integer parity calibration closes the critical planar gap}

The residual-packing question is unnecessary for the critical
geometric lower bound.  Integer cable widths retain a parity
constraint in a half-integral locking flow.  The following
argument converts that constraint into a positive Euclidean
cost gap, without a planar graph assumption or a path-packing
conjecture.

\begin{theorem}[SH-37: the critical integer-network gap]
\label{thm:critical-parity-calibration}
Let a finite integer-capacity network have terminal positions
$a=(1,0)$, $b=(0,1)$, $c=(-1,0)$ and $d=(0,-1)$.
Suppose each terminal has integer degree $b$, each singleton
mincut is $b$, and
\[
 \lambda(ab)=2b,\qquad \lambda(ad)=L\text{ is odd}.
                                                        \tag{SH-37.1}
\]
Then its Euclidean capacity-length cost satisfies
\[
 w\ge2\sqrt2\,b+(2-\sqrt2)L+2-\sqrt{2+\sqrt2}.
                                                        \tag{SH-37.G}
\]
In particular, $b=6,L=5$ gives
\[
 w\ge12+7\sqrt2-\sqrt{2+\sqrt2}
   =20.051735871589\ldots>20.                         \tag{SH-37.2}
\]
Internal vertices, crossings, multiple edges and zero-length
edges are allowed.  The assertion concerns integer capacities
and is independent of the tensor coefficient field.
\end{theorem}

\begin{proof}
\emph{Step 1: obtain four half-integral commodity flows.}
The family consisting of the four singleton terminal sets
and the two adjacent cuts $ab\mid cd$ and $ad\mid bc$ is
$3$-cross-free: a singleton crosses none of the other sets,
and only two nonsingleton cuts are present.  Double every
capacity.  The doubled graph is inner Eulerian because
every internal weighted degree is even.  The integral
multiflow locking theorem therefore applies to this family
(Hirai, 2010, Theorem~1.1, author preprint p.~2, restating
the Karzanov--Lomonosov theorem).  Divide its integral path
flow by two.  We obtain a feasible half-integral multiflow
locking all six original cuts.

Let $p_{uv}$ be the total flow between terminals $u,v$.
Locking the four singleton cuts gives endpoint load $b$
at each terminal and total multiflow value $2b$.
Locking $ab\mid cd$ at value $2b$ consequently implies
$p_{ab}=p_{cd}=0$.  The four endpoint equations and the
locked $ad\mid bc$ cut then give exactly
\[
 p_{ad}=p_{bc}=b-\frac L2,\qquad
 p_{ac}=p_{bd}=\frac L2.                              \tag{SH-37.3}
\]
The degree-$b$ terminal constraint leaves no capacity for
additional paths passing internally through a terminal.

Orient the four commodities as $a\to d$, $b\to c$,
$a\to c$, and $b\to d$, in that order.  Give every graph
edge an arbitrary reference orientation and write
$f_i(e)\in\tfrac12\mathbb Z$ for the resulting signed
\emph{net} flow of commodity $i$.  Oppositely directed
uses of the same commodity may cancel, so feasibility gives
\[
 \sum_{i=1}^4|f_i(e)|\le c_e,\qquad c_e\in\mathbb Z_{\ge0}.
                                                        \tag{SH-37.4}
\]
Its signed divergence is still the demand in (SH-37.3).

\emph{Step 2: compare every edge with a Euclidean calibration.}
The four terminal direction vectors are
\[
 u_1=u_2=\frac{(-1,-1)}{\sqrt2},\qquad
 u_3=(-1,0),\qquad u_4=(0,-1).
\]
For a positively oriented edge displacement $\Delta_e$,
flow conservation and telescoping give
\[
 \sum_e\sum_i f_i(e)\,u_i\cdot\Delta_e
       =\sqrt2(2b-L)+2L.                            \tag{SH-37.5}
\]
For an edge of positive length, put
$t_e=\Delta_e/\|\Delta_e\|$ and
$\alpha_i(e)=2f_i(e)\pmod2\in\mathbb F_2$.
Write $h=\sum_i\alpha_i(e)$ and
$|f_i(e)|=n_i+\alpha_i(e)/2$, with $n_i$ nonnegative
integers.  Since $c_e$ is an integer, (SH-37.4) implies
$c_e\ge\sum_i n_i+\lceil h/2\rceil$.  Hence
\begin{align*}
 c_e-\sum_i f_i(e)u_i\cdot t_e
 &\ge c_e-\sum_i|f_i(e)|\,|u_i\cdot t_e|\\
 &\ge \left\lceil\frac h2\right\rceil
     -\frac12\sum_{i:\alpha_i(e)=1}|u_i\cdot t_e|\\
 &\hspace{1.4em}
     +\sum_i n_i\bigl(1-|u_i\cdot t_e|\bigr)\\
 &\ge g_{\alpha(e)}(t_e),
 \end{align*}
where
\[
 g_\alpha(t)=\left\lceil\frac{|\alpha|}2\right\rceil
                  -\frac12\sum_{i:\alpha_i=1}|u_i\cdot t|.
                                                        \tag{SH-37.6}
\]
Every $g_\alpha$ is nonnegative.  More precisely, let
\[
 \gamma=1-\frac{\sqrt{2+\sqrt2}}2>0.
\]
Then
\[
 \alpha\notin\{0000,1100\}
       \quad\Longrightarrow\quad g_\alpha(t)\ge\gamma
       \quad\text{for every unit }t.                 \tag{SH-37.7}
\]
Here is the complete parity-case check.  If $|\alpha|$ is
odd, the bound $|u_i\cdot t|\le1$ gives $g_\alpha\ge1/2$.
If $|\alpha|=2$ and $\alpha\ne1100$, its two directions
are either perpendicular or at angle $\pi/4$.  For any
unit vectors $u,v$,
\[
 |u\cdot t|+|v\cdot t|
 \le\max\{\|u+v\|,\|u-v\|\}
 =\sqrt{2+2|u\cdot v|}
 \le\sqrt{2+\sqrt2},
\]
which gives (SH-37.7).  If $|\alpha|=4$, then
\[
 2|u_1\cdot t|+|u_3\cdot t|+|u_4\cdot t|
       \le2+\sqrt2,
\]
and $g_\alpha\ge1-1/\sqrt2>\gamma$.
The remaining labels $0000$ and $1100$ merely give the
nonnegative bounds $0$ and $1-|u_1\cdot t|$.

\emph{Step 3: parity forces a positive-length support.}
Take the quotient group
\[
 Q=\mathbb F_2^4/\langle e_1+e_2\rangle
\]
and retain precisely those edges whose parity label
$\alpha(e)$ is nonzero in $Q$.  Call the resulting
subgraph $R$.  At every nonterminal its $Q$-valued
boundary is zero, since the doubled net flows have zero
integer divergence there.  The terminal charges, using
the odd doubled demands $2b-L,2b-L,L,L$, are
\[
 a,c:\ X=e_1+e_3,\qquad
 b,d:\ Y=e_1+e_4
 \quad\text{in }Q.
\]
The two vectors $X,Y$ are linearly independent in $Q$.
The sum of the boundary charges in each connected
component of $R$ is zero: each internal edge contributes
twice, and discarded edges have zero quotient label.
Consequently the component containing $a$ must also
contain $c$.  It follows that
\[
 \sum_{e\in E(R)}\|\Delta_e\|\ge\|a-c\|=2.           \tag{SH-37.8}
\]
This argument retains zero-length edges in $R$; they can
join coincident vertices but cannot reduce the length of
an $a$-to-$c$ path below the endpoint distance.  On such
edges both cost and calibration contributions are zero,
so no unit direction need be defined.

Summing (SH-37.6)--(SH-37.8) and using (SH-37.5) proves
\[
 w-\bigl(\sqrt2(2b-L)+2L\bigr)\ge2\gamma,
\]
which is (SH-37.G).  In the particular case $b=6,L=5$,
this is (SH-37.2).  To verify its strict final comparison
without decimal approximations, note that $7\sqrt2-8>0$
and
\[
 (7\sqrt2-8)^2-(2+\sqrt2)=160-113\sqrt2>0,
\]
because $160^2-2\cdot113^2=62>0$.
\end{proof}

The finite angular and parity calculations are independently
reproducible with \path{programs/sh_parity_calibration.py},
\path{inputs/sh_parity_calibration.json}, and
\path{outputs/sh_parity_calibration.json}.  The program
checks all sixteen parity labels and every sign choice
that determines each angular maximum, using exact
quadratic-number comparisons.  It also verifies that the
only zero-sum terminal subsets in $Q$ are the empty set,
$\{a,c\}$, $\{b,d\}$ and all four terminals, and checks
the integer margin $62$.  The graph argument above,
rather than numerical optimization, is the certificate
for the universal lower bound.

\begin{corollary}[SH-38: planar widths six and seven]
\label{cor:planar-width-six-seven}
Over every field, a strict shortening in dimension $D=2$
has width at least $7$.  If it has exactly two summands,
its width is at least $8$.
\end{corollary}

\begin{proof}
Widths at most five are already excluded by SH-28d.
We first treat two summands at widths $m=6,7$, writing
$q=2^m$.  Attach zero-length identity cables of width $m$
at all four outputs of each term network.  This does not
alter the computed tensors or cost and makes every
terminal degree exactly $m$.  Both adjacent network
mincuts are therefore at most $2m$.  Their sum is at
most the network cost, which is strictly below $4m-4$.

On either crossing cut the target has rank $2^{2m}$.
If both term rank ceilings on one cut were at most
$2m-1$, both ranks would have to equal $2^{2m-1}$.
On the other cut some term has ceiling at least $2m-1$;
that network would have adjacent capacities summing to
at least $4m-2$, a contradiction.  Thus each cut has a
term of ceiling $2m$, and these high cuts must belong
to opposite terms.  Each such network has high mincut
exactly $2m$.  Since this cut is at most the sum of
the singleton mincuts on either side, all four singleton
mincuts equal $m$.

Let $L$ be its other adjacent mincut.  The continuous
bound (RT-W6.1) excludes $L\ge2m-6$, because then
\[
 w\ge2\sqrt2m+(2-\sqrt2)(2m-6)
       =4m-12+6\sqrt2>4m-4.
\]
Hence $L\le2m-7$.  Equality is odd and is excluded by
SH-37, whose lower bound at this value exceeds $4m-4$
by the same positive constant
$7\sqrt2-8-\sqrt{2+\sqrt2}$ as in (SH-37.2).
Both low capacities are consequently at most $2m-8$.
For $m=6$, their ranks sum to at most $16+16<64$,
contradicting the elementary mixed-slice inequality.
For $m=7$, their ranks sum to at most $64+64=128$,
contradicting PG-SLICE-CLIFFORD, which requires at least
$q+1=129$.  These arguments allow arbitrary term networks.

It remains to exclude three summands at $m=6$; four or
more are excluded by RT-W6-PARAM and one by Crofton.
The cost budget is
\[
 w<24-4\log_2 3<\frac{53}{3},
\]
using $3^{12}>2^{19}$.  For the adjacent rank ceilings
$(H,L)$ of a term, with $H\ge L$, (RT-W6.1) implies
\[
 H=12\ \Longrightarrow\ L\le1,\qquad
 H=11\ \Longrightarrow\ L\le3.                       \tag{SH-38.1}
\]
Indeed the first excluded costs are $4+10\sqrt2$ and
$8+7\sqrt2$, respectively, both greater than $53/3$.
If no term has ceiling $12$, each has sum of the two
adjacent ranks at most $2^{11}+2^3=2056$: for maximum
ceiling at most $10$, the coarser bound $2\cdot2^{10}$
already suffices.  The three terms would then contribute
at most $6168<8192$ to the sum of the two required
target flattening ranks.  Therefore one term $U$ has
ceiling $12$ on one cut and rank at most $2$ on the other.

If no other term has ceiling $12$ on that other cut,
the two remaining terms must both have ceiling $11$
there, to supply rank at least $4096-2=4094$.
By (SH-38.1) their ranks on the first cut are at most
$8$ each.  This contradicts PG-CROSSCUT-RESIDUAL:
removing a tensor of rank at most $2$ on the other cut
leaves rank at least $64\cdot63/2=2016$ on the first.
Thus a second term $S$ has the opposite ceiling $12$
and also has low rank at most $2$.

Apply PG-CROSSCUT-RESIDUAL separately to $U$ and $S$.
The third term must have rank at least $2016-2=2014$
on \emph{both} cuts, so both ceilings are at least $11$.
The planar frame bound gives cost at least $22$, again
contradicting $w<53/3$.  This excludes all term counts
at planar width six and completes the proof.
\end{proof}

For planar width six with two summands there is also an
independent algebraic proof: the continuous bound leaves
low capacities at most $5$, so the sum of the two cheaper
ranks is at most $64$, already contrary to the strict
$q+1$ bound of PG-SLICE-CLIFFORD.  SH-37 supplies a
stronger geometric statement and is needed above to
exclude the next two-summand width, $m=7$.

\subsection{Rectangular folds and the first higher-dimensional exclusion}

\begin{lemma}[SH-39: rectangular crossing constraints]
\label{lem:rectangular-fold-constraints}
Let $C_{p,q}$ be the tensor product of opposite identity wires
of dimensions $p,q$, and let $C_{p,q}=S+U$ have ranks $r,s$
on its two different adjacent cuts.  Over every field,
\[
 r+t^2s\ge tp,\qquad t^2r+s\ge tp
       \quad(1\le t\le q).                          \tag{SH-39.1}
\]
For $q\ge2$, also $r+3s\ge2p$ and $3r+s\ge2p$.
If $2^q>2p$, then $r+s\ge p+1$.

For two sign cuts of a $D$-diameter tensor of width $m$ at
folded Hamming distance $h\in\{1,\ldots,\lfloor D/2\rfloor\}$,
these inequalities apply with
\[
 p=2^{m(D-h)},\qquad q=2^{mh}.                        \tag{SH-39.2}
\]
In particular the strict sum bound applies whenever
$2^{mh}>m(D-h)+1$.
\end{lemma}

\begin{proof}
Contract the two $q$-dimensional opposite parties against
basis vectors $i,j$.  Rank decompositions of $S,U$ give
matrices satisfying
\[
 X_iY_j^T+Z_jW_i^T=\delta_{ij}I_p,
\]
where $X_i,Y_i$ have $r$ columns and $Z_i,W_i$ have $s$
columns.  Stack $t$ selected $X_i$ and $Y_i$.  Their
product has rank at most $r$, while its difference from
$I_{tp}$ is a $t$-by-$t$ block matrix whose individual
blocks have rank at most $s$.  Its rank is at most
$t^2s$, proving the first inequality in (SH-39.1).
Exchange the two cut orientations for the second.
For $i\ne j$, the diagonal identities give
$\operatorname{rank}X_i,\operatorname{rank}Y_j\ge p-s$,
while the off-diagonal identity gives
$\operatorname{rank}(X_iY_j^T)\le s$.
Sylvester's inequality yields $r+3s\ge2p$, and symmetry
gives its partner.  This is the rectangular version of
PG-CROSSING-UNBALANCED.

For the strict sum assertion, the slice construction
gives homogeneous linear matrices $L(x,z),R(x,z)$ of
sizes $p$-by-$(r+s)$ and $(r+s)$-by-$p$ such that
$LR=(x\cdot z)I_p$, with $x,z\in F^q$.
The ordinary slice bound gives $r+s\ge p$.  If equality
held, the matrices would be square and would also satisfy
$RL=(x\cdot z)I_p$.  The creation--annihilation proof in
PG-SLICE-CLIFFORD applies on the doubled space of dimension
$2p$: its $q$ commuting idempotents have $2^q$ joint
eigenspaces of the same positive dimension.  Therefore
$2^q$ divides $2p$, in particular $2^q\le2p$, a contradiction.
This proof includes characteristic two.

Finally, merge original parties according to the pair of
signs assigned by the two cuts.  The $D-h$ agreeing axes
form one opposite wire of dimension $p$, and the $h$
disagreeing axes form the other, of dimension $q$.
Merging within cut blocks preserves the relevant
flattening ranks, giving (SH-39.2).
\end{proof}

\begin{theorem}[SH-40: width six, dimension three, two summands]
\label{thm:width-six-dimension-three}
Over every field there is no strict shortening with
$m=6,D=3,A=2$.
\end{theorem}

\begin{proof}
There are four sign cuts, modulo simultaneous negation,
and every pair of distinct cuts has folded Hamming
distance one.  Write $k_S(s),k_U(s)$ for the two summands'
integer rank ceilings on these cuts.  They are at most
$18$ by the output dimensions.  SH-23a gives, for each
network of cost $w<30$,
\[
 w\ge\frac24\sum_s k(s),
 \qquad\text{hence}\qquad \sum_s k(s)\le59.          \tag{SH-40.1}
\]
This is the sign tight-frame/coarea bound and allows
arbitrary internal vertices and layouts.

The combined ceiling sum is at most $118$, so at some
cut its two ceilings sum to at most $29$.  The target
rank there is $2^{18}$.  If neither ceiling were $18$,
the two ranks could reach this value only if both
ceilings were $17$, whose sum would be $34$.
Thus one ceiling is $18$ and the other is at most $11$.
Choose the absolute minimum $t$ among all eight ceilings;
after relabelling terms and cuts let $k_S(0)=t\le11$.
The target rank forces $k_U(0)=18$.

For any of the other three cuts $j$, SH-39 folds the
identity to $C_{4096,64}$ and applies to ranks bounded
above by $2^t$ and $2^{k_U(j)}$.
If $t\le8$, use (SH-39.1) with parameter $8$:
\[
 8^2\,2^t+2^{k_U(j)}\ge8\cdot4096
 \quad\Longrightarrow\quad k_U(j)\ge14.
\]
Then $\sum_s k_U(s)\ge18+3\cdot14=60$, contradicting
(SH-40.1).

If $t=9$, parameter $4$ instead gives
$2^{k_U(j)}\ge4\cdot4096-16\cdot512=8192$, so all
three other $U$ ceilings are at least $13$.
Their baseline sum with $k_U(0)$ is $57$.
The budget $59$ prohibits any of them from being $17$
or larger, which would add at least four.  Thus all
three are at most $16$, forcing all three corresponding
$S$ ceilings to be $18$.  Consequently
$\sum_s k_S(s)\ge9+54=63>59$.

Finally let $t=10$ or $11$.  Since
$2^{64}>2\cdot4096$, the strict rectangular sum bound
forbids both prescribed ceilings from being at most
$11$.  Hence all other $U$ ceilings are at least $12$.
Their baseline total is $18+3\cdot12=54$.
The budget $59$ excludes any other ceiling $18$, which
would add six, and permits at most one ceiling $17$,
since two would add ten.  At least two corresponding
$S$ ceilings are therefore $18$, and the third is at
least $17$.  Their total with $t$ is at least
$t+18+18+17>59$, the final contradiction.
\end{proof}

The strict rectangular inequalities alone were also tested
in the saved necessary-profile programs
\path{programs/sh_strict_sign_rank_profiles.py} and
\path{programs/sh_strict_radial_rank_certificate.py}.
They admit exact feasible integer profiles in dimensions
three and four; thus they do not by themselves prove
SH-40.  The unbalanced inequalities used in the human
proof are essential to this deduction.  Later necessary
profile experiments use
\path{programs/sh_unbalanced_radial_rank_certificate.py}.
Only independently verified exact certificates or explicit
proofs are treated as exclusions; a solver's status for
the dual-certificate search is not an infeasibility
certificate for the original tensor problem.

\begin{theorem}[SH-41: width six, dimension four, two summands]
\label{thm:width-six-dimension-four}
Over every field there is no strict shortening with
$m=6,D=4,A=2$.
\end{theorem}

\begin{proof}
The eight sign cuts modulo negation form the folded
four-cube, whose distance-one graph is $K_{4,4}$.
Write $S,U$ for the summands and $k_S,k_U$ for their
rank ceilings.  The output dimensions bound every
ceiling by $24$.  Each network would have cost less
than $40$, so the SH-23a frame bound gives
\[
 \sum_s k_S(s)\le159,\qquad
 \sum_s k_U(s)\le159.                         \tag{SH-41.1}
\]
At some cut the two ceilings therefore sum to at most
$39$.  Since the target rank is $2^{24}$, one ceiling
there must be $24$ and the other at most $15$.
Choose the absolute minimum $t$ among all sixteen
ceilings, and arrange $k_S(0)=t\le15$.
Then $k_U(0)=24$.  From the marked cut there are four
cuts at distance one and three at distance two.

First suppose $t\le8$.  On every distance-one cut,
SH-39 with $p=2^{18},q=64$ and parameter $64$ gives
\[
 2^{k_U(s)}\ge2^{24}-2^{12+t}>2^{23},
\]
so these four ceilings are $24$.  On each distance-two
cut use $p=q=2^{12}$ and parameter $2^{11-t}$:
\[
 2^{k_U(s)}\ge 2^{23-t}-2^{22-t}=2^{22-t}.
\]
Thus $\sum_s k_U(s)\ge186-3t>159$.

For $t=9$ or $10$, the same five ceilings of $U$ are
$24$.  Its other three ceilings sum to at most $39$;
each is at least $t$, hence each is at most $21$, and
one is at most $13$.  The target ranks force all three
corresponding ceilings of $S$ to be $24$.
Every distance-two cut is adjacent to all four
distance-one cuts.  Apply SH-39 from a $U$ ceiling
at most $13$, using $p=2^{18}$ and parameter $16$.
All four of those $S$ ceilings are at least $21$.
Consequently $\sum_s k_S(s)\ge t+4\cdot21+3\cdot24
=t+156>159$.

It remains to treat $11\le t\le15$.  Across an edge
of $K_{4,4}$, any two ceilings $a,b$ belonging to
different summands satisfy
\[
                         a+b\ge34.          \tag{SH-41.2}
\]
Indeed, this is immediate if both are at least $17$.
Otherwise let $b=\min(a,b)\in[11,16]$ and apply
SH-39 with integer parameter $2^{17-b}\in[2,64]$.
It gives $2^a\ge2^{34-b}$, proving (SH-41.2).

At each cut at least one ceiling is $24$, or both are
$23$.  Let $x$ count the latter exceptional cuts and
$y$ count cuts where both ceilings are $24$.
For a cut with exactly one $24$, call the other ceiling
the low ceiling and its difference from $t$ its
deviation.  Relative to the baseline $24+t$ at each
of the eight cuts, the combined budget allows excess
\[
 R=126-8t.
\]
An exceptional cut uses $22-t$ of this excess and a
double-$24$ cut uses $24-t$.  The remaining excess is
the sum of all low deviations.  On an edge joining
opposite single-$24$ types, (SH-41.2) requires that
the two low deviations sum to at least
\[
 d=34-2t.                                      \tag{SH-41.3}
\]

Let $a$ be the number of $24$ ceilings of either
fixed summand.  Its individual budget gives
\[
 (24-t)a+(23-t)x\le159-8t.                    \tag{SH-41.4}
\]
Together with the total excess bound, this leaves
only
\[
 (x,y)\in\{(0,0),(0,1),(0,2),
                 (1,0),(1,1),(2,0)\}.
\]
For completeness, $x\ge4$ already exceeds $R$;
$x=3$ requires $t\le12$, when (SH-41.4) permits
at most two $24$ ceilings per summand although at
least five are needed in total.  For $x=2$ each
summand has at most three $24$ ceilings, forcing
$y=0$; for $x=1$ each has at most four, forcing
$y\le1$; and for $x=0$ each has at most five,
forcing $y\le2$.

The two sets of opposite single-$24$ types are
disjoint vertex sets in $K_{4,4}$.  Disjoint sets
of sizes $4,4$ have a matching of size $4$; sizes
$4,3$ have one of size $3$; and sizes $3,3$ have
one of size at least $2$.  These assertions follow
by counting the vertices of each set in the two
bipartition classes: the matching size is the sum
of the two corresponding minimum class sizes.
Summing (SH-41.3) over such a matching counts every
low deviation at most once.  The following table
lists the resulting excess over the entire budget
$R$; every entry is positive for $11\le t\le15$.
\[
\begin{array}{c|c|c|c}
(x,y)&\text{single-type sizes}&\text{matching size}
 &\text{required excess}-R\\\hline
(0,0)&4,4&4&10\\
(0,1)&4,3&3&t\\
(0,2)&3,3&2&2t-10\\
(1,0)&4,3&3&t-2\\
(1,1)&3,3&2&2t-12\\
(2,0)&3,3&2&2t-14
\end{array}
\]

The only omitted case has $(x,y)=(0,0)$ with
single-type sizes $5,3$.  Suppose $S$ has five
ceilings equal to $24$.  The deviations at its
other three cuts have total at most $39-3t$;
the five low deviations of $U$ have total at most
$87-5t$.  If the smaller type has $b$ vertices in
one bipartition class, its cross-type edge graph
is the disjoint union
$K_{b,b+1}\sqcup K_{3-b,4-b}$, for $b=0,1,2,3$.
One can select four edges with distinct endpoints
in the larger type and with each endpoint in the
smaller type used at most twice.  (For $b=0,3$
use four edges of $K_{3,4}$; for $b=1,2$ use two
from each nonempty component.)  Summing (SH-41.3)
over them shows that the five low deviations of
$U$ total at least
\[
 4d-2(39-3t)=58-2t>87-5t,
\]
where the last difference is $3t-29>0$.
This final contradiction proves the theorem.
\end{proof}

The saved radial branch certificates at width six and
dimension four cover every minimum $t=0,\ldots,15$
except $t=12$; they are supplementary checks, not
premises of the analytic proof SH-41.  Their inputs,
exact rational outputs, and batch manifest are
\path{inputs/sh_unbalanced_minimum_batch_D4_D6.json},
\path{outputs/sh_unbalanced_minimum_batch_D4_D6.json},
and the files with prefix
\path{sh_unbalanced_m6_D4_min} in the corresponding
input and output directories.

\begin{theorem}[SH-42: width six in dimensions five and six]
\label{thm:width-six-dimensions-five-six}
Over every field there is no strict shortening with
$m=6,A=2$ and $D\in\{5,6\}$.  This assertion concerns
arbitrary finite network layouts, including internal
vertices and zero-length edges.
\end{theorem}

\begin{proof}
Use the same necessary rank-profile model as SH-28,
augmented by SH-39.  There are $M=2^{D-1}$ sign cuts,
and each nonzero summand has integer rank ceilings
$k(s)\in\{0,\ldots,6D\}$.  The sign-frame bound gives
\[
 \sum_s k(s)\le5DM-1.                         \tag{SH-42.1}
\]
Changing one sign can change a ceiling by at most
$12$.  For two cuts at folded distance $h$, the
mixed rank-sum bounds and all selected rectangular
inequalities of SH-39 are necessary.

The combined budget implies that the absolute
minimum of the two profiles is an integer
$t\in\{0,\ldots,4D-1\}$.  Indeed, at some cut the
sum of the two ceilings is at most $10D-1$.
The target rank $2^{6D}$ then forces one ceiling
to be $6D$, since the alternative with neither
full ceiling requires both to equal $6D-1$.
The other ceiling is at most $4D-1$.
Mark a cut attaining the absolute minimum and
average threshold indicators over its coordinate
permutation stabilizer.  Each orbit is a folded
Hamming layer, with its exact binomial weight.
For a surviving branch also mark a minimum of
the other profile.  If its value is $y$, then
$t\le y\le5D-1$ by (SH-42.1).  The orbits of
the two marked cuts are indexed by weights in
the two coordinate blocks, modulo complementation.
This averaging preserves every linear necessary
inequality used in the certificates.

The new rectangular rows have the following
explicit meaning.  If the upper rank bounds
$2^a,2^b$ violate an instance of SH-39, the two
ceilings cannot simultaneously be at most $a,b$.
For threshold indicators $z_{j,s,k}=1[k_j(s)\ge k]$,
the resulting row is
\[
 -z_{0,s,a+1}-z_{1,u,b+1}\le-1.                \tag{SH-42.2}
\]
Its label records the folded distance, both
ceiling bounds, the inequality family, its
orientation, and its integer parameter.  The
strict rectangular sum row is likewise included
only when its explicit hypothesis holds.
All other rows are the budget, monotonicity,
box constraints, the rank-change bound, and the
marked global-minimum constraints already
described in SH-28.

The exact certificate coverage is below.  The last
three columns count radial certificates, two-minimum
certificates, and pruned two-minimum branches,
respectively.
\[
\begin{array}{c|c|c|c|c}
D&\text{first minima}&c_{\rm rad}&c_{\rm two}&c_{\rm prune}\\\hline
5&0,\ldots,19&18&17&37\\
6&0,\ldots,23&17&121&215
\end{array}
\]
For $D=5$ the radial exceptions are $t=15,17$;
for $D=6$ they are $t=15,16,\ldots,21$.
Every remaining admissible second-minimum value
and folded distance has a certificate.  The
listed pruned branches violate the ordinary
mixed-cut rank-sum inequality at the two marked
minima, checked with integer powers of two.
Each certificate is an explicit nonnegative
rational combination of necessary rows whose
left sides cancel and whose right side is $-1$.
It therefore rules out its branch over every
field, without relying on a numerical solver's
infeasibility status.

An independent root verifier reconstructs the
folded-cube orbit geometry and every row from
its label, checks each rectangular violation
with integer arithmetic, and recomputes all
rational cancellations and branch coverage.
The complete checks passed for both dimensions.
This covers every possible absolute minimum,
and proves the theorem.
\end{proof}

The generating programs are
\path{programs/sh_unbalanced_radial_rank_certificate.py}
and \path{programs/sh_unbalanced_two_minimum_certificate.py}.
The two-minimum manifests are
\path{outputs/sh_unbalanced_two_min_D5_m6_manifest.json}
and \path{outputs/sh_unbalanced_two_min_D6_m6_manifest.json};
the selected radial inputs and outputs have prefixes
\path{sh_unbalanced_m6_D5_min} and
\path{sh_unbalanced_m6_D6_min}.
The independent checker is
\path{programs/rt_verify_unbalanced_certificates.py},
with complete input/output pairs
\path{inputs/rt_unbalanced_D5_verification.json},
\path{outputs/rt_unbalanced_D5_verification.json},
\path{inputs/rt_unbalanced_D6_verification.json}, and
\path{outputs/rt_unbalanced_D6_verification.json}.

\begin{theorem}[SH-43: width six, dimension seven, two summands]
\label{prop:width-six-dimension-seven-certificates}
For $m=6,A=2,D=7$, the necessary model in the proof
of SH-42 has complete exact certificate coverage.
There is no strict shortening over any field or
any finite network layout with these parameters.
\end{theorem}

\begin{proof}[Exact certificate proof]
The possible absolute minima are $t=0,\ldots,27$.
Twenty-one radial certificates exclude all values
except $20,21,\ldots,26$.  Marking the second
profile's minimum gives $336$ cases before pruning.
Of these, $226$ violate the ordinary mixed-cut
rank sum, one violates its strict rectangular
version, and $33$ violate a recorded unbalanced
rectangular inequality from SH-39.  The other
$76$ cases each have an exact nonnegative rational
combination of necessary rows with complete
cancellation and right side $-1$.
All pruning decisions use integer powers of two;
the rectangular records include the inequality
family, orientation, and integer parameter.
Thus these certificates cover every first minimum,
second minimum, and folded distance permitted by
the geometric rank budgets.
\end{proof}

The complete manifest is
\path{outputs/sh_unbalanced_two_min_D7_m6_pruned_manifest.json},
generated from
\path{inputs/sh_unbalanced_two_min_D7_m6_pruned_batch.json}
by \path{programs/sh_unbalanced_two_min_pruned_batch.py}.
The radial outputs have prefix
\path{outputs/sh_unbalanced_m6_D7_min}.
All $76$ rational outputs pass the generator's
exact checker.  The root's separate independent
checker also reconstructed every row and every
rectangular pruning decision and verified full
coverage of $t=0,\ldots,27$.  Its input and output
are \path{inputs/rt_unbalanced_D7_verification.json}
and \path{outputs/rt_unbalanced_D7_verification.json}.
No numerical solver status is used as a premise.

\begin{corollary}[SH-44: stronger rectangular rank-profile rows]
\label{cor:singh-sign-cut-constraints}
Consider two distinct sign cuts at folded distance $h$
for a width-six diameter tensor in dimension $D$.
Put $L=6(D-h)$.  Every decomposition into arbitrary
summands $T_1,\ldots,T_A$ obeys
\[
 \sum_{j=1}^A\min\{
   \operatorname{rank}_{s}(T_j),
   \operatorname{rank}_{u}(T_j)\}
 \ge\left\lceil\frac{67\,2^L}{49}\right\rceil.
                                                        \tag{SH-44.1}
\]
For two summands, prescribed opposite-cut ceiling
pairs $(L,L-2)$ and $(L-2,L)$ are therefore impossible.
\end{corollary}

\begin{proof}
The rectangular fold of SH-39 has
$p=2^{6(D-h)}$ and $q=2^{6h}$, both at least $64$.
Apply the independently audited rectangular extension
of Singh's seven-index subspace certificate,
Corollary~\ref{pg:singh-rectangular}.
For a sum of more than two tensors, group each
summand according to its cheaper cut.  Rank
subadditivity bounds the two resulting prescribed
ranks by the sum on the left of (SH-44.1), so the
same lower bound applies.  Finally,
$2^L+2^{L-2}=5\cdot2^L/4<67\cdot2^L/49$.
\end{proof}

The stronger rows in SH-44 were established after
the completed SH-42/SH-43 certificate runs.  Those
exclusions use only SH-39 and do not depend on this
later strengthening.  Further searches can include
the forbidden ceiling pairs and, for three or more
summands, the full cheaper-rank sum inequality.

\begin{lemma}[SH-45: a smaller three-summand minimum]
\label{lem:three-summand-minimum-eight}
In a hypothetical strict shortening with $m=6,D=3,A=3$,
the absolute minimum of the twelve sign-cut rank
ceilings is at most $8$.
\end{lemma}
\begin{proof}
There are four sign cuts, and each ceiling is at most
$18$.  The cost budget and SH-23a give, for each term,
\[
 \sum_s k(s)<72-12\log_2 3<53,
 \qquad\text{so}\qquad \sum_s k(s)\le52.
\]
The strict comparison uses $3^{12}>2^{19}$.
Give a ceiling weight $0$ when it is at most $15$,
weight $1$ at $16$, weight $2$ at $17$, and weight
$4$ at $18$.  At each cut the three ranks must sum
to at least $2^{18}$, so their weights sum to at
least $4$.  Indeed, without a ceiling $18$ there
must be either two ceilings $17$ or a triple
dominating $(17,16,16)$.

Suppose all twelve ceilings were at least $9$.
A term with a ceiling $18$ can have no second $18$
and no $17$, since their minimum row sums would
be $54$ and $53$, respectively.  It can have at
most one $16$, so its weight is at most $5$.
A term without an $18$ can have at most two
ceilings at least $16$, since three would cost
at least $3\cdot16+9=57$; its weight is at most
$4$.  Thus the total weight is at most $15$,
contrary to the four cuts' required total $16$.
\end{proof}

\begin{theorem}[SH-46: width six, dimension three, three summands]
\label{thm:width-six-dimension-three-three-summands}
Over every field, no strict shortening has
$m=6,D=3,A=3$, for arbitrary finite network layouts.
\end{theorem}

\begin{proof}
Write the three rank-ceiling profiles as $S,U,V$.
Each has four entries in $\{0,\ldots,18\}$ and
sum at most $52$.  By SH-45 their absolute minimum
is $t\in\{0,\ldots,8\}$; put it at $S(0)$.
The four sign cuts have pairwise folded distance
one, so all distinct pairs fold to $C_{4096,64}$.

We first specify the finite reduction and its exact
certificate.  At a cut, the three ceilings must
dominate one of the nine distinct permutations of
\[
 (18,t,t),\qquad(17,17,t),\qquad(17,16,16).
                                                        \tag{SH-46.1}
\]
This follows directly from their rank sum being at
least $2^{18}$.  At the marked cut, after exchanging
$U,V$ if necessary, only the lower patterns
$(t,18,t)$ and $(t,17,17)$ are needed.
The other three cuts may be permuted, since every
distinct pair has the same folded distance and all
the necessary constraints are invariant under that
permutation.  Choose their pattern indices in
nondecreasing order.  There are
$2\binom{11}{3}=330$ branches for each $t$, and
$2970$ branches in total.  These are lower-bound
patterns; they do not assume that an actual ceiling
equals its selected lower pattern.

The exact propagation certificate maintains integer
lower and upper bounds for every ceiling.  Initially
the lower bounds are the selected patterns, the
upper bounds are $18$, and $S(0)=t$ is fixed.
Every recorded deduction is one of the following:
\begin{enumerate}
\item A row budget gives an upper bound of $52$
minus the other three lower bounds.  The rank-change
bound across two distinct cuts transfers bounds
with an additive allowance of $12$.
\item The single-cut rank sum is at least $2^{18}$.
Upper bounds for two terms therefore give a residual
rank lower bound for the third, rounded upward to
an integer binary logarithm.
\item For distinct cuts $s,u$, the rectangular
inequalities imply
\[
 \sum_{j=1}^3\min\{2^{k_j(s)},\lambda2^{k_j(u)}\}
 \ge T,                                           \tag{SH-46.2}
\]
with $(\lambda,T)=(1,5601),(3,8192)$, or
$(z^2,4096z)$ for any integer $2\le z\le64$.
For the weighted forms, group each summand by its
cheaper weighted cut and apply SH-39; for the
first form apply SH-44.  Upper bounds for the
other two terms leave a required contribution
$R$ from the selected term.  If $R>0$, both its
rank and its weighted rank must be at least $R$,
giving the two recorded ceiling lower bounds.
\end{enumerate}
The certificate stops a branch only when a lower
bound exceeds its upper bound, a row's lower
bounds exceed $52$, or the upper bounds violate
one of these necessary rank inequalities.  All
operations use integers and upward binary-logarithm
rounding; no numerical optimization is involved.

The independently reconstructed certificate checks
$4364$ such deductions and excludes $2967$ branches.
All $330$ branches are excluded for each $t=0,\ldots,5$.
For each $t=6,7,8$, exactly one lower pattern survives:
\[
\begin{array}{c|cccc}
 &0&1&2&3\\\hline
S&t&s&18&18\\
U&18&u&a_2&a_3\\
V&v&18&b_2&b_3
\end{array}
\qquad t\le s\le16-t.                         \tag{SH-46.3}
\]
The remaining entries have only their global
lower bound $t$ and output upper bound $18$ at
this point.  We finish these three cases analytically.

The last two row budgets in (SH-46.3) give
\[
 u+a_2+a_3\le34,\qquad v+b_2+b_3\le34.
                                                        \tag{SH-46.4}
\]
Put $B_t=5601-2^t$ and $C_t=5601-2^{16-t}$.
Both exceed $4096$ for $t=6,7,8$.
The cheaper-rank bound on the pairs of cuts
$(0,1),(0,j),(1,j)$ gives, for $j=2,3$,
\begin{align}
 2^u+2^v&\ge B_t,\nonumber\\
 2^{a_j}+\min\{2^v,2^{b_j}\}&\ge B_t,\nonumber\\
 2^{b_j}+\min\{2^u,2^{a_j}\}&\ge C_t.
                                                    \tag{SH-46.5}
\end{align}
If $v\le11$, the first two lines force
$u,a_2,a_3\ge12$, contradicting (SH-46.4).
If $u\le11$, the first and third lines force
$v,b_2,b_3\ge12$, giving the other contradiction.
Thus $u,v\ge12$.

Apply (SH-46.2) to cuts $j,0$ with
$z=2^{11-t}\in\{32,16,8\}$.  The contribution
of $S$ is at most $z^2 2^t$, whereas those of
$U,V$ are at most $2^{a_j},2^{b_j}$.  Therefore
\[
 2^{a_j}+2^{b_j}
 \ge4096z-z^2 2^t=2^{22-t}.                    \tag{SH-46.6}
\]
Since $a_j,b_j\ge t$, this forces
$a_j+b_j\ge22$.  Equality is possible only for
the pair $(t,22-t)$ or its reverse: if neither
ceiling reaches $22-t$, both must be $21-t$,
whose sum is greater than $22$ in all three cases.
Combining the two lower bounds with (SH-46.4)
and $u,v\ge12$ forces $u=v=12$ and equality
for each pair $(a_j,b_j)$.

If $a_j=t$, the second line of (SH-46.5) would
require $4096+2^{t+1}\ge5601$, which is false
for $t=6,7,8$.  If $b_j=t$, its third line
would require
\[
 4096+2^t+2^{16-t}\ge5601.
\]
The three left sides are $5184,4736,4608$,
respectively, again impossible.  This eliminates
the last three branches and proves the theorem.
\end{proof}

The exact reduction is generated by
\path{programs/sh_three_term_pattern_propagation.py},
with \path{inputs/sh_three_term_pattern_all.json}
and \path{outputs/sh_three_term_pattern_all.json}.
Every branch, initial pattern, integer update, and
final contradiction is preserved.  The separate
root program
\path{programs/rt_verify_three_term_propagation.py}
independently reconstructs all $2970$ branches,
replays every semantic deduction, and checks the
three final arithmetic cases.  Its complete result
is \path{outputs/rt_three_term_propagation_verification.json}.
The earlier AND linear and integer searches were
exploratory and are not premises of SH-46.

\begin{theorem}[SH-47: width six in dimensions eight and nine]
\label{thm:width-six-dimensions-eight-nine}
Over every field there is no strict shortening with
$m=6,A=2$ and $D\in\{8,9\}$, for arbitrary finite
network layouts.
\end{theorem}
\begin{proof}
Use the exact rank-profile and marked-minimum
certificate scheme of SH-42, with the additional
necessary rectangular inequalities of SH-44.
The complete coverage is
\[
\begin{array}{c|c|c|c|c}
D&\text{first minima}&c_{\rm rad}&c_{\rm two}&c_{\rm prune}\\\hline
8&0,\ldots,31&21&224&601\\
9&0,\ldots,35&26&177&598
\end{array}
\]
The three certificate-count columns have the
same meanings as in SH-42.  The radial exceptions
are $20,\ldots,30$ in dimension eight and
$25,\ldots,34$ in dimension nine.  Every
admissible choice of the other minimum and its
folded distance is either excluded by a recorded
integer rank inequality or covered by an exact
nonnegative rational certificate with full
variable cancellation and right side $-1$.

The additional Singh row is used only when
its number-of-slices hypothesis holds.  A recorded
violation, for example
\[
 49(2^a+2^b)<67\,2^{6(D-h)},\qquad 2^{6h}\ge7,
\]
gives the same threshold exclusion (SH-42.2).
The independent checker verifies these two
integer conditions directly, as well as every
earlier row, orbit weight, pruning decision,
and complete branch coverage.  Both dimensions
passed that separate reconstruction.  Thus
no numerical solver status is a premise.
\end{proof}

The new generators are
\path{programs/sh_singh_radial_rank_certificate.py}
and \path{programs/sh_singh_two_minimum_certificate.py}.
The complete manifests are
\path{outputs/sh_singh_two_min_D8_m6_pruned_manifest.json}
and \path{outputs/sh_singh_two_min_D9_m6_pruned_manifest.json};
their radial outputs have prefixes
\path{outputs/sh_singh_m6_D8_min} and
\path{outputs/sh_singh_m6_D9_min}.
The independent semantic checker and coverage audit
are recorded by the input/output pairs
\path{inputs/rt_singh_D8_verification.json},
\path{outputs/rt_singh_D8_verification.json},
\path{inputs/rt_singh_D9_verification.json}, and
\path{outputs/rt_singh_D9_verification.json}.

\begin{lemma}[SH-48: a dimension-four three-summand reduction]
\label{lem:dimension-four-three-summand-minimum}
For $m=6,D=4,A=3$, a hypothetical strict shortening
has an absolute sign-cut rank-ceiling minimum at
most $13$.
\end{lemma}
\begin{proof}
There are eight cuts, each ceiling is at most $24$,
and the frame/cost bound gives each term a ceiling
sum at most $141$.  For example,
$\log_2 3>19/12$ gives the strict upper bound
$32(6-\log_2 3)<141+1/3$ before integer rounding.
Give ceilings $22,23,24$ weights $1,2,4$,
respectively, and smaller ceilings weight zero.
Each cut requires total weight at least $4$.

If all ceilings were at least $14$, each row
would have only $141-8\cdot14=29$ units above
that baseline.  Three $24$ entries already use
$30$ units.  Two $24$ entries leave nine units,
allowing at most one further high entry and
total weight at most $10$.  One $24$ leaves
$19$ units, allowing at most two entries at
least $22$, and total weight at most $8$.
With no $24$, at most three entries are at
least $22$, for total weight at most $6$.
Thus three rows supply weight at most $30$,
less than the eight cuts' demand $32$.
\end{proof}

SH-48 is a necessary branch reduction, not an
exclusion in dimension four.  A bounded exact
pattern-tree exploration is preserved in
\path{programs/sh_three_term_dimension_four_tree.py}
with input/output
\path{inputs/sh_three_term_D4_pattern_tree.json} and
\path{outputs/sh_three_term_D4_pattern_tree.json}.
It completed the $t=13$ exploration but stopped
at its explicit node limit during $t=12$, with
surviving interval boxes.  These boxes are not
tensor constructions.  No full dimension-four
exclusion is inferred from that incomplete search
or from subsequent numerical searches of individual
boxes.

\subsection{Operational limits of the explicit local-unitary pair}

\begin{proposition}[SH-49: rank-one binary weights give a cut metric]
\label{prop:rank-one-binary-l1}
Let $P_v$ be binary matrices of rank at most one, of any fixed
rectangular size.  Then $d(v,w)=\operatorname{rank}(P_v-P_w)$
is a nonnegative combination of cut pseudometrics.  In particular,
the weights in Proposition~\ref{rt:lu-rank-failure} cannot certify
a one-term graph-tensor resource advantage, even if the conversion
uses arbitrary local maps before and after its non-Clifford local
unitary.
\end{proposition}
\begin{proof}
Every nonzero binary rank-one matrix has a unique factorization
$P=ab^{\mathsf T}$ with $a,b$ nonzero: its column and row lines
each contain exactly one nonzero vector.  For two such matrices,
\[
 \operatorname{rank}(ab^{\mathsf T}-cd^{\mathsf T})
   =\boldsymbol 1_{a\ne c}+\boldsymbol 1_{b\ne d}.       \tag{SH-49.1}
\]
If both factors differ, $a,c$ and $b,d$ are two independent
pairs over $\mathbb F_2$, so the matrix has rank two.  If exactly
one factor differs it has rank one, and if neither differs it
is zero.  The distance from a nonzero matrix to zero is one.
Thus
\[
 0\longmapsto0,\qquad
 ab^{\mathsf T}\longmapsto\tfrac12(e_a,e_b)
                                                               \tag{SH-49.2}
\]
is an isometry into a finite-dimensional $\ell_1$ space.  Equivalently,
give weight $1/2$ to each set of vertices having a specified
nonzero left factor, and to each set having a specified nonzero
right factor; zero matrices lie in none of these sets.

For a graph tensor with binary wire widths $b_{uv}$, the logarithm
of its flattening rank across a cut $C$ is
$\sum_{uv}b_{uv}\boldsymbol 1_{|\{u,v\}\cap C|=1}$.
A restriction cannot increase this rank.  Summing those
inequalities with the nonnegative coefficients in (SH-49.2)
proves
\[
 \sum_{uv\in H}b_{uv}d(u,v)
 \le \sum_{uv\in G}b_{uv}d(u,v)
 \quad\text{whenever }T_H\le T_G.                      \tag{SH-49.3}
\]
The composite of preparation, local unitary, and distillation
is still a restriction, so it obeys the same inequality.
This proof concerns one restriction, not a sum of restrictions.
\end{proof}

For the displayed $27$ weights, the zero sites are $3,25,26$.
The exact certificate contains $30$ cuts of weight $1/2$ and
checks (SH-49.1)--(SH-49.2) for all $351$ unordered site pairs.
The weighted stabilizer-rank increase $36\to37$ remains valid;
it is not a graph-resource inequality for these weights.

\begin{lemma}[SH-50: grouping without entanglement loss]
\label{lem:lu-graph-grouping}
Fix any grouping of the qubits into parties.  If two locally
unitary equivalent pure states are, after party-local isometries
and the removal of local pure factors, graph tensors with binary
edge multiplicities $b_{ij}$ and $c_{ij}$, then $b_{ij}=c_{ij}$
for every pair of parties.  Consequently a grouping that turns
both states in the explicit pair into graph tensors without
postselection or loss cannot produce a graph-resource advantage.
\end{lemma}
\begin{proof}
Normalize each Bell pair to a unit vector and let $E(C)$ be
the base-two entropy across the party cut $C$.  For a graph tensor,
$E(C)$ is its total crossing edge multiplicity.  In particular,
\[
 2b_{ij}=E(\{i\})+E(\{j\})-E(\{i,j\}).              \tag{SH-50.1}
\]
Party-local isometries, local pure factors, and local unitaries
preserve all these entropies.  Formula (SH-50.1) therefore
recovers the same multiplicities from either state.  The same
argument can instead use logarithmic flattening ranks, since
every Bell-pair Schmidt spectrum is uniform.
\end{proof}

\paragraph{Exact systematic presentation for an operational search.}
The source data of Proposition~\ref{rt:lu-rank-failure} admit
pivot sites $P=\{1,2,3,4,8,16\}$.  Row reduction gives
the same support subspace in the form $[I_6\ M]$, with the
columns placed at their original physical sites.  Applying
Hadamards on the remaining $21$ qubits sends the source to
the graph state having one edge from pivot $i$ to nonpivot $v$
exactly when $M_{iv}=1$.  There are $70$ such edges.  The target
has the same edges and the $11$ pivot-pivot edges already listed
in Proposition~\ref{rt:lu-rank-failure}.

To check this description, write a support word as
$(y,yM)$, $y\in\mathbb F_2^6$.  Fourier transformation of the
nonpivot coordinates gives the amplitude
\[
 2^{-27/2}(-1)^{yMz^{\mathsf T}}
 \quad\text{at }(y,z).
\]
The target multiplies this by $(-1)^{q(y,yM)}$.
All edges defining $q$ join pivot sites, so its polynomial is
exactly the listed quadratic form on $y$, with no linear
correction.  This verifies the graph-state representation
without expanding $2^{27}$ amplitudes.  It does not identify
either graph state with the graph tensor on its graph: the
former has one qubit at a site, while the latter has one qubit
for every incident wire.

The calculations and complete graph data are preserved in
\path{programs/sh_lu_operational_reduction.py},
\path{inputs/sh_lu_operational_reduction.json}, and
\path{outputs/sh_lu_operational_reduction.json}.
They provide a finite preparation/distillation search model,
not an EPR-resource conversion or a shortening.

\begin{proposition}[SH-51: the standard preparation is too costly]
\label{prop:lu-standard-preparation-obstruction}
Let $G_0$ be the $70$-edge preparation graph just specified.
Group its $27$ sites in any fixed way into parties $I$, and
let $n_i$ be the number of sites in party $i$.  For every
pseudometric $d$ on $I$, write $L_d(G_0)$ for the sum of its
edge lengths, counting the original edges with multiplicity.

Suppose a graph tensor $T_H$ is covered by $A$ tensors,
each of which is a restriction of a tensor having local
dimension at most $2^{Nn_i}$ at party $i$.  This includes
$N$ copies of either explicit LU state, followed by arbitrary
party-local operations.  Put
\[
 \pi_i=\frac1{27}\sum_{v=1}^{27}d(i,\gamma(v)),
                                                        \tag{SH-51.1}
\]
where $\gamma(v)$ is the party containing site $v$.
If the width of an edge of $H$ means the base-two logarithm
of its dimension, then
\[
 L_d(H)\le N L_d(G_0)+(\log_2 A)\sum_{i\in I}\pi_i.
                                                        \tag{SH-51.2}
\]
In particular, this preparation template cannot give a
one-term cost advantage for any pseudometric, or a positive
linear cost advantage using $2^{o(N)}$ terms as $N\to\infty$.
\end{proposition}
\begin{proof}
The exact rational certificate described below routes demand
$2/27$ between every pair of distinct vertices of $G_0$,
with total path load at most one on every edge.
Triangle inequality along each path therefore gives
\[
 \frac2{27}\sum_{1\le u<v\le27}d(\gamma(u),\gamma(v))
       \le L_d(G_0).                                  \tag{SH-51.3}
\]
The proof remains valid for coincident parties and zero
distances.  Averaging the triangle inequality through all
$27$ sites also gives $d(i,j)\le\pi_i+\pi_j$.

Let $b_i$ be the total width incident to party $i$ in $H$.
Its singleton flattening rank is exactly $2^{b_i}$.
Each of the $A$ covering tensors has singleton rank at most
$2^{Nn_i}$, by the stated local-dimension bound and monotonicity
under restriction.  Rank subadditivity implies
$b_i\le Nn_i+\log_2 A$.  Consequently
\begin{align*}
 L_d(H)&\le\sum_i b_i\pi_i\\
       &\le N\sum_i n_i\pi_i+(\log_2 A)\sum_i\pi_i\\
       &=\frac{2N}{27}\sum_{u<v}d(\gamma(u),\gamma(v))
                     +(\log_2 A)\sum_i\pi_i,
\end{align*}
which proves (SH-51.2) using (SH-51.3).
The party set, grouping, and metric are fixed in the
asymptotic assertion, so the last sum is a finite constant.
\end{proof}

Here is the complete finite certificate for (SH-51.3).
The program \path{programs/sh_lu_preparation_routing.py}
enumerated $15152$ simple paths of length at most four and
searched for uniform all-pairs flow.  The numerical optimum
$1/13$ was used only to obtain slack: the proof uses $414$
explicit paths with positive rational weights, whose
exact demand on each of the $351$ pairs is $2/27$ and whose
load on each of the $70$ edges is at most one.  The paths
and rational coefficients are in
\path{outputs/sh_lu_preparation_routing.json}; the input is
\path{inputs/sh_lu_preparation_routing.json}.

The separate program
\path{programs/sh_verify_lu_preparation_routing.py}
imports no optimization or discovery code.  It enumerates
the $64$ original support words to reconstruct the six
systematic rows, reconstructs the $70$ graph edges, checks
each path, and sums every demand and capacity using exact
fractions.  Its output is
\path{outputs/sh_lu_preparation_routing_verified.json}.
Thus the LP status and its floating-point optimum are not
premises.  An initial checker assumption that the supplied
rows were already systematic was corrected by the support
enumeration; the preserved diagnostic is
\path{outputs/sh_lu_preparation_verifier_initial_diagnostic.json}.

SH-51 concerns this specified preparation graph, even with
arbitrary subsequent local maps and grouping.  It does not
exclude a cheaper preparation of the same state, a different
resource state, or a sum of different preparation layouts.

\paragraph{A smaller checked preparation is also obstructed.}
Local complementation at sites $15,1,15$, in that order,
gives a $66$-edge graph state locally Clifford equivalent
to the source.  For the graph-state generators $X_iZ_{N(i)}$,
complementation at $v$ is realized, up to local Pauli signs,
by the binary symplectic map $(x_v,z_v)\mapsto(x_v+z_v,z_v)$
at $v$ and $(x_u,z_u)\mapsto(x_u,z_u+x_u)$ at each neighbor.
Adding the transformed generator at $v$ to those of its
neighbors restores an identity $X$ block and toggles exactly
the edges among the neighbors.  This verifies the stated
local-Clifford equivalence directly.

This $66$-edge graph also has an exact uniform $2/27$
all-pairs routing, now on $410$ paths of length at most four.
Consequently the entire conclusion of SH-51 holds with this
smaller preparation graph as well.  Its exact path data and
separate verification are
\path{outputs/sh_lu_preparation_lc_routing.json} and
\path{outputs/sh_lu_preparation_lc_routing_verified.json};
the same verifier replays the complementation sequence before
checking all demands and capacities.

The bounded search that found it is preserved in
\path{programs/sh_lu_preparation_lc_search.py} with its
same-named input/output files.  It tried $90000$ single
complementation moves and $96000$ edge-pivot moves.  The
smallest edge count encountered was $66$; this is not a
certificate that no smaller preparation exists.

\begin{proposition}[SH-52: spin-sector labels alone do not lift Cauchy]
\label{prop:cauchy-central-spin-obstruction}
Work over a field of characteristic zero or greater than three.
At each party of the binary $r=4$ crown construction, let
$S=\frac16\sum_{\pi\in S_3}P_\pi$ be the projector onto
$\operatorname{Sym}^3(\mathbb F^2)$ and let $M=I-S$.
There is no cover identity for $H_4$ by crown tensors
$J_\sigma$, $\sigma\in S_4$, using any number of terms whose
local maps all lie in $\operatorname{span}\{S,M\}$.
Thus a correction that records only the two local spin
sectors, with scalar action inside each multiplicity space,
already fails at $r=4$.
\end{proposition}
\begin{proof}
Use the crown port and alternating-edge conventions of
Proposition~\ref{pg:cauchy-r3-integral}, with four labels
instead of three.  A local binary basis index is
$b_0+2b_1+4b_2$.  Define a two-row map $R_1$ by the covectors
\[
 u_0=\langle1|-\langle2|,
 \qquad u_1=\langle2|-\langle4|,
\]
and define $R_2$ by replacing every basis index $a$ by $7-a$.
Each row has coefficient sum zero within its Hamming-weight
class.  Therefore $R_jS=0$ and $R_jM=R_j$.
Apply $R_1,R_1,R_2,R_2$ in that order to the left parties,
and the same four maps to the right parties.  Write $h$ for
the projected target and $c_\sigma$ for the $24$ projected
crowns, all in $(\mathbb F^2)^{\otimes8}$.

Any local map $aS+bM$ becomes multiplication by $b$ after
this projection.  Consequently a proposed identity with
arbitrarily many such terms would imply
$h\in\operatorname{span}\{c_\sigma:\sigma\in S_4\}$.
The following integer functional refutes that implication.
Index the projected coordinates by
$z=z_0+2z_1+\cdots+128z_7$ and use the listed nonzero
coefficients; all others are zero:
\[
\begin{array}{c|rrrrrrrrrr}
z&0&1&5&7&13&16&17&19&20&21\\\hline
w_z&3&-4&-4&10&-4&-8&-4&-2&-4&-4\\[2pt]
z&22&25&28&65&69&76&77&85&86&88\\\hline
w_z&4&4&4&-4&-8&16&-8&4&4&-4
\end{array}
\]
Exact coefficient evaluation gives
\[
 \sum_z w_z(c_\sigma)_z=0\quad(\sigma\in S_4),
 \qquad \sum_z w_zh_z=16.                       \tag{SH-52.1}
\]
These evaluations are over the integers, and $16$ is
nonzero in the stated characteristics.  This is the desired
contradiction.

For an explicit verification of (SH-52.1), expand each of
the eight selected local covectors into its two signed
basis covectors.  For each of the $256$ boundary choices,
evaluate a graph coefficient as the product of its twelve
$\epsilon$ entries, where
$\epsilon_{01}=1$, $\epsilon_{10}=-1$, and the other
entries vanish.  Multiply by the eight local signs and
then by $w_z$, and sum.  This rule defines every coefficient
using only the displayed finite data.
\end{proof}

The discovery program
\path{programs/sh_cauchy_spin_sector_probe.py}, with
\path{inputs/sh_cauchy_spin_sector_probe.json}, records the
full $256$-coordinate arrays and rational ranks $22$ and
$23$ before and after adjoining $h$, in
\path{outputs/sh_cauchy_spin_sector_probe.json}.
The separate coefficient checker
\path{programs/sh_verify_cauchy_spin_sector.py} does not
read those arrays or import symbolic linear algebra.  It
evaluates boundary coefficients by the rule above and
records the $24$ zeros and target value $16$ in
\path{outputs/sh_cauchy_spin_sector_verified.json}.

This integer witness vanishes on the target modulo two,
and the tested projected target is in the crown span over
$\mathbb F_2$.  No characteristic-two obstruction is inferred
from this projection.  More generally, SH-52 does not
exclude local maps that act nontrivially on the
two-dimensional multiplicity spaces, nor does it exclude
multiple dense noncentral terms within a matching class.

\begin{corollary}[SH-52b: the full mixed-sector obstruction over $\mathbb F_2$]
\label{prop:cauchy-central-spin-f2}
The same central-sector-only cover is impossible over
$\mathbb F_2$ when the symmetric projector is defined by
$S=I+C+C^2$, where $C$ cyclically permutes the three ports,
and the complementary projector is $M=I-S$.
\end{corollary}
\begin{proof}
The cyclic permutation has order three, which is invertible
in $\mathbb F_2$.  Its fixed subspace is exactly the symmetric
subspace: on each Hamming-weight class of three binary ports,
the cyclic group is transitive.  Hence $S$ is its rank-four
projector and $M$ is the complementary rank-four projector.
The class sum $C+C^2$ commutes with every port permutation.

At every party now retain all four covectors, in the order
$(u_0,u_1,\bar u_0,\bar u_1)$, where the bars complement
the three basis bits.  Call this map $R$.  As before,
$RS=0$ and $RM=R$, so every allowed local map acts as a
scalar after projection.  Index a coordinate in
$(\mathbb F_2^4)^{\otimes8}$ by
$z=\sum_{v=0}^7z_v4^v$.  Sum, over $\mathbb F_2$, the
coordinates in the following set:
\[
\begin{split}
\mathcal Z=\{&442,443,491,494,506,555,559,619,\\
             &635,639,651,667,675,827,875,6939\}.
\end{split}
\]
This functional is zero on all $24$ projected crowns and
one on the projected target.  The boundary evaluation rule
in SH-52 verifies these $25$ values directly, now with all
signs reduced modulo two.  It therefore contradicts any
proposed linear combination of projected crowns, and hence
any central-sector-only cover.
\end{proof}

The full-sector witness is extracted by
\path{programs/sh_cauchy_spin_sector_f2.py}, with the
same-named input/output JSON files.  Its exact integer
source arrays are
\path{inputs/sh_cauchy_all_mixed_sector_arrays.npz}.
The independent checker
\path{programs/sh_verify_cauchy_spin_sector_f2.py} reads
no projected arrays and performs no elimination: it checks
the $16$ listed boundary-coordinate evaluations directly.
Its output is
\path{outputs/sh_cauchy_spin_sector_f2_verified.json}.
The difference from the earlier characteristic-two test is
material: this functional retains every local spin weight,
whereas the first test retained just one weight at each site.

Noncentral multiplicity maps remain available.  For example,
on either retained weight slice, the two port transpositions
act by
\[
 A=\begin{pmatrix}-1&0\\1&1\end{pmatrix},\qquad
 B=\begin{pmatrix}1&1\\0&-1\end{pmatrix}.
\]
The matrices $I,A,B,AB$ span all two-by-two matrices over
every field of characteristic different from three.
Indeed, for $\bigl(\begin{smallmatrix}a&b\\c&d\end{smallmatrix}\bigr)$,
put $t=(a-d-2b+2c)/3$, $x=a-b+c-t$, $y=c-t$, $z=b+t$;
then it equals $xI+yA+zB+tAB$.
Thus the obstruction does not come from an absence of
local noncentral operators.  It concerns assembling them
into a full cover with a sufficiently small number of terms.

\paragraph{A bounded noncentral construction attempt.}
To test the remaining freedom directly, an exploratory
real-valued fit retained the entire all-mixed sector,
including all compatible spin weights.  It used two terms
for each of the $24$ crown classes and arbitrary two-by-two
multiplicity maps at every party, acting identically on
the two spin weights.  After $161$ iterations and an
$80$-second optimization limit, the best relative squared
coefficient residual was $0.01401756627140251$.
This is not an identity, a border construction, or a
nonexistence certificate.  The other spin sectors were
not part of this fit.

The explicit reverse-mode derivatives passed two
finite-difference checks with absolute errors below
$1.1\cdot10^{-11}$.  The program, configuration, trace,
and best numerical maps are
\path{programs/sh_cauchy_noncentral_fit.py},
\path{inputs/sh_cauchy_noncentral_fit.json},
\path{outputs/sh_cauchy_noncentral_fit.json}, and
\path{outputs/sh_cauchy_noncentral_fit_best.npz}.
The exact integer projected source arrays are retained
separately in the input NPZ file cited above.  A missing
optional runtime package and an initial array-contraction
syntax error were corrected before optimization; both
diagnostics are preserved in the output directory.

\begin{theorem}[SH-53: width six in dimension ten]
\label{thm:width-six-dimension-ten}
Over every field, no strict shortening with $m=6$, $A=2$,
and $D=10$ exists, even with arbitrary finite network layouts.
\end{theorem}
\begin{proof}
Apply the necessary rank-profile inequalities and the
marked-minimum averaging construction of SH-42, including
the rectangular inequalities of SH-44.  The first global
minimum lies in $0,\ldots,39$.  Exact radial certificates
exclude $0,\ldots,24$ and $39$.  For each of the remaining
fourteen values $25,\ldots,38$, mark also the minimum of
the second profile and its folded-cube distance from the
first mark.  Complete enumeration gives $405$ admissible
two-minimum cases and $1149$ cases excluded directly by
integer rank inequalities.  Every admissible case has an
exact nonnegative rational combination of necessary rows
whose variable coefficients cancel and whose right side
is $-1$.  Thus each branch implies $0\leq-1$.

The separate semantic verifier reconstructs the folded-cube
orbits, row meanings, rectangular witnesses, and all
pruning decisions without importing the certificate
generator.  It checks rational coefficients with exact
fractions, checks each certificate's hash, and verifies
coverage of every marked-minimum branch.  Its final
output records $26$ radial certificates, $405$
two-minimum certificates, $1149$ prunes, and complete
coverage.  No numerical optimization status is used
as a mathematical premise.
\end{proof}

The authoritative input manifests are
\path{inputs/rt_singh_D10_radial_manifest.json} and
\path{outputs/rt_D10_two_minimum_manifest.json}.
The independent audit is
\path{programs/rt_verify_unbalanced_certificates.py},
with input \path{inputs/rt_singh_D10_verification.json}
and output \path{outputs/rt_singh_D10_verification.json}.
The latter records every selected certificate path and
SHA-256 hash.  Earlier incomplete and timed-out outputs
are retained but are not premises of this theorem.
Together with the preceding exclusions, the remaining
width-six parameter ranges are
\[
 A=2:\quad 11\leq D\leq19,
 \qquad
 A=3:\quad 4\leq D\leq16.
\]
These are remaining necessary possibilities, not
constructions.

\begin{lemma}[SH-54: an invariant projection with polynomial cover cost]
\label{lem:polynomial-spin-zero-projection}
Work over $\mathbb C$.  Let $P_d$ be the orthogonal
projector onto the vectors fixed by the diagonal
$\mathrm{SU}(2)$ action on $(\mathbb C^2)^{\otimes d}$.
There are matrices $g_1,\ldots,g_K\in\mathrm{SU}(2)$ and
nonnegative real numbers $a_1,\ldots,a_K$ with
\[
 P_d=\sum_{j=1}^K a_j g_j^{\otimes d},\qquad
 \sum_j a_j=1,\qquad
 K\leq {d+3\choose3}+1.
 \tag{SH-54.1}
\]
For any grouping of the $d$ binary ports into parties,
each summand in (SH-54.1) is a product of party-local
maps.  Thus applying $P_d$ to an $A$-term cover increases
its number of terms by at most the displayed factor $K$.
\end{lemma}
\begin{proof}
Write a unit quaternion as $x=(a,b,c,e)\in S^3$ and set
\[
 g(x)=
 \begin{pmatrix}a+bi&c+ei\\-c+ei&a-bi\end{pmatrix}.
\]
Every matrix entry of $g(x)^{\otimes d}$ is a complex
linear combination of the ${d+3\choose3}$ real monomials
of total degree $d$ in $a,b,c,e$.  Let $v(x)$ be the
vector of these monomials and let $\bar v$ be its
average over the uniform probability measure on $S^3$.
The vector $\bar v$ lies in the convex hull of the
compact set $v(S^3)$.  In a real space of dimension $k$,
any finite convex combination of more than $k+1$ points
has an affine dependence; subtracting a suitable multiple
of that dependence removes one positive coefficient
without changing the combination.  Consequently the
convex hull of a compact set consists of combinations
of at most $k+1$ points and is compact.  This also
justifies passage from finite average approximations
to the integral.  Applying the argument with
$k={d+3\choose3}$ gives
$\bar v=\sum_{j=1}^K a_jv(x_j)$ and hence the same
identity for every matrix entry.

Quaternion multiplication by a fixed unit quaternion
is an orthogonal transformation of $\mathbb R^4$.
Therefore the spherical average of $g(x)^{\otimes d}$
is invariant under left multiplication by every
$g\in\mathrm{SU}(2)$; it fixes every invariant vector.
The average is self-adjoint because inversion preserves
the spherical measure.  It is therefore exactly $P_d$.
Finally $g_j^{\otimes d}$ factors over any grouping of
ports.  The scalar $a_j$ can be absorbed into a map at
one party, proving the cover statement.
\end{proof}

To apply this lemma to the balanced graph-cover search,
first gauge every binary EPR edge to the antisymmetric
tensor $\epsilon=e_0\otimes e_1-e_1\otimes e_0$.
This is an invertible local change of basis, performed
separately on the corresponding wire factors.  The
target $H_r=K_r(2)\sqcup K_r(2)$ is then fixed by the
common $\mathrm{SU}(2)$ action on all left ports.
There are $d=r(r-1)$ such ports.  If a proposed sum of
restrictions $\sum_{j=1}^A T_j$ satisfies only
\[
 P_{\rm left}\sum_{j=1}^A T_j=H_r,
 \tag{SH-54.2}
\]
then SH-54 turns it into an actual cover with at most
$A\bigl({r(r-1)+3\choose3}+1\bigr)$ terms.  Its
additional base-two logarithmic cost is
$6\log_2r+O(1)$, so an $o(r^2)$ cover exponent remains
$o(r^2)$.  For edge dimension $q=2^N$, the same binary
factorization uses $d=Nr(r-1)$ and has factor
$O(N^3r^6)$.  For fixed $r$, its additional
$\log_q$ cover exponent tends to zero as $N\to\infty$.
Equation (SH-54.2) is still a substantial identity to
prove; the lemma does not supply it.  No
characteristic-two version of this averaging argument
is asserted.

\paragraph{An explicit small projector and the numerical scope.}
For four effective spin-$\tfrac12$ factors, the
$24$ unit quaternions
\[
 \{\pm1,\pm\mathbf i,\pm\mathbf j,\pm\mathbf k\}
 \ \cup\
 \{(\pm1\pm\mathbf i\pm\mathbf j\pm\mathbf k)/2\}
\]
give the exact projector by averaging their fourth
tensor powers.  The signs in the second set vary
independently.  The checker
\path{programs/sh_verify_spinzero_projection.py}
verifies all $256$ matrix entries over $\mathbb Q(i)$,
the $768$ local intertwiner entries for the retained
mixed-sector covectors, and the integer projected
target and $24$ crown arrays.  In particular, projection
to left spin zero also projects these crown tensors
to right spin zero.  Its input/output pair is
\path{inputs/sh_spinzero_projection.json} and
\path{outputs/sh_spinzero_projection_verified.json}.
The denominators $2$ and $3$ matter; this finite
formula is not used over $\mathbb F_2$.

The deterministic coefficient least-squares refinement
of the earlier $48$-term fit reached relative squared
residual $0.0009475329570313699$ in the unprojected
full mixed sector.  Adding the checked spin-zero
projector and refining for another bounded
$180$ seconds reached
$6.245644609848444\cdot10^{-7}$, measured again in
all original tensor coordinates after that projector.
Neither value is zero.  These runs establish neither
an identity nor a border degeneration, and do not
include the other local Schur sectors.  The shared
program is
\path{programs/sh_cauchy_noncentral_varproj.py};
the two configurations and traces have stems
\path{sh_cauchy_noncentral_varproj} and
\path{sh_cauchy_spinzero_varproj} in the input and
output directories.  Best numerical maps are
preserved in their corresponding \path{_best.npz}
files.  The exact projector identity is independent
of these numerical fits.

\begin{lemma}[SH-55: a second-moment geometric rank bound]
\label{lem:sign-rank-second-moment}
Let a finite network of Euclidean cost $w$ have terminal
parties at $\{\pm e_i:1\leq i\leq D\}$.  All wire
widths are nonnegative integers.  For a uniformly random
$s\in\{\pm1\}^D$, take the cut choosing the terminal
$s_i e_i$ from each antipodal pair.  If $T$ is the
network tensor, put
\[
 \kappa(s)=\left\lceil\log_2\max\{1,
              \operatorname{rank}_s T\}\right\rceil.
\]
Then, over every field,
\[
                 \mathbb E\kappa(s)^2\leq w^2/4.
 \tag{SH-55.1}
\]
If each terminal output has dimension $2^m$, then
\[
 \kappa(-s)=\kappa(s),\qquad
 \operatorname{Var}\kappa\leq m^2D/2.
 \tag{SH-55.2}
\]
\end{lemma}
\begin{proof}
Let $k_s$ be the integer minimum cut capacity separating
the two terminal sets.  A network flattening factors
through the wires of any such cut, so
$\kappa(s)\leq k_s$.  For every $t\in(-1,1)$, the
hyperplane $s\cdot x=t$ induces precisely this terminal
partition.  Integrating the capacities crossed by these
hyperplanes gives
\[
 2\kappa(s)\leq2k_s
 \leq\sum_e b_e\,|s\cdot\Delta_e|,
\]
where $\Delta_e$ is the displacement vector of edge $e$.
Its portion of the projection interval inside $[-1,1]$
can only be shorter than $|s\cdot\Delta_e|$.  Since
$\mathbb E(s\cdot\Delta_e)^2=\|\Delta_e\|_2^2$,
the triangle inequality in $L^2(\{\pm1\}^D)$ yields
\[
 \|\kappa\|_{L^2}
 \leq\frac12\left\|\sum_e b_e|s\cdot\Delta_e|\right\|_{L^2}
 \leq\frac12\sum_e b_e\|\Delta_e\|_2=w/2.
\]
Zero-length edges cause no difficulty in this argument.

Complementary flattenings have equal rank, proving
evenness.  Flipping one coordinate moves two terminal
factors, each of dimension $2^m$, across the flattening.
Its rank therefore changes by a factor of at most
$2^{2m}$, so the ceiling changes by at most $2m$.
In the Fourier expansion on the sign cube, every
nonconstant coefficient of an even function has degree
at least two.  Consequently
\[
 \operatorname{Var}\kappa
 \leq\frac18\sum_{i=1}^D
       \mathbb E\bigl(\kappa(s)-\kappa(s^{(i)})\bigr)^2
 \leq\frac18D(2m)^2,
\]
where $s^{(i)}$ flips coordinate $i$.  This proves
(SH-55.2).  A zero tensor has the constant ceiling zero
under the stated convention.
\end{proof}

\begin{theorem}[SH-55b: bulk width-six exclusions]
\label{thm:width-six-moment-bulk}
For arbitrary finite layouts and over every field, no
strict shortening exists with either
\[
 m=6,\ A=2,\ D\geq17,
 \qquad\text{or}\qquad
 m=6,\ A=3,\ D\geq13.
\]
\end{theorem}
\begin{proof}
Write $\mu_j=\mathbb E\kappa_j$ and
$\sigma_j^2=\operatorname{Var}\kappa_j$ for the
summand profiles.  Every sign cut has target rank
$2^{6D}$.  For either two or three summands, at least
one summand on that cut has ceiling at least
$T=6D-1$: otherwise even three ranks sum to at most
$3\cdot2^{6D-2}<2^{6D}$.  Hence some profile has
$\Pr\{\kappa_j\geq T\}\geq1/A$.

The one-sided variance inequality gives
\[
 T\leq\mu_j+\sqrt{A-1}\,\sigma_j.
 \tag{SH-55.3}
\]
Here $\mu_j<T$ follows from the cost bounds below;
if $\sigma_j=0$, the displayed positive tail is already
impossible.  For completeness, the variance inequality
follows by applying Markov's inequality to
$(\kappa_j-\mu_j+u)^2$ and minimizing over $u>0$.

For $A=2$, strict shortening gives $w_j<10D$.
Thus SH-55 implies
\[
 \mu_j^2+\sigma_j^2<25D^2,
 \qquad \sigma_j^2\leq18D.
\]
For $D\geq2$, the function
$\sqrt{25D^2-x^2}+x$ is increasing on
$0\leq x\leq\sqrt{18D}$.  Therefore
\[
 \mu_j+\sigma_j
 <\sqrt{25D^2-18D}+\sqrt{18D}
 <5D-\frac95+\sqrt{18D}.
\]
For $D\geq17$, one has
$\sqrt{18D}<D+4/5$: after squaring, the required
polynomial is $25D^2-410D+16$, whose value at $17$
is $271$ and which is increasing thereafter.
The last upper bound is therefore less than
$6D-1=T$, contradicting (SH-55.3).

For $A=3$, the integer inequality $3^{12}>2^{19}$
gives $\log_2 3>19/12$.  Hence
\[
 w_j<2D(6-\log_2 3)<\frac{53D}{6},
 \quad
 \mu_j^2+\sigma_j^2<\left(\frac{53D}{12}\right)^2,
 \quad \sigma_j^2\leq18D.
\]
The corresponding increasing-function bound is
\[
 \mu_j+\sqrt2\,\sigma_j
 <\sqrt{(53D/12)^2-18D}+6\sqrt D.
 \tag{SH-55.4}
\]
At $D=13$, the two radicals on the right satisfy
\[
 \sqrt{441025/144}<\frac{55342}{1000},
 \qquad
 \sqrt{468}<\frac{21634}{1000}.
\]
The exact squared cross-product differences are
$9\,122\,816$ and $29\,956$, respectively; the
upper bounds sum to $76.976<77=T$.
For $D\geq14$, put $c=53/12$ and use
$\sqrt{c^2D^2-18D}<cD-9/c$.  It suffices that
\[
 \frac{19D}{12}+\frac{55}{53}>6\sqrt D.
\]
The difference is positive at $D=14$; the difference
of its two squared sides there is
$3483145/101124>0$.  Its derivative is
$19/12-3/\sqrt D>0$ for $D\geq14$.
Thus (SH-55.4) is less than $T$ throughout this range,
again contradicting (SH-55.3).
\end{proof}

These bulk exclusions reduce the remaining width-six
ranges to
\[
       A=2:\ 11\leq D\leq16,
       \qquad A=3:\ 4\leq D\leq12.
\]
They also supply a stronger exact row for the two-term
rank-profile certificate model.  If
$z_{s,k}=\mathbf1[\kappa(s)\geq k]$ and
$M=2^{D-1}$ is the number of complementary sign cuts,
then at general width $m$ a strict two-term shortening
must satisfy, separately for each summand,
\[
 \sum_s\sum_{k=1}^{mD}(2k-1)z_{s,k}
 \leq D^2(m-1)^2M-1.
 \tag{SH-55.5}
\]
Indeed $\kappa^2=\sum_{k=1}^{\kappa}(2k-1)$,
and the strict second-moment bound has an integer left
side.  Averaging within folded-cube orbits preserves
this linear inequality; an orbit of size $w_a$
contributes coefficient $w_a(2k-1)$ to its averaged
threshold variable.  No equality between tensor ranks
and network cut capacities is required.

\begin{corollary}[SH-55c: the two-term bulk bound improves to dimension fifteen]
\label{cor:width-six-two-term-coupled-moment}
There is no strict shortening with $m=6$, $A=2$, and
$D\geq15$, over any field and for arbitrary finite layouts.
\end{corollary}
\begin{proof}
The terminal dimensions give $0\leq\kappa_j\leq n=6D$.
Put $Z_j=n-\kappa_j$.  At every sign cut, the rank sum
condition forces either one ceiling to equal $n$, or
both ceilings to equal $n-1$.  Indeed, if both are at
most $n-1$ and one is at most $n-2$, their rank upper
bounds sum to less than $2^n$.  Therefore
\[
                 Z_1Z_2\leq1\quad\text{pointwise}.
\]
Let $a_j=\mathbb EZ_j$ and $A=a_1+a_2$.
The strict moment bound implies
$\mathbb E\kappa_j<5D$, so $a_j>D$.
By the variance bound and Cauchy--Schwarz,
\[
 a_1a_2\leq1+\sqrt{\operatorname{Var}Z_1
                         \operatorname{Var}Z_2}
          \leq18D+1.
 \tag{SH-55.6}
\]
On the other hand,
\[
\begin{split}
 50D^2
 &>\mathbb E(\kappa_1^2+\kappa_2^2)\\
 &=72D^2-12DA+\mathbb E(Z_1^2+Z_2^2)\\
 &\geq72D^2-12DA+A^2-2.
\end{split}
\]
Hence $(A-6D)^2<14D^2+2$.  For $D\geq15$,
$\sqrt{14D^2+2}<15D/4$, since $D^2>32$.
Consequently $A>9D/4$, and
\[
 a_1a_2=D(A-D)+(a_1-D)(a_2-D)
       >5D^2/4>18D+1.
\]
The final inequality is equivalent to
$5D^2-72D-4>0$; its value at $D=15$ is $41$
and it increases thereafter.  This contradicts
(SH-55.6).
\end{proof}

\begin{theorem}[SH-56: moment certificates in dimensions eleven through thirteen]
\label{thm:width-six-moment-eleven-thirteen}
Over every field, no strict shortening with $m=6$,
$A=2$, and $D\in\{11,12,13\}$ exists, for arbitrary
finite network layouts.
\end{theorem}
\begin{proof}
Use the radial marked-minimum scheme of SH-42 and
SH-44, adding the necessary quadratic-moment row
(SH-55.5).  For dimension eleven, the previously
generated radial certificates exclude first minima
$0,\ldots,29$ and $43$; the new row gives exact
certificates for all thirteen remaining values
$30,\ldots,42$.  For dimensions twelve and thirteen,
the new radial model supplies certificates for all
$48$ and $52$ possible first minima, respectively.
Thus the complete radial coverage counts are
\[
\begin{array}{c|c|c|c}
 D&\text{first minima}&\text{old rows only}&
                   \text{moment model}\\\hline
 11&0,\ldots,43&31&13\\
 12&0,\ldots,47&0&48\\
 13&0,\ldots,51&0&52
\end{array}
\]
Every selected certificate is a nonnegative rational
combination with exact variable cancellation and a
negative right side.  The independent root verifier
reconstructs every row semantically, including the
new coefficient $w_a(2k-1)$ and right side
$D^2(m-1)^2M-1$, checks the exact fractions and file
hashes, and verifies complete coverage of the first
minimum.  All three dimensions passed this separate
audit.  Two-minimum search outputs and numerical
solver statuses are unnecessary for this proof.
\end{proof}

The new model is implemented separately in
\path{programs/sh_moment_radial_rank_certificate.py};
the earlier generators and their outputs are preserved.
The new complete-case manifests are
\path{outputs/sh_moment_D11_batch.json},
\path{outputs/sh_moment_D12_batch.json}, and
\path{outputs/sh_moment_D13_batch.json}.
Their selected old and new certificates, exact hashes,
and complete coverage are checked by
\path{programs/rt_verify_moment_certificates.py}
and independently recorded in
\path{outputs/rt_moment_D11_verification.json},
\path{outputs/rt_moment_D12_verification.json}, and
\path{outputs/rt_moment_D13_verification.json}.
Combined with SH-55b and SH-55c, the remaining
width-six ranges are now
\[
             A=2:\ D=14,
             \qquad A=3:\ 4\leq D\leq12.
\]
The dimension-fourteen radial tests leave three
uncertified first-minimum cases, $43,44,45$.
Those tests do not establish either feasibility or
an exclusion in that dimension.

\begin{theorem}[SH-55d: a joint edge bound closes the two-term case]
\label{thm:width-six-two-term-joint-edge}
Over every field, a strict shortening with $m=6$ and
$A=2$ does not exist in any dimension, for arbitrary
finite network layouts.
\end{theorem}
\begin{proof}
It suffices first to exclude every $D\geq13$.
Put $n=6D$, $Z_j=n-\kappa_j$, $a_j=\mathbb EZ_j$,
and $B=a_1+a_2$.  The same-cut rank condition, used in
SH-55c, says that every pair $(Z_1(s),Z_2(s))$ has one
of the forms
\[
                  (a,0),\qquad(0,b),\qquad(1,1),
                  \qquad 0\leq a,b\leq n.
 \tag{SH-55.7}
\]
In particular $Z_1Z_2\leq1$.  The strict geometric
moment bounds give $a_1,a_2>D$ and, by the calculation
in SH-55c,
\[
 (B-6D)^2<14D^2+2,
 \qquad B>\frac94D.
 \tag{SH-55.8}
\]
The second inequality follows already for $D\geq6$,
since $D^2>32$ implies
$\sqrt{14D^2+2}<15D/4$.

Write $T=Z_1+Z_2$ and $p=\Pr\{T\leq12\}$.
By (SH-55.7), $0\leq T\leq6D$, and therefore
\[
 B=\mathbb ET\leq12p+6D(1-p),
 \qquad
 p\leq\frac{6D-B}{6D-12}
       <\frac{5D}{8(D-2)}.
 \tag{SH-55.9}
\]
We now use the joint geometry of the two profiles.
Each $Z_j$ is even on the sign cube and changes by at
most $12$ along a cube edge.  For $s'=s^{(i)}$, the
allowed states (SH-55.7) imply
\[
 \bigl(Z_1(s)-Z_1(s')\bigr)^2+
 \bigl(Z_2(s)-Z_2(s')\bigr)^2
 \leq145+143\,\mathbf1_{
            \{T(s)\leq12,\ T(s')\leq12\}}.
 \tag{SH-55.10}
\]
Indeed, two states on the same coordinate axis have
combined squared difference at most $144$.  Between
an axis state and $(1,1)$ the bound is $144+1=145$.
Between opposite axes, the separate Lipschitz bounds
force both nonzero coordinates to be at most $12$,
so both endpoint totals are at most $12$ and the
combined squared difference is at most $288$.
These cases include zero coordinates and exhaust all
possibilities.

Apply the even-cube Poincare inequality separately to
$Z_1$ and $Z_2$, add the inequalities, and use
$\Pr\{T(s)\leq12,T(s^{(i)})\leq12\}\leq p$.
This gives
\[
 \operatorname{Var}Z_1+\operatorname{Var}Z_2
 \leq\frac D8(145+143p)
 <\frac D8\left(145+\frac{715D}{8(D-2)}\right).
 \tag{SH-55.11}
\]
Conversely, $\operatorname{Var}(Z_1+Z_2)\geq0$ and
$\mathbb EZ_1Z_2\leq1$ give
\[
\begin{split}
 \operatorname{Var}Z_1+\operatorname{Var}Z_2
 &\geq2a_1a_2-2\mathbb EZ_1Z_2\\
 &\geq2a_1a_2-2
 >2D(B-D)-2
 >\frac52D^2-2.
\end{split}
 \tag{SH-55.12}
\]
Here $a_1a_2=D(B-D)+(a_1-D)(a_2-D)>D(B-D)$.
The lower bound in (SH-55.12) exceeds the upper bound
in (SH-55.11) for every integer $D\geq13$.
After multiplication by the positive denominator
$64(D-2)$, the difference is
\[
             160D^3-2195D^2+2192D+256.
 \tag{SH-55.13}
\]
Its value at $D=13$ is $9317>0$.  For $D\geq14$,
the expression equals
$D^2(160D-2195)+2192D+256>0$.
This proves the contradiction for $D\geq13$.

The lower dimensions are already excluded: SH-38
covers $D=2$; SH-40 and SH-41 cover $D=3,4$;
SH-42 and SH-43 cover $D=5,6,7$; SH-47 covers
$D=8,9$; SH-53 covers $D=10$; and SH-56 covers
$D=11,12$.  If dimension one is allowed, each term
of cost $w$ has two-terminal rank at most
$2^{w/2}$ by the minimum-cut and path-length bound.
Two terms of cost strictly below $2(6-1)=10$
cannot sum to the target rank $64$.  This also
excludes $D=1$.
\end{proof}

The joint-edge proof is analytic and uses no solver
status or numerical profile.  It has been independently
checked by the root and exponent-audit agents, including
all transitions in (SH-55.10), the Poincare factor,
and the integer comparison (SH-55.13).  The earlier
dimension-fourteen radial diagnostics remain preserved;
the analytic argument now excludes that dimension.
Together with the other width-six exclusions, the
remaining parameter range at width six is precisely
\[
                        A=3,\qquad4\leq D\leq12.
\]
These are unexcluded necessary possibilities, not
constructions or assertions of existence.

\begin{theorem}[SH-57: three summands in dimension twelve]
\label{thm:width-six-three-term-dimension-twelve}
Over every field and for arbitrary finite layouts, no
strict shortening exists with $m=6$, $A=3$, and $D=12$.
\end{theorem}
\begin{proof}
We use an integer quadratic majorant for the exponential
rank bound.  If $a$ is an integer, $N-a=M\geq32$, and
$k\leq N$ is an integer, then
\[
       2^{k-N}\leq\frac{(k-a)^2}{M^2}+2^{-M}.
 \tag{SH-57.1}
\]
To prove this, put $t=k-a\leq M$.  At $t=0$ the
additive term is exact.  For $1\leq t\leq M$,
the function $2\log t+(M-t)\log2$ is concave, so
$t^2 2^{M-t}$ is at least the smaller of its endpoint
values $2^{M-1}$ and $M^2$.  The former is at least
$M^2$ for every $M\geq32$: this holds at $32$, and
$2M^2\geq(M+1)^2$ for $M\geq3$ proves induction.
For $t\leq-1$, one instead has
$t^2 2^{M-t}\geq2^{M+1}\geq M^2$.
Thus the quadratic term alone bounds $2^{t-M}$
whenever $t\neq0$, proving (SH-57.1).

For one of the three term profiles, write
$\mu=\mathbb E\kappa$ and $v=\operatorname{Var}\kappa$.
Here $0\leq\kappa\leq72$.  The geometric bounds
SH-55 and $\log_2 3>19/12$ give
\[
                 v\leq216,
                 \qquad\mathbb E\kappa^2<53^2=2809.
 \tag{SH-57.2}
\]
If $\mu\leq40$, choose $a=\lfloor\mu\rfloor$,
so $M=72-a\geq32$.  Taking expectations in
(SH-57.1) yields
\[
 \mathbb E2^{\kappa-72}
 \leq\frac{v+(\mu-a)^2}{M^2}+2^{-M}
 <\frac{217}{1024}+2^{-32}<\frac13.
 \tag{SH-57.3}
\]
If $\mu\geq40$, choose $a=40$ and $M=32$.
The centered second moment satisfies
\[
 \mathbb E(\kappa-40)^2
 \leq\min\{(\mu-40)^2+216,\ 4409-80\mu\}.
\]
For $\mu\geq40$, the first function is increasing
and the second decreasing.  They meet at
$\mu=\sqrt{2593}$; hence the minimum is at most
$4409-80\sqrt{2593}<337$, since
$2593>(509/10)^2$.
Consequently (SH-57.1) gives
\[
       \mathbb E2^{\kappa-72}
       <\frac{337}{1024}+2^{-32}<\frac13.
 \tag{SH-57.4}
\]
The last strict comparison is rational:
$1/3-337/1024=13/3072>2^{-32}$.

Finally the target has rank $2^{72}$ on every sign cut.
Subadditivity of flattening rank therefore implies
$\sum_{j=1}^3 2^{\kappa_j(s)-72}\geq1$ pointwise.
Taking expectations contradicts (SH-57.3)--(SH-57.4)
for the three profiles.
\end{proof}

The exponential-moment argument was found by the
exponent-audit agent and independently checked by the
shortening agent.  It requires no numerical certificate.
The remaining width-six parameter range is now
\[
                         A=3,\qquad4\leq D\leq11.
\]

\begin{theorem}[SH-58: the moment row closes dimension four]
\label{thm:width-six-three-term-dimension-four}
Over every field and for arbitrary finite network layouts,
there is no strict shortening with $m=6$, $A=3$, and $D=4$.
\end{theorem}
\begin{proof}
There are eight complementary sign cuts.  Write their
three rank-ceiling profiles as $\kappa_{j,s}$, with
$j\in\{1,2,3\}$ and $s\in\{0,1\}^3$, using the
representative with fourth sign fixed.  The folded-cube
distance is $h(s,t)=\min\{|s\triangle t|,4-|s\triangle t|\}$.
The geometric first and second moments imply
\[
 \sum_s\kappa_{j,s}\leq141,
 \qquad
 \sum_s\kappa_{j,s}^2\leq2496.
 \tag{SH-58.1}
\]
For the second bound, SH-55 gives the strict integer
inequality $144\sum_s\kappa_{j,s}^2<2809\cdot4^2\cdot8$;
division and rounding yield $2496$.
The other necessary rows are the terminal-dimension
bounds $0\leq\kappa\leq24$, the Lipschitz bounds
$|\kappa_{j,s}-\kappa_{j,t}|\leq12h(s,t)$,
and the single-cut inequalities
$\sum_j2^{\kappa_{j,s}}\geq2^{24}$.
For each ordered pair of distinct cuts, put
$p=2^{6(4-h)}$.  The rectangular results SH-39 and
SH-44 give
\[
 \sum_j\min\{2^{\kappa_{j,s}},w2^{\kappa_{j,t}}\}
 \geq L_{h,w},
 \qquad
 \begin{cases}
 L_{h,1}=\lceil67p/49\rceil,\\
 L_{h,3}=2p,\\
 L_{h,z^2}=zp& (z=2,4,\ldots,2^{6h}).
 \end{cases}
 \tag{SH-58.2}
\]
All these inequalities use upper bounds on the actual
flattening ranks, so no equality of ranks and minimum
cut capacities is assumed.

By SH-48 the absolute minimum ceiling $t$ belongs to
$\{0,\ldots,13\}$.  Cube and term symmetries put this
minimum at the first term and cut zero.  The other two
ceilings at that cut either include a $24$, which may
be placed in the second term, or both equal $23$.
These give exactly twenty-eight root cases.  At any
other cut a permissible triple must dominate one of
the nine distinct permutations of
\[
             (24,t,t),\qquad(23,23,t),\qquad(23,22,22).
 \tag{SH-58.3}
\]
Indeed if no ceiling equals $24$, at least one equals
$23$; unless a second equals $23$, the two others
must both equal $22$ for their rank upper bounds
to sum to $2^{24}$.

The retained finite proof tree splits only by these
necessary patterns.  Between splits it propagates
integer interval bounds using (SH-58.1)--(SH-58.2).
For example, if the current lower bounds in one row
are $l_s$, its moment row yields the upper bound
\[
 \kappa_s\leq
 \left\lfloor\sqrt{2496-\sum_u l_u^2+l_s^2}\right\rfloor.
\]
A weighted-rank row raises one lower bound only after
subtracting the upper rank bounds of all other terms.
Every such update, omitted pattern, child interval box,
and final contradiction is recorded as an integer
certificate.

The full tree has $7013$ reachable nodes, consisting
of $1286$ complete splits and $5727$ contradictions.
The contradictions comprise $3190$ first-moment
violations, $574$ second-moment violations, $1273$
single-cut rank violations, and $690$ weighted
mixed-cut violations.  There are no surviving or
interrupted branches.  A separate verifier, importing
neither the generator nor an optimizer, reconstructed
the folded-cube geometry and every necessary row,
replayed all $127524$ integer deductions, checked
all twenty-eight roots and all children and omissions,
and verified reachability and complete coverage.
Thus every possible rank profile is excluded.
\end{proof}

The generator is
\path{programs/sh_three_term_moment_tree.py}, with
input \path{inputs/sh_three_term_moment_D4_tree.json}
and exact proof output
\path{outputs/sh_three_term_moment_D4_tree.json}.
The independent semantic replay is
\path{programs/rt_verify_three_term_moment_tree.py};
its input and verified output are
\path{inputs/rt_three_term_moment_D4_verification.json}
and \path{outputs/rt_three_term_moment_D4_verification.json}.
The older incomplete dimension-four searches are
preserved but are not premises of this theorem.

\begin{theorem}[SH-59: cube isoperimetry excludes four more dimensions]
\label{thm:width-six-three-term-harper}
Over every field and for arbitrary finite network layouts,
no strict shortening with $m=6$, $A=3$, and
$D\in\{8,9,10,11\}$ exists.
\end{theorem}
\begin{proof}
At each sign cut at least one of the three ceilings is
at least $6D-1$.  Evenness of each profile allows us
to fix the last sign, retaining its probability
distribution on an ordinary cube $Q_d$ with $d=D-1$.
Some profile $\kappa$ therefore has a high set
\[
 H=\{s\in Q_d:\kappa(s)\geq6D-1\},
 \qquad |H|\geq K=\lceil2^d/3\rceil.
\]
This restriction remains $12$-Lipschitz along each
ordinary cube edge.  Consequently
\[
 \kappa(s)^2\geq f_D(\operatorname{dist}(s,H)),
 \qquad f_D(r)=\max\{0,6D-1-12r\}^2.
 \tag{SH-59.1}
\]

We use Harper's vertex-isoperimetric theorem in the
form stated by R\"aty~\cite{Raty2020}, Theorem~1 and
the paragraph following it (arXiv version~2, p.~2).
Order subsets first by cardinality and then by the
lexicographic rule that the least element of their
symmetric difference belongs to the earlier set.
If $I_K$ is the initial segment of size $K$, every
closed $r$-neighborhood of $H$ has at least the size
of the corresponding neighborhood of $I_K$.
For $|H|>K$, apply the theorem to any $K$-element
subset of $H$.  The extension from radius one to
every radius follows because the closed neighborhood
of an initial segment is itself an initial segment.

Thus distance to $H$ is stochastically no larger than
distance to $I_K$.  Since $f_D$ is decreasing,
\[
 \mathbb E\kappa^2\geq
 2^{-d}\sum_{r\geq0}L_r f_D(r),
 \qquad
 L_r=|\{s:\operatorname{dist}(s,I_K)=r\}|.
 \tag{SH-59.2}
\]
This expectation comparison also follows directly by
writing a decreasing function as the sum of its
nonnegative downward differences times indicators
of closed neighborhoods.

The finite distance counts and exact moment gaps are
as follows.  The last column subtracts the strict
upper bound $2809D^2/144$ supplied by SH-55.
\[
\begin{array}{c|r|l|r}
 D&K&(L_0,L_1,\ldots)&\text{moment gap}\\\hline
 8&43&(43,41,30,12,2)&12299/576\\
 9&86&(86,75,58,28,8,1)&20223/256\\
 10&171&(171,148,118,57,16,2)&450299/2304\\
 11&342&(342,272,227,126,46,10,1)&2637959/9216
\end{array}
\]
Every gap is positive, contradicting the strict
geometric second-moment bound.  The counts were
independently reconstructed in two ways: breadth-first
search on the ordinary cube, and exhaustive direct
Hamming-distance minimization.  Each computed closed
neighborhood was also checked to equal the initial
simplicial segment of its own cardinality.
\end{proof}

The retained exact calculation is
\path{programs/sh_harper_rank_moments.py}, with
\path{inputs/sh_harper_rank_moments.json} and
\path{outputs/sh_harper_rank_moments.json}.
The independent checker and its full arithmetic output
are \path{programs/rt_harper_moment_bounds.py} and
\path{outputs/rt_harper_moment_bounds.json}.
An additional independent layer calculation is
\path{programs/ex_verify_harper_layers.py}, with
\path{inputs/ex_harper_layers_parameters.json} and
\path{outputs/ex_harper_layers_independent.json}.
Together with SH-57 and SH-58, the remaining
width-six parameter possibilities are now only
\[
                         A=3,\qquad D\in\{5,6,7\}.
\]
The same isoperimetric estimate has negative margins
in those three dimensions; it does not decide them.

\begin{theorem}[SH-60: full tail counts close dimensions five, six, and seven]
\label{thm:width-six-three-term-tail-counts}
Over every field and for arbitrary finite network layouts,
there is no strict shortening with $m=6$, $A=3$, and
$D\in\{5,6,7\}$.
\end{theorem}
\begin{proof}
Put $M=2^{D-1}$ and $n=6D$.  As in SH-59, each
even ceiling profile has the same distribution on
the section with last sign fixed.  Rank subadditivity
on all cuts gives
\[
 \sum_{j=1}^3\sum_s2^{\kappa_j(s)}\geq M2^n.
\]
Consequently one profile $\kappa$ satisfies
\[
 \sum_s2^{\kappa(s)}\geq R_D,
 \qquad R_D=\left\lceil\frac{M2^n}{3}\right\rceil.
 \tag{SH-60.1}
\]
For $1\leq k\leq n$, let
$q_k=|\{s:\kappa(s)\geq k\}|$.  Introduce the
integer indicators $y_{k,v}=\mathbf1[q_k\geq v]$,
where $1\leq v\leq M$.  They satisfy
\[
\begin{split}
 &0\leq y_{k,v}\leq1,\qquad
 y_{k,v+1}\leq y_{k,v},\qquad
 y_{k+1,v}\leq y_{k,v},\\
 &\sum_{k,v}y_{k,v}\leq
       \left\lfloor\frac{53DM-1}{12}\right\rfloor,\\
 &\sum_{k,v}(2k-1)y_{k,v}\leq
       Q_D:=\left\lfloor\frac{2809D^2M-1}{144}\right\rfloor,\\
 &M+\sum_{k,v}2^{k-1}y_{k,v}\geq R_D.
\end{split}
 \tag{SH-60.2}
\]
The first two weighted sums are exactly
$\sum_s\kappa(s)$ and $\sum_s\kappa(s)^2$.
Their strict geometric bounds explain the subtracted
ones before integer division.  The last equality
uses $2^a=1+\sum_{k=1}^a2^{k-1}$ for each integer
$a\geq0$.

Let $g_r(v)$ be a certified lower bound on the size
of the closed radius-$r$ neighborhood of any
$v$-element set in the relevant section graph.
The $12$-Lipschitz property gives the additional
linear implications
\[
       y_{k,v}\leq y_{k-12r,g_r(v)}
       \qquad(k-12r\geq1).
 \tag{SH-60.3}
\]
For $D=6,7$, the graph is the ordinary cube
$Q_{D-1}$ and $g_r(v)$ is the exact neighborhood
size of its initial simplicial segment, by Harper's
theorem.  For $D=5$ we retain the complement edge
of the folded cube: on $\{0,1\}^4$, vertices are
adjacent when their difference is one coordinate
vector or $(1,1,1,1)$.  The complement edge is
legitimate because flipping the fifth original sign
and then reversing all signs preserves the ceiling.
An exhaustive enumeration of its $2^{16}=65536$
vertex subsets gives, for $v=0,\ldots,16$,
\[
\begin{split}
 (g_1(v))_{v=0}^{16}
   ={}&(0,6,10,12,12,14,14,15,15,15,15,\\
       &16,16,16,16,16,16).
\end{split}
 \tag{SH-60.4}
\]
For every nonempty set, $g_r(v)=16$ when $r\geq2$,
as the graph has diameter two.  The enumeration
forms each closed neighborhood by a bitwise union
of the sixteen explicitly known closed vertex balls
and minimizes its cardinality separately at each
subset cardinality.  An independent implementation
repeated this enumeration in full.

Only finitely many values of the top count $q_n$
need to be considered.  The following implications
come directly from the same neighborhood bounds;
the middle column is the resulting lower bound on
$\sum_s\kappa(s)^2$:
\[
\begin{array}{c|l|r|r}
 D&\text{hypothesis}&\text{moment lower}&Q_D\\\hline
 5&q_{30}\geq6&8064&7802\\
 6&q_{36}\geq11&22752&22471\\
 7&q_{42}\geq22&61884&61173
\end{array}
 \tag{SH-60.5}
\]
For instance, the dimension-five row uses six values
at least $30$, eight further values at least $18$,
and two further values at least $6$.
The dimension-six row uses the ordinary-cube
distance layers $(11,13,7,1)$ at threshold $36$;
the dimension-seven row uses $(22,20,15,6,1)$ at
threshold $42$.  Thus these hypotheses are impossible.

In dimension six, the final top-count case
$q_{36}=10$ is refined once.  If $q_{35}\geq11$,
the neighborhood bounds force the following tail
lower bounds:
\[
\begin{array}{c|rrrrrr}
 k&1\!:\!11&12&13\!:\!23&24&25\!:\!35&36\\\hline
 q_k\text{ lower}&31&30&24&22&11&10.
\end{array}
\]
Here the size-ten initial segment in $Q_5$ has
distance layers $(10,12,8,2)$, and the size-eleven
segment has layers $(11,13,7,1)$.  The resulting
second-moment lower bound is
\[
\begin{split}
 &31\cdot121+30\cdot23+24\cdot385
       +22\cdot47+11\cdot649+10\cdot71\\
 &\hspace{35mm}=22564>22471.
\end{split}
 \tag{SH-60.6}
\]
Hence $q_{35}=10$ whenever $q_{36}=10$.

It remains to check six cases $q_{30}=0,\ldots,5$
in dimension five; ten cases $q_{36}=0,\ldots,9$
and the case $q_{36}=q_{35}=10$ in dimension six;
and twenty-two cases $q_{42}=0,\ldots,21$ in
dimension seven.  An exact count $q_k=c$ is
enforced by $y_{k,c}=1$ when $c>0$ and
$y_{k,c+1}=0$ when $c<M$, together with monotonicity.
For each of these thirty-nine cases, a retained
nonnegative rational combination of the rows
(SH-60.2)--(SH-60.3) and the fixed-count rows has
every variable coefficient zero and right side
exactly $-1$.  Thus even the continuous relaxation
$0\leq y\leq1$ is impossible in each case.

The independent exponent-audit verifier rebuilt every
row from its mathematical label, recomputed all
ordinary-cube neighborhoods, exhaustively recomputed
the folded dimension-five growth, and checked every
rational coefficient, variable cancellation, negative
right side, file hash, and branch-coverage condition.
It imported neither a generator nor an optimization
routine.  All thirty-nine certificates and the
prunes (SH-60.5)--(SH-60.6) passed this audit.
The covered branches exhaust every integer tail
profile, contradicting (SH-60.1).
\end{proof}

The common generator is
\path{programs/rt_harper_tail_certificate.py}; its
outputs are used only through the exact certificates.
The dimension-five and dimension-seven manifests are
\path{outputs/rt_folded_D5_top_batch.json} and
\path{outputs/rt_harper_D7_top_batch.json}.
The dimension-six selected certificates, hashes and
both direct prunes are recorded in
\path{outputs/sh_D6_tail_coverage.json}, constructed
by \path{programs/sh_D6_tail_coverage.py} from
\path{inputs/sh_D6_tail_coverage.json}.
The independent verifier is
\path{programs/ex_verify_top_tail_certificates.py},
with input
\path{inputs/ex_top_tail_verification_parameters.json}
and full result
\path{outputs/ex_top_tail_certificates_independent.json}.
Earlier failed or unrefined dual attempts are retained
but are not premises of the theorem.  A corrected
nonsemantic dimension-six diagnostic string, with
unchanged certificate coefficients, is documented in
\path{outputs/sh_D6_tail_metadata_correction.json}.

\begin{corollary}[SH-60b: strict shortening requires width at least seven]
\label{cor:strict-shortening-width-at-least-seven}
Over every field, including $\mathbb C$ and $\mathbb F_2$,
there is no strict shortening of cable width $m\leq6$,
in any dimension and with any finite number of summands
or arbitrary finite term layouts.  In dimension two,
a shortening with exactly two summands has width at
least eight.
\end{corollary}
\begin{proof}
SH-28d excludes all widths at most five.  At width
six, one summand is excluded by the one-network
Crofton bound, and four or more are excluded by
Proposition~\ref{rt:width-six-parameters}.  SH-55d
excludes two summands in every dimension.
For three summands, SH-38 covers $D=2$, SH-46
covers $D=3$, SH-58 covers $D=4$, SH-60 covers
$D=5,6,7$, SH-59 covers $D=8,9,10,11$, SH-57
covers $D=12$, and SH-55b covers every $D\geq13$.
Dimension one, if allowed, is excluded directly:
each term of cost $w$ has rank at most $2^{w/2}$,
so $A$ terms of cost strictly below
$2(m-\log_2 A)$ cannot sum to rank $2^m$.
The final planar assertion is SH-38.
\end{proof}

This finite-width theorem leaves width seven as the
first unexcluded width in the unrestricted problem.
It supplies neither a strict shortening at a larger
width nor a proof that shortening is impossible at
every width, and it does not settle the remaining
flat-versus-contracting classification.
